%%%%%%%%%%%%%%%%%%%%%%%%%%%%%%%%%%%%%%%%%%%%%%%%%%%%%%%%%%%%%%%%%%%%%%%%%%%%%%%
%
% APPENDIX B: COMPLEMENTARITY STRUCTURE (CORE-HOW)
% 
% The Complete Specification of How Utility Dimensions Interact
%
% EBF Framework – Behavioral Change Model
% FehrAdvice & Partners AG / BEATRIX Research Group
%
% Status: CORE Appendix (5C-CORE Standard)
% Version: 2.1
% Category: LIT
% Language: English
% Last Updated: 2026-01-06
%
%%%%%%%%%%%%%%%%%%%%%%%%%%%%%%%%%%%%%%%%%%%%%%%%%%%%%%%%%%%%%%%%%%%%%%%%%%%%%%%

%==============================================================================
% 1.1 HEADER BLOCK (Required)
%==============================================================================

\begin{tcolorbox}[colback=red!5!white, colframe=red!75!black, 
    title=\textbf{CORE Appendix: Complementarity Structure}]

\begin{tabular}{@{}ll@{}}
\textbf{Appendix:} & B (Complementarity Structure) \\
\textbf{Category:} & CORE \\
\textbf{CORE Question:} & \textbf{HOW} --- How do utility dimensions interact? \\
\textbf{Scope:} & A3 (Methodological Framework) \\
\textbf{Version:} & 2.1 \\
\textbf{Last Updated:} & January 2026 \\
\textbf{Dependencies:} & Appendix WHO (Aggregation Levels) \\
 & Appendix WAT (FEPSDE Utility Structure) \\
 & Appendix CTW (8Ψ Context Dimensions) \\
\textbf{Outputs:} & $\Gamma$-matrix (Utility$\times$Utility: 15 parameters) \\
 & $\Lambda$-matrix (Context$\times$Context: 28 parameters) \\
 & $\delta$-tensor (Context modulation: 120 parameters) \\
\end{tabular}

\end{tcolorbox}

%==============================================================================
% 1.2 CHAPTER LINKAGE BOX (Required)
%==============================================================================

\begin{tcolorbox}[colback=purple!5!white, colframe=purple!75!black, 
    title=\textbf{Chapter Linkages}]

\textbf{Primary Chapter:} Chapter 4: Complementarity and Interaction Effects
\begin{itemize}[nosep]
    \item Section 4.1: Foundations of Complementarity $\leftrightarrow$ Part I
    \item Section 4.2: Utility Interactions $\leftrightarrow$ Part II
    \item Section 4.3: Context Effects $\leftrightarrow$ Parts III--IV
    \item Section 4.4: Strategic Implications $\leftrightarrow$ Parts V--VIII
\end{itemize}

\textbf{Secondary Chapters:}
\begin{itemize}[nosep]
    \item Chapter 3: Model Specification --- uses $\Gamma$-matrix in utility function
    \item Chapter 5: Interventions --- mechanism design under complementarity
    \item Chapter 7: Predictions --- derives testable implications from $\Gamma$
    \item Chapter 8: Applications --- calibrated $\Gamma$ for specific contexts
\end{itemize}

\textbf{Reading Guide:}
\begin{itemize}[nosep]
    \item For conceptual overview: Read Chapter 4 first
    \item For mathematical details: This appendix provides complete specification
    \item For empirical values: See Appendix EST (Parameter Calibration)
\end{itemize}

\end{tcolorbox}

%==============================================================================
% 1.3 CROSS-REFERENCE MAP (Required)
%==============================================================================

\begin{tcolorbox}[colback=white, colframe=black, 
    title=\textbf{Cross-Reference Map}]

\begin{center}
\small
\begin{tabular}{@{}c|cccccc@{}}
\toprule
& \textbf{AAA} & \textbf{C} & \textbf{V} & \textbf{BBB} & \textbf{AN} & \textbf{AU} \\
\midrule
\textbf{B $\leftarrow$} & Levels $L$ & FEPSDE & 8Ψ & Calibration & LLM-MC & Axioms \\
\textbf{B $\rightarrow$} & $\Gamma$ by level & $\gamma_{dd'}$ & $\gamma(\Psi)$ & $\delta$-tensor & Estimation & B-axioms \\
\midrule
\textbf{Link} & \cellcolor{red!20}Strong & \cellcolor{red!20}Strong & \cellcolor{red!20}Strong & \cellcolor{red!20}Strong & \cellcolor{yellow!20}Medium & \cellcolor{yellow!20}Medium \\
\bottomrule
\end{tabular}
\end{center}

\textbf{Legend:}
\colorbox{red!20}{Strong} = Bidirectional CORE dependency \quad
\colorbox{yellow!20}{Medium} = One-way dependency

\textbf{New in 2026:} See also \textbf{Appendix HIE (CORE-HIERARCHY)}---Decision Hierarchy \& Strategic Coherence. HI formalizes how the $\Gamma$ complementarity matrix creates stratification into L0-L1-L2-L3 decision levels via Theorem T1 (Universal L2 Complexity Formula). HI Axiom A1 grounds the Complementarity Principle in Milgrom \& Roberts (1995).

\textbf{Critical Link:} \textbf{Appendix NC (FORMAL-NC)}---Non-Concavity, Fit, and Coherence. NC formalizes \textit{why} $\gamma > 0$ creates multiple stable configurations (peaks) and establishes the $K$ (coherence) vs.\ $Q$ (quality) distinction central to Chapter 7.

\textbf{The 5C CORE Framework:}
\begin{center}
\begin{tabular}{@{}cllll@{}}
\toprule
\textbf{C} & \textbf{Question} & \textbf{Appendix} & \textbf{Output} \\
\midrule
WHO & Who has utility? & AAA & Level $L$ \\
WHAT & What is utility? & C & FEPSDE $d$ \\
\textbf{HOW} & \textbf{How do they interact?} & \textbf{B (this)} & $\Gamma, \Lambda, \delta$ \\
WHEN & When does context matter? & V & $\boldsymbol{\Psi}$ \\
WHERE & Where do numbers come from? & BBB & $\Theta, E(\theta)$ \\
\bottomrule
\end{tabular}
\end{center}

\end{tcolorbox}

%==============================================================================
% FORMAL COMPLIANCE BLOCK
%==============================================================================

% --- Terminology Reference ---
\begin{tcolorbox}[colback=gray!5!white, colframe=gray!50!black]
\textbf{Terminology:} See Appendix GLS (Glossary) for standard definitions.
Key terms in this appendix: complementarity ($\gamma > 0$), substitutability ($\gamma < 0$),
$\Gamma$-matrix (utility interactions), $\Lambda$-matrix (context interactions),
$\delta$-tensor (modulation coefficients), supermodularity.
\end{tcolorbox}

% --- Epistemic Status Tags ---
\begin{tcolorbox}[colback=blue!5!white, colframe=blue!75!black,
    title=\textbf{Epistemic Status: Key Parameters}]

\renewcommand{\arraystretch}{1.3}
\begin{tabular}{@{}llcc@{}}
\toprule
\textbf{Parameter} & \textbf{Description} & \textbf{Value/Range} & \textbf{Tag} \\
\midrule
$\gamma_{FE}$ & Financial-Emotional & 0.15--0.35 & [EMP] \\
$\gamma_{FS}$ & Financial-Social & 0.20--0.40 & [EMP] \\
$\gamma_{EP}$ & Emotional-Physical & 0.25--0.45 & [EMP] \\
$\gamma_{PS}$ & Physical-Social & 0.10--0.25 & [EMP] \\
$\gamma_{FD}$ & Financial-Digital & 0.05--0.15 & [LLM] \\
$\gamma_{Eco}$ & Ecological interactions & 0.05--0.20 & [LLM] \\
\addlinespace
\multicolumn{4}{@{}l}{\textit{Matrix dimensions:}} \\
$\Gamma$ & Utility$\times$Utility & 15 parameters & [THR] \\
$\Lambda$ & Context$\times$Context & 28 parameters & [THR] \\
$\delta$ & Context modulation & 120 parameters & [THR] \\
\bottomrule
\end{tabular}

\vspace{0.5em}
\textbf{Tag Legend:} [EMP] = Empirically validated; [THR] = Theoretically derived;
[LLM] = LLM Monte Carlo estimated

\end{tcolorbox}

% --- Bibliography Reference ---
\noindent\textbf{Full citations:} See \texttt{bibliography/bcm\_master.bib} \\
\textbf{Key sources:} Milgrom \& Roberts (1990); Topkis (1998); Samuelson \& Zeckhauser (1988)

\vspace{1em}

%==============================================================================
% 1.4 ABSTRACT (Required, max 100 words)
%==============================================================================

\begin{tcolorbox}[colback=white, colframe=black,
    title=\textbf{Abstract}]

Economic dimensions do not operate independently---they interact through 
\textbf{complementarity}. This appendix provides the complete mathematical 
specification of how utility dimensions (FEPSDE) reinforce or offset each other, 
how context (8Ψ) modulates these interactions, and how context dimensions 
themselves interact. We derive 163 interaction parameters organized in three 
structures: the $\Gamma$-matrix (utility$\times$utility), the $\Lambda$-matrix 
(context$\times$context), and the $\delta$-tensor (context modulation). 
Applications span mechanism design, organizational economics, matching theory, 
and large supermodular games.

\textit{(87 words)}

\end{tcolorbox}

%==============================================================================
% 1.4b METHODOLOGICAL PRINCIPLE (Fehr-Compatibility)
%==============================================================================

\begin{tcolorbox}[colback=red!5!white, colframe=red!75!black,
    title=\textbf{Methodological Principle: $\gamma = 0$ as Null Hypothesis}]

\textbf{The Exclusion Principle} (see Appendix FRM: LIT-FEHR-METHOD):

Following Fehr's methodological critique, complementarity coefficients $\gamma_{ij}$
must be understood as \textbf{exclusion} parameters, not integration parameters:

\begin{enumerate}[noitemsep]
    \item \textbf{Null Hypothesis:} $H_0: \gamma_{ij} = 0$ (additivity assumed)
    \item \textbf{Alternative:} $H_1: \gamma_{ij} \neq 0$ (complementarity claimed)
    \item \textbf{Scientific Value:} A complementarity claim gains value precisely
          where it can be \textit{refuted}---not where it can explain everything
\end{enumerate}

\textbf{What this means for EBF:}
\begin{itemize}[noitemsep]
    \item EBF identifies \textit{where} $\gamma \approx 0$ (boundary of applicability)
    \item EBF does \textit{not} claim universal complementarity ($\gamma > 0$ everywhere)
    \item Each $\gamma_{ij} \neq 0$ claim requires empirical justification against the null
\end{itemize}

\textit{``Complementarity is not a default assumption, but a risky hypothesis.''}

\textbf{Cross-Reference:} Appendix FRM provides the full philosophy of science
foundation for this principle (Popper, Lakatos, Cartwright).

\end{tcolorbox}

%==============================================================================
% 1.5 QUICK REFERENCE BOX (Required)
%==============================================================================

\begin{tcolorbox}[colback=yellow!10!white, colframe=yellow!75!black,
    title=\textbf{Quick Reference}]

\textbf{Core Equation --- Complementarity:}
\[
\gamma_{dd'} = \frac{\partial^2 U}{\partial x_d \partial x_{d'}} \quad
\begin{cases}
> 0 & \text{complements (reinforce)} \\
< 0 & \text{substitutes (offset)} \\
= 0 & \text{independent}
\end{cases}
\]

\textbf{The Three Interaction Matrices:}
\begin{center}
\small
\begin{tabular}{@{}llcc@{}}
\toprule
\textbf{Matrix} & \textbf{Interaction} & \textbf{Size} & \textbf{Parameters} \\
\midrule
$\Gamma = [\gamma_{dd'}]$ & Utility $\times$ Utility & $6 \times 6$ & 15 (symmetric) \\
$\Lambda = [\lambda_{kl}]$ & Context $\times$ Context & $8 \times 8$ & 28 (symmetric) \\
$\delta = [\delta_{dd',k}]$ & Context modulates $\Gamma$ & $15 \times 8$ & 120 \\
\midrule
& & \textbf{Total:} & \textbf{163} \\
\bottomrule
\end{tabular}
\end{center}

\textbf{Context Modulation:}
\[
\gamma_{dd'}(\boldsymbol{\Psi}) = \bar{\gamma}_{dd'} + \sum_{k=1}^{8} \delta_{dd',k} \cdot (\Psi_k - 0.5)
\]

\textbf{Jump to:}
\begin{itemize}[nosep]
    \item Mathematical Foundations $\rightarrow$ Part I (\S\ref{sec:B-foundations})
    \item Utility$\times$Utility $\Gamma$ $\rightarrow$ Part II
    \item Utility$\times$Context $\rightarrow$ Part III
    \item Context$\times$Context $\Lambda$ $\rightarrow$ Part IV
    \item Mechanism Design $\rightarrow$ Part V
    \item Large Games $\rightarrow$ Part VI
    \item Matching Theory $\rightarrow$ Part VII
    \item Organizations $\rightarrow$ Part VIII
    \item Axiom Index $\rightarrow$ \S\ref{sec:B-axiom-index}
\end{itemize}

\end{tcolorbox}





%==============================================================================
% PART I: FOUNDATIONS
%==============================================================================

\subsection{Part I: Mathematical Foundations of Complementarity}
\label{sec:B-foundations}

\subsubsection{The Central Question}

The EBF framework specifies utility as a function of multiple dimensions.
But these dimensions do not operate independently. The central question of 
this appendix is:

\begin{center}
\fbox{\parbox{0.8\textwidth}{
\textbf{CORE-HOW Question:} Given utility dimensions $d_1, d_2, \ldots, d_n$,
how does changing one dimension affect the marginal value of changing another?
}}
\end{center}

The answer lies in the concept of \textbf{complementarity}: the systematic 
pattern of interactions between choice variables that determines whether 
they reinforce or substitute for each other.

\subsubsection{Fundamental Definition}

\begin{definition}[Complementarity]
\label{def:complementarity}
Two variables $x_i$ and $x_j$ are \textbf{complements} if increasing one 
raises the marginal return to increasing the other:
\begin{equation}
    \frac{\partial^2 U}{\partial x_i \partial x_j} \geq 0
    \label{eq:complementarity-def}
\end{equation}
They are \textbf{substitutes} if the cross-partial is negative:
\begin{equation}
    \frac{\partial^2 U}{\partial x_i \partial x_j} < 0
    \label{eq:substitutes-def}
\end{equation}
They are \textbf{independent} if the cross-partial equals zero.
\end{definition}

\paragraph{Notation.}
We denote the complementarity coefficient between dimensions $i$ and $j$ as:
\begin{equation}
    \gamma_{ij} \equiv \frac{\partial^2 U}{\partial x_i \partial x_j}
    \label{eq:gamma-def}
\end{equation}

The complete set of pairwise complementarities forms the \textbf{$\gamma$-matrix}:
\begin{equation}
    \Gamma = \begin{pmatrix}
        \gamma_{11} & \gamma_{12} & \cdots & \gamma_{1n} \\
        \gamma_{21} & \gamma_{22} & \cdots & \gamma_{2n} \\
        \vdots & \vdots & \ddots & \vdots \\
        \gamma_{n1} & \gamma_{n2} & \cdots & \gamma_{nn}
    \end{pmatrix}
    \label{eq:gamma-matrix}
\end{equation}

\paragraph{Symmetry.}
By Young's theorem (equality of mixed partials), the $\gamma$-matrix is symmetric:
\begin{equation}
    \gamma_{ij} = \gamma_{ji} \quad \forall i, j
    \label{eq:gamma-symmetry}
\end{equation}

\subsubsection{Supermodularity: The Mathematical Foundation}

The theory of supermodular functions provides the rigorous foundation for 
complementarity analysis. This section draws heavily on \citet{topkis1978}
and \citet{milgromroberts1990compare}.

\begin{definition}[Supermodularity]
\label{def:supermodularity}
A function $f: \mathbb{R}^n \to \mathbb{R}$ is \textbf{supermodular} if 
for all $x, y \in \mathbb{R}^n$:
\begin{equation}
    f(x \vee y) + f(x \wedge y) \geq f(x) + f(y)
    \label{eq:supermodular-lattice}
\end{equation}
where $x \vee y$ denotes the component-wise maximum and $x \wedge y$ the 
component-wise minimum.
\end{definition}

\begin{theorem}[Supermodularity and Cross-Partials]
\label{thm:supermod-crosspartial}
If $f$ is twice continuously differentiable, then $f$ is supermodular if 
and only if:
\begin{equation}
    \frac{\partial^2 f}{\partial x_i \partial x_j} \geq 0 \quad \forall i \neq j
    \label{eq:supermod-condition}
\end{equation}
\end{theorem}

\begin{proof}
See \citet{topkis1978}, Theorem 2.6.1, or \citet{milgromroberts1990compare}, 
Lemma 1.
\end{proof}

\paragraph{Interpretation.}
Supermodularity means all variables are pairwise complements. The lattice 
formulation (\ref{eq:supermodular-lattice}) captures the intuition that 
``moving up together'' in all dimensions is better than ``mixing'' high 
and low values across dimensions.

\subsubsection{Submodularity: The Opposite Case}

\begin{definition}[Submodularity]
\label{def:submodularity}
A function $f: \mathbb{R}^n \to \mathbb{R}$ is \textbf{submodular} if:
\begin{equation}
    f(x \vee y) + f(x \wedge y) \leq f(x) + f(y)
    \label{eq:submodular-lattice}
\end{equation}
Equivalently, for smooth functions:
\begin{equation}
    \frac{\partial^2 f}{\partial x_i \partial x_j} \leq 0 \quad \forall i \neq j
    \label{eq:submod-condition}
\end{equation}
\end{definition}

Submodularity means all variables are pairwise substitutes. This implies
diminishing returns to ``doing more of everything.''

\subsubsection{Substitutes: The Opposite Case}

\paragraph{When Actions Offset Each Other.}
While much of the EBF framework emphasizes complementarity, substitutability
is equally important. Actions are \textbf{substitutes} when doing more of one
\textit{reduces} the return to doing more of another:

\begin{equation}
\gamma_{ij} = \frac{\partial^2 f}{\partial x_i \partial x_j} < 0
\label{eq:substitutes-detailed}
\end{equation}

\paragraph{Canonical Examples of Substitutes.}
\begin{itemize}
    \item \textbf{Cournot competition:} If you produce more, my marginal revenue
          falls---your output crowds out mine
    \item \textbf{Arms races:} If you arm more, I must arm more (but arms are
          costly)---a negative-sum dynamic
    \item \textbf{Status goods:} Your status gain is my status loss---positional
          competition creates substitutability
    \item \textbf{Time allocation:} More time on work means less time for family---
          a fixed resource creates substitutes
\end{itemize}

\paragraph{Consequences of Substitutability.}
When $\gamma_{ij} < 0$:
\begin{enumerate}
    \item \textbf{Best responses are decreasing:} If you do more, my optimal
          response is to do less: $\partial x_i^* / \partial x_j < 0$
    \item \textbf{Equilibria tend to be unique:} Unlike complements, substitutes
          typically yield a single equilibrium
    \item \textbf{Coordination is less critical:} Agents naturally differentiate
          rather than converge
    \item \textbf{Stability is more robust:} Systems with substitutes resist
          cascading dynamics
\end{enumerate}

\paragraph{Comparison: Complements vs. Substitutes.}
\begin{center}
\renewcommand{\arraystretch}{1.3}
\begin{tabular}{|l|c|c|}
\hline
\textbf{Property} & \textbf{Complements} ($\gamma > 0$) & \textbf{Substitutes} ($\gamma < 0$) \\
\hline
Best response & Increasing & Decreasing \\
Equilibria & Multiple possible & Unique (typically) \\
Coordination & Critical & Less important \\
Dynamics & Cascading, tipping & Dampening, stable \\
Examples & Network effects & Cournot competition \\
\hline
\end{tabular}
\end{center}

\subsubsection{Mixed Systems: The General Case}

Real utility functions are typically neither purely supermodular nor purely 
submodular. Some pairs of dimensions are complements ($\gamma_{ij} > 0$), 
others are substitutes ($\gamma_{ij} < 0$), and some are approximately 
independent ($\gamma_{ij} \approx 0$).

\begin{definition}[Mixed Complementarity System]
\label{def:mixed-system}
A utility function exhibits \textbf{mixed complementarity} if its $\gamma$-matrix 
contains both positive and negative off-diagonal elements:
\begin{equation}
    \exists \, i, j, k, l : \gamma_{ij} > 0 \text{ and } \gamma_{kl} < 0
    \label{eq:mixed-system}
\end{equation}
\end{definition}

The EBF framework accommodates mixed systems. The key analytical challenge 
is identifying which pairs are complements and which are substitutes, and how 
these relationships vary across contexts and aggregation levels.

\subsubsection{Monotone Comparative Statics}

A powerful result from supermodularity theory concerns how optimal choices 
respond to parameter changes.

\begin{theorem}[Monotone Comparative Statics]
\label{thm:monotone-statics}
Let $f(x; \theta)$ be supermodular in $x$ for each $\theta$, and suppose 
$f$ has increasing differences in $(x; \theta)$. Then the set of maximizers 
$X^*(\theta) = \arg\max_x f(x; \theta)$ is increasing in $\theta$ in the 
strong set order.
\end{theorem}

\begin{proof}
See \citet{milgromshannon1994} or \citet{topkis1998}.
\end{proof}

\paragraph{EBF Application.}
When context $\Psi$ shifts, optimal behavior $x^*$ moves monotonically if:
\begin{enumerate}
    \item The utility function is supermodular in the choice variables
    \item Utility has increasing differences between choices and context
\end{enumerate}

This provides the foundation for predicting how behavioral patterns shift 
with context changes (see Appendix CTW, CORE-WHEN).

\subsubsection{The Complete $\gamma$-Tensor}

The simple $\gamma$-matrix notation is insufficient for EBF because 
complementarities vary across:
\begin{itemize}
    \item Aggregation levels $L$ (CORE-WHO)
    \item Context states $\Psi$ (CORE-WHEN)
    \item Time $t$ (dynamic evolution)
\end{itemize}

\begin{definition}[The $\gamma$-Tensor]
\label{def:gamma-tensor}
The complete EBF complementarity structure is specified by the tensor:
\begin{equation}
    \gamma^L_{ij}(\Psi, t)
    \label{eq:gamma-tensor}
\end{equation}
where:
\begin{itemize}
    \item $i, j \in \{F, E, P, S, D, Ec\}$ index FEPSDE dimensions (from CORE-WHAT)
    \item $L \in \{0, 1, 2, 3, 4, 5, \infty\}$ indexes aggregation level (from CORE-WHO)
    \item $\Psi = (\Psi_1, \ldots, \Psi_8)$ is the context vector (from CORE-WHEN)
    \item $t$ indexes time
\end{itemize}
\end{definition}

\paragraph{Interpretation.}
$\gamma^L_{ij}(\Psi, t)$ answers: ``At level $L$, in context $\Psi$, at time $t$, 
how does an increase in dimension $i$ affect the marginal value of dimension $j$?''

\subsubsection{Seven Domains of Complementarity}

The EBF framework identifies seven distinct domains where complementarity 
relationships must be specified:

\begin{center}
\renewcommand{\arraystretch}{1.4}
\begin{tabular}{|c|l|l|c|}
\hline
\textbf{\#} & \textbf{Domain} & \textbf{Notation} & \textbf{Unique Pairs} \\
\hline
1 & Utility $\times$ Utility & $\gamma_{dd'}$ & 15 (FEPSDE) \\
2 & Utility $\times$ Context & $\gamma_{d\Psi}$ & 48 (6$\times$8) \\
3 & Context $\times$ Context & $\gamma_{\Psi\Psi'}$ & 28 (8$\times$8) \\
4 & Awareness $\times$ Willingness & $\gamma_{AW}$ & 4 quadrants \\
5 & PARC $\times$ Strategy & $\gamma_{PARC}$ & 6 + 4 \\
6 & Technology $\times$ All & $\gamma_{T*}$ & $\sim$20 \\
7 & Level $\times$ Level & $\gamma_{LL'}$ & 15 (6 levels) \\
\hline
& \textbf{Total} & & $\sim$130+ \\
\hline
\end{tabular}
\end{center}

Each domain is developed in detail in Parts II--VII of this appendix.

\subsubsection{Coherence and the $K$ Index}

When complementary elements are aligned, the system exhibits \textbf{coherence}.
When they are misaligned, performance suffers.

\begin{definition}[Coherence Index $K$]
\label{def:coherence-K}
The coherence index measures the degree of alignment across complementary 
dimensions:
\begin{equation}
    K = \frac{\text{Actual system performance}}{\text{Performance at optimal alignment}}
    \label{eq:K-def}
\end{equation}
where $K \in [0, 1]$, with $K = 1$ indicating perfect coherence.
\end{definition}

\paragraph{Roberts' ``Fit'' as $K$.}
\citet{roberts2004} describes organizational success as depending on ``fit'' 
between environment, strategy, and organization. In EBF terms:
\begin{equation}
    \text{Roberts' Fit} \equiv K \approx 1
\end{equation}

Poor fit (misalignment, ``mix-and-match'') corresponds to $K < 1$.

\subsubsection{Foundational Axioms}

We now state the structural axioms that govern complementarity in EBF.

\begin{axiom}[Symmetry]
\label{ax:B-symmetry}
\textbf{(B-SYM)} The complementarity relationship is symmetric:
\begin{equation}
    \gamma_{ij} = \gamma_{ji} \quad \forall i, j
\end{equation}
\end{axiom}

\begin{axiom}[Sign Stability]
\label{ax:B-sign}
\textbf{(B-SGN)} The sign of $\gamma_{ij}$ is stable across most contexts, 
though magnitude may vary:
\begin{equation}
    \text{sign}(\gamma_{ij}(\Psi)) = \text{sign}(\gamma_{ij}(\Psi')) \quad 
    \text{for most } \Psi, \Psi'
\end{equation}
\end{axiom}

\begin{axiom}[Context Modulation]
\label{ax:B-context}
\textbf{(B-CTX)} Context $\Psi$ modulates the magnitude of complementarities:
\begin{equation}
    |\gamma_{ij}(\Psi)| \neq |\gamma_{ij}(\Psi')| \quad \text{in general}
\end{equation}
\end{axiom}

\begin{axiom}[Level Variation]
\label{ax:B-level}
\textbf{(B-LVL)} Complementarity coefficients vary systematically with 
aggregation level:
\begin{equation}
    \gamma^L_{ij} \neq \gamma^{L'}_{ij} \quad \text{for } L \neq L'
\end{equation}
\end{axiom}

\begin{axiom}[Coherence Optimality]
\label{ax:B-coherence}
\textbf{(B-COH)} Maximum utility is achieved when complementary elements 
are aligned:
\begin{equation}
    U^* = U(x^*) \text{ where } x^* \text{ satisfies coherence conditions}
\end{equation}
\end{axiom}

\subsubsection{Why Complementarity Implies Coherence}

\paragraph{The Logic.}
When choices are complementary, a systematic pattern emerges:
\begin{enumerate}
    \item \textbf{High levels reinforce high levels:} Doing more of $x_i$ raises
          the return to doing more of $x_j$, creating upward momentum
    \item \textbf{Low levels reinforce low levels:} Doing less of $x_i$ reduces
          the return to $x_j$, creating downward momentum
    \item \textbf{Mixed configurations are unstable:} Being high on some dimensions
          but low on others wastes the complementarity gains
\end{enumerate}

Therefore, systems with strong complementarity tend toward \textbf{coherent
configurations}---states where all elements are aligned, either all high or all low.

\begin{proposition}[Coherence from Complementarity]
\label{prop:coherence-from-complementarity}
If the payoff function $f$ is supermodular and $\Phi$ is the best-response mapping,
then:
\begin{enumerate}
    \item $\Phi$ is monotone increasing
    \item The set of fixed points forms a complete lattice
    \item Highest and lowest Nash equilibria exist
    \item Starting from any point, best-response dynamics converge to a fixed point
\end{enumerate}
\end{proposition}

\begin{proof}
Standard result from \citet{topkis1998} and \citet{milgromroberts1990rationalizability}.
\end{proof}

\paragraph{Interpretation.}
Complementarity creates \textbf{natural attractors}---coherent configurations
that the system tends toward. These attractors are the $C^*$ of the EBF framework.
The proposition explains why firms, institutions, and societies tend toward
internally consistent configurations rather than random mixtures.

\subsubsection{Literature Foundation}

The mathematical theory of complementarity draws on several foundational
contributions:

\paragraph{Edgeworth (1881).}
First identified complementarity in consumption: the marginal utility of 
coffee increases with sugar \citep{edgeworth1881}.

\paragraph{Topkis (1978, 1998).}
Established the mathematical theory of supermodularity and its applications 
to optimization \citep{topkis1978, topkis1998}.

\paragraph{Milgrom \& Roberts (1990).}
Applied supermodularity to games and organizations, establishing that:
\begin{enumerate}
    \item Nash equilibria exist in supermodular games (no convexity needed)
    \item Equilibria form a complete lattice
    \item Comparative statics are monotone
\end{enumerate}
See \citet{milgromroberts1990compare, milgromroberts1990rationalizability}.

\paragraph{Milgrom \& Shannon (1994).}
Developed monotone comparative statics for general optimization problems 
\citep{milgromshannon1994}.

\paragraph{Roberts (2004).}
Applied complementarity theory to organizational design, introducing the 
concepts of ``fit,'' ``coherent patterns,'' and the dangers of ``mix-and-match'' 
\citep{roberts2004}. See Appendix FND for detailed coverage.

\paragraph{Ichniowski et al. (1997).}
Provided empirical evidence for complementarities in human resource management 
practices \citep{ichniowski1997}.

\paragraph{Bloom et al. (2012).}
Demonstrated complementarities between information technology and organizational 
practices \citep{bloom2012}.


%==============================================================================
% PART I EXTENSION: THE ORDINAL STATUS OF SUPERMODULARITY
%==============================================================================

\subsubsection{The Ordinal Status of Supermodularity}
\label{sec:ordinal-status}

A fundamental question for any theory built on complementarity is: 
What is the empirical content of supermodularity? Is it a strong assumption 
that restricts behavior, or a weak property that many preferences satisfy?

\citet{chambers2009supermodularity} provide a definitive answer that has 
profound implications for EBF.

\paragraph{The Historical Debate.}
Economists have long debated whether supermodularity---the formal expression 
of complementarity---has ordinal content:

\begin{itemize}
    \item \textbf{Allen, Hicks, Samuelson, Stigler (1930s--1970s):} 
    Supermodularity is a \emph{cardinal} property (depends on the specific 
    utility function, not just the preference ordering). Therefore, it has 
    no testable implications for behavior.
    
    \item \textbf{Milgrom \& Shannon (1994):} Introduced \emph{quasisupermodularity}, 
    an ordinal weakening of supermodularity that preserves monotone comparative 
    statics.
    
    \item \textbf{Chambers \& Echenique (2009):} Under monotonicity, 
    supermodularity and quasisupermodularity are \emph{ordinally equivalent}. 
    The earlier economists were right---but for the wrong reasons.
\end{itemize}

%------------------------------------------------------------------------------
\paragraph{The Chambers-Echenique Theorem.}

\begin{theorem}[Chambers-Echenique 2009]
\label{thm:chambers-echenique}
Let $(X, \preceq)$ be a finite lattice. A binary relation $R$ on $X$ has a 
\textbf{weakly increasing and quasisupermodular} representation if and only 
if it has a \textbf{weakly increasing and supermodular} representation.
\end{theorem}

\textbf{Translation:} Under monotonicity, supermodularity imposes \emph{no 
additional restrictions} beyond quasisupermodularity. Any monotonic preference 
that satisfies the ordinal property (quasisupermodularity) can be represented 
by a cardinal supermodular function.

\paragraph{The Constructive Proof.}
The proof is elegant and constructive. Given a quasisupermodular function $u$:
\begin{enumerate}
    \item Label the values of $u(X)$ as $\{u_1, \ldots, u_N\}$ with $u_1 < \cdots < u_N$
    \item Define the transformation $f(u_j) = 2^{j-1}$
    \item Then $g = f \circ u$ is \textbf{supermodular}
\end{enumerate}

The exponential transformation ``spreads out'' the utility values sufficiently 
that supermodularity is guaranteed.

%------------------------------------------------------------------------------
\paragraph{The Bombshell Corollary.}

\begin{corollary}[Chambers-Echenique Corollary 5]
\label{cor:strongly-increasing}
If a binary relation has a \textbf{strongly increasing} representation, 
then it has a \textbf{supermodular} representation.
\end{corollary}

\textbf{Implication:} \emph{Every} strictly monotonic preference---where more 
is always better---can be represented by a supermodular utility function!

This means:
\begin{itemize}
    \item Supermodularity is a \textbf{weak} assumption, not a strong one
    \item It is \textbf{not testable} with consumption data alone
    \item The $\gamma$-matrix \textbf{always exists} under monotonicity
\end{itemize}

%------------------------------------------------------------------------------
\paragraph{Implications for EBF.}

\textbf{The Good News:}
\begin{enumerate}
    \item The complementarity framework is \textbf{mathematically guaranteed} 
    to exist for any monotonic preference
    \item We are not imposing an arbitrary functional form---we are revealing 
    structure that is always present
    \item The $\gamma$-tensor is not ``invented'' but \emph{discovered}
\end{enumerate}

\textbf{The Challenge:}
\begin{enumerate}
    \item Supermodularity \emph{alone} has no empirical content
    \item We cannot \emph{test} whether $\gamma_{dd'} > 0$ or $\gamma_{dd'} < 0$ 
    from revealed preference alone
    \item The \emph{signs} and \emph{magnitudes} of $\gamma$ must come from 
    theory and direct measurement, not from consumption data
\end{enumerate}

%------------------------------------------------------------------------------
\paragraph{The Solution: Supermodularity Plus Additional Structure.}

Chambers \& Echenique note that supermodularity \emph{combined with other 
assumptions} does have testable implications:

\begin{theorem}[Chipman 1977; Quah 2007]
\label{thm:chipman-quah}
If a utility function is both \textbf{supermodular} and \textbf{concave}, 
then demand is \textbf{normal} (demand increases with income for all goods).
\end{theorem}

This means:
\begin{itemize}
    \item Supermodularity alone: not testable
    \item Concavity alone: not testable (Afriat 1967)
    \item Supermodularity + Concavity: \textbf{testable!} (normal demand)
\end{itemize}

\textbf{EBF Application:}
The FEPSDE utility function combines:
\begin{itemize}
    \item Supermodularity (complementarity structure via $\gamma$)
    \item Reference-dependence (Prospect Theory curvature)
    \item Context-dependence (modulation by $\Psi$)
\end{itemize}

These joint assumptions \emph{do} have testable implications, even though 
supermodularity alone does not.

%------------------------------------------------------------------------------

%------------------------------------------------------------------------------
\paragraph{Extension: Infinite Supermodularity.}

\citet{chateauneuf2017infinite} extend the Chambers-Echenique result even further, 
showing that the ordinal equivalence extends to \emph{infinite supermodularity}---a 
much stronger condition than standard supermodularity.

\begin{definition}[$\infty$-Supermodularity]
\label{def:inf-supermodular}
A function $v : X \to \mathbb{R}$ on a finite lattice $(X, \leq)$ is 
\textbf{$\infty$-supermodular} if for all $n \geq 2$ and all $x_1, \ldots, x_n \in X$:
\begin{equation}
v\left(\bigvee_{k=1}^{n} x_k\right) \geq 
\sum_{\emptyset \neq I \subseteq \{1,\ldots,n\}} (-1)^{|I|+1} \, 
v\left(\bigwedge_{i \in I} x_i\right)
\end{equation}
\end{definition}

Standard supermodularity is the special case $n = 2$. Infinite supermodularity 
requires the inequality to hold for \emph{any} number of elements simultaneously.

\begin{theorem}[Chateauneuf-Vergopoulos-Zhang 2017]
\label{thm:inf-supermodular}
A binary relation $\succeq$ on $X$ has a weakly increasing and quasisupermodular 
representation if and only if it has a weakly increasing and 
\textbf{$\infty$-supermodular} representation.
\end{theorem}

\textbf{The Complete Equivalence Chain.}
Under monotonicity on finite lattices, the following are all ordinally equivalent:

\begin{center}
\fbox{
\begin{minipage}{0.85\textwidth}
\centering
\vspace{0.3em}
Weak Quasi $\Leftrightarrow$ Quasi $\Leftrightarrow$ Strong Quasi 
$\Leftrightarrow$ Super $\Leftrightarrow$ $\infty$-Super
\vspace{0.3em}
\end{minipage}
}
\end{center}

The strongest version ($\infty$-supermodular) is just as ``free'' as the weakest!

\paragraph{EBF Implication: Higher-Order Complementarity.}
The existence of an $\infty$-supermodular representation means that 
complementarity extends beyond pairs:

\begin{itemize}
    \item $\gamma_2$: Pairwise complementarity (e.g., $F \times E$)
    \item $\gamma_3$: Three-way complementarity (e.g., $F \times E \times P$)
    \item $\gamma_4$: Four-way complementarity (e.g., $F \times E \times P \times S$)
    \item $\vdots$
    \item $\gamma_n$: $n$-way complementarity for all $n$ dimensions
\end{itemize}

This justifies modeling \emph{interaction effects} among multiple FEPSDE 
dimensions simultaneously, not just pairwise.

\begin{axiom}[Infinite Supermodular Existence]
\label{ax:B-INF}
\textbf{(B-INF)} Under monotonicity, a preference relation admits an 
$\infty$-supermodular representation. Higher-order complementarities 
($\gamma_3, \gamma_4, \ldots$) are therefore well-defined:
\begin{equation}
    \text{Monotonic preference} \Rightarrow \exists \text{ representation with all } 
    \gamma_n \text{ for } n \geq 2
\end{equation}
\end{axiom}



\paragraph{What Determines the Signs of $\gamma$?}

If supermodularity is ``free'' under monotonicity, what determines whether 
specific dimension pairs are complements ($\gamma > 0$) or substitutes 
($\gamma < 0$)?

\begin{enumerate}
    \item \textbf{Production Technology:} Physical constraints on joint 
    production (e.g., time cannot be spent on both work and leisure)
    
    \item \textbf{Psychological Mechanisms:} Mind-body connections, 
    hedonic adaptation, social comparison
    
    \item \textbf{Institutional Design:} Rules that align or misalign 
    individual and collective interests
    
    \item \textbf{Direct Measurement:} Experimental and survey evidence 
    on cross-domain effects
\end{enumerate}

The $\gamma$-matrix in EBF is \emph{calibrated} from these sources, 
not derived from revealed preference.

%------------------------------------------------------------------------------
\paragraph{Axiom: Ordinal Foundation of Complementarity.}

\begin{axiom}[Ordinal Equivalence]
\label{ax:B-ORD}
\textbf{(B-ORD)} Under weak monotonicity on finite lattices, supermodularity 
is ordinally equivalent to quasisupermodularity:
\begin{equation}
    \text{Monotonic + Quasisupermodular} \Leftrightarrow 
    \text{Monotonic + Supermodular representation exists}
\end{equation}
The existence of a $\gamma$-tensor representation is guaranteed for any 
monotonic preference.
\end{axiom}

\begin{axiom}[Testability Requires Joint Assumptions]
\label{ax:B-TEST}
\textbf{(B-TEST)} The signs and magnitudes of $\gamma_{dd'}$ are not 
identifiable from revealed preference alone. Testable predictions require 
supermodularity combined with additional structure (concavity, 
reference-dependence, or context-specificity).
\end{axiom}

%------------------------------------------------------------------------------
\paragraph{Literature: The Ordinal Debate.}

\begin{center}
\renewcommand{\arraystretch}{1.3}
\begin{tabular}{|l|l|l|}
\hline
\textbf{Reference} & \textbf{Contribution} & \textbf{Implication} \\
\hline
\citet{samuelson1947} & Supermodularity is cardinal & Not testable (claimed) \\
\citet{milgrom1994monotone} & Quasisupermodularity is ordinal & MCS preserved \\
\citet{chambers2009supermodularity} & Equivalence under monotonicity & Both right! \\
\citet{chipman1977} & Supermodular + concave & Normal demand \\
\citet{quah2007} & Generalized Chipman & Broader MCS results \\
\citet{afriat1967} & Concavity alone & Not testable \\
\hline
\end{tabular}
\end{center}



\subsubsection{What Classical Economics Assumes}

\begin{proposition}[The Classical Independence Assumption]
\label{prop:classical-independence}
Classical economic analysis typically assumes $\gamma_{ij} = 0$ for all $i \neq j$.
\end{proposition}

This assumption allows optimization to proceed dimension-by-dimension, 
greatly simplifying analysis. However, it misses the interaction effects 
that are central to behavioral economics and organizational design.

\paragraph{EBF Contribution.}
By explicitly modeling $\gamma_{ij} \neq 0$, EBF captures:
\begin{itemize}
    \item Why isolated interventions often fail
    \item Why coherent ``bundles'' of changes succeed
    \item Why context matters for intervention effectiveness
    \item Why organizational change is difficult (complementarity creates inertia)
\end{itemize}

\subsubsection{Part I Summary}

This section established:

\begin{enumerate}
    \item \textbf{Definition}: Complementarity as $\gamma_{ij} = \partial^2 U / \partial x_i \partial x_j$
    \item \textbf{Theory}: Supermodularity provides the mathematical foundation
    \item \textbf{Notation}: The $\gamma^L_{ij}(\Psi, t)$ tensor captures full structure
    \item \textbf{Domains}: Seven domains of complementarity identified
    \item \textbf{Coherence}: The $K$ index measures alignment
    \item \textbf{Axioms}: Five foundational axioms (B-SYM, B-SGN, B-CTX, B-LVL, B-COH)
    \item \textbf{Literature}: Topkis, Milgrom-Roberts, Roberts as key sources
\end{enumerate}

The following parts develop each complementarity domain in detail.

%==============================================================================
% CROSS-REFERENCES TO OTHER PARTS
%==============================================================================

\vspace{1cm}
\begin{tcolorbox}[colback=gray!10, colframe=gray!50, title=Appendix HOW Structure]
\textbf{Part I: Foundations} (This section) \\
Part II: Utility $\times$ Utility Complementarity ($\gamma_{dd'}$) \\
Part III: Utility $\times$ Context Complementarity ($\gamma_{d\Psi}$) \\
Part IV: Awareness $\times$ Willingness Complementarity ($\gamma_{AW}$) \\
Part V: PARC $\times$ Strategy Complementarity ($\gamma_{PARC}$) \\
Part VI: Technology Complementarity ($\gamma_{T*}$) \\
Part VII: Level Complementarity ($\gamma^L$, $\gamma_{LL'}$) \\
Part VIII: Integration, Critical Foundations, and Open Issues \\

\vspace{0.3cm}
\textit{Cross-references:}
\begin{itemize}
    \item For formal supermodularity theory: See Appendix FND (Milgrom-Roberts)
    \item For FEPSDE dimensions: See Appendix WAT (CORE-WHAT)
    \item For aggregation levels: See Appendix WHO (CORE-WHO)
    \item For context specification: See Appendix CTW (CORE-WHEN)
\end{itemize}
\end{tcolorbox}


%==============================================================================
% PART II: UTILITY × UTILITY COMPLEMENTARITY (γ_dd')
%==============================================================================

\subsection{Part II: Utility × Utility Complementarity}
\label{sec:B-utility-utility}

This part specifies the complementarity relationships between the six FEPSDE 
utility dimensions. With 6 dimensions, there are $\binom{6}{2} = 15$ unique 
pairs to analyze.


%==============================================================================
% PART II EXTENSION: INTER-CATEGORY AND CATEGORY-SPECIFIC COMPLEMENTARITY
%==============================================================================

\subsubsection{The Three Levels of Utility Complementarity}
\label{sec:three-levels-complementarity}

The EBF framework distinguishes three levels at which utility complementarities 
operate. This subsection provides the overview; subsequent sections develop each level.


%------------------------------------------------------------------------------
\subsubsection{Computational Aspects: r-Decomposability and Optimization}
\label{sec:r-decomposable}

While the preceding sections establish \emph{existence} of supermodular representations, 
practical applications require \emph{computational tractability}. Recent work by 
\citet{chen2025supermodular} provides crucial results on optimizing supermodular 
functions under constraints---directly relevant to EBF intervention design.

\paragraph{The r-Decomposability Concept.}

\begin{definition}[r-Decomposability]
\label{def:r-decomposable}
A supermodular function $f: 2^V \to \mathbb{R}_+$ is \textbf{r-decomposable} if there 
exist subsets $V_1, \ldots, V_m \subseteq V$ with $|V_i| \leq r$ for each $i$, and 
corresponding supermodular functions $f_i: 2^{V_i} \to \mathbb{R}_+$, such that:
\begin{equation}
f(S) = \sum_{i=1}^{m} f_i(S \cap V_i) \quad \text{for all } S \subseteq V
\end{equation}
\end{definition}

The parameter $r$ captures the ``interaction order'' of the function:
\begin{itemize}
    \item $r = 1$: Function is \emph{modular} (additively separable, no interactions)
    \item $r = 2$: Pairwise interactions only (graph structure)
    \item $r = 3$: Three-way interactions (hypergraph with rank 3)
    \item $r = n$: Trivial decomposition (arbitrary interactions)
\end{itemize}

\textbf{EBF Connection.} The FEPSDE utility structure with $\gamma$-matrix is 
naturally 2-decomposable when only pairwise complementarities are modeled:
\begin{equation}
U(\mathbf{d}) = \sum_{d} u_d + \sum_{d < d'} \gamma_{dd'} \cdot u_d \cdot u_{d'}
\end{equation}
Each nonzero $\gamma_{dd'}$ corresponds to a component $V_i = \{d, d'\}$ with 
$|V_i| = 2$. Higher-order interactions ($\gamma_3, \gamma_4, \ldots$ from 
$\infty$-supermodularity) would require larger $r$.

\paragraph{Approximation Guarantees.}

\begin{theorem}[Chen et al. 2025]
\label{thm:chen-approx}
For monotone nonnegative $r$-decomposable supermodular functions:
\begin{enumerate}
    \item \textbf{Cardinality-constrained maximization} (find $k$-subset maximizing $f$):
    \[
    O(n^{(r-1)/2})\text{-approximation in } O(n^{r+1}) \text{ time}
    \]
    \item \textbf{Connected $k$-subset maximization}:
    \[
    O(n^{(r-1)/2})\text{-approximation in } O(mn^{r+1}) \text{ time}
    \]
\end{enumerate}
\end{theorem}

For EBF with pairwise interactions ($r = 2$):
\begin{itemize}
    \item Approximation ratio: $O(\sqrt{n})$
    \item Runtime: $O(n^3)$ for cardinality constraints
\end{itemize}

This is significant: \emph{intervention portfolio optimization}---selecting $k$ 
dimensions to target for maximum welfare gain---is computationally tractable!

\paragraph{Two Complementary Algorithms.}

\citet{chen2025supermodular} achieve their bounds through two strategies:

\begin{enumerate}
    \item \textbf{Greedy Peeling} (top-down): Start with $S = V$, iteratively remove 
    element with smallest marginal contribution until $|S| = k$. Achieves 
    $O(n^{r-1}/k^{r-1})$ approximation.
    
    \item \textbf{Batch-Greedy Augmenting} (bottom-up): Start with $S = \emptyset$, 
    iteratively add $r$-subset with maximum value increase until $|S| \geq k$. 
    Achieves $O(k^{r-1})$ approximation.
\end{enumerate}

The combination achieves $\min\{k^{r-1}, n^{r-1}/k^{r-1}\} \leq n^{(r-1)/2}$.

\begin{axiom}[Computational Tractability]
\label{ax:B-COMP}
\textbf{(B-COMP)} For $r$-decomposable supermodular utility functions, constrained 
welfare maximization admits polynomial-time $O(n^{(r-1)/2})$-approximation algorithms:
\begin{equation}
    \text{Pairwise } \gamma \text{ (}r=2\text{)} \Rightarrow O(\sqrt{n})\text{-approximation}
\end{equation}
\end{axiom}

\paragraph{Implications for Intervention Design.}

The decomposability parameter $r$ has direct policy implications:

\begin{center}
\begin{tabular}{lll}
\toprule
\textbf{Structure} & \textbf{r} & \textbf{Approximation} \\
\midrule
Independent utilities & 1 & Exact (select top-$k$) \\
Pairwise $\gamma_{dd'}$ & 2 & $O(\sqrt{n})$ \\
Three-way interactions & 3 & $O(n)$ \\
Full FEPSDE (6 dims) & $\leq 6$ & $O(n^{2.5})$ \\
\bottomrule
\end{tabular}
\end{center}

This suggests a trade-off: richer complementarity models (larger $r$) provide better 
behavioral predictions but make optimization harder. The EBF framework with 
primarily pairwise $\gamma$ interactions ($r = 2$) occupies a ``sweet spot'' of 
expressive power and computational tractability.



\paragraph{Level 1: Inter-Category Complementarity.}
How do the three utility \emph{categories} (INU, KNU, IDN) interact with each other?

\begin{equation}
Q = \sum_{c} \omega_c \cdot U^c + \sum_{c < c'} \gamma_{cc'} \cdot U^c \cdot U^{c'}
\end{equation}

where $c, c' \in \{INU, KNU, IDN\}$. This yields 3 unique pairs:
$\gamma_{INU,KNU}$, $\gamma_{INU,IDN}$, $\gamma_{KNU,IDN}$.

\paragraph{Level 2: Intra-Category Complementarity.}
Within each category, how do the six FEPSDE dimensions interact?

\begin{equation}
U^c = \sum_{d} u^c_d + \sum_{d < d'} \gamma^c_{dd'} \cdot u^c_d \cdot u^c_{d'}
\end{equation}

where $d, d' \in \{F, E, P, S, D, Ec\}$ and $c \in \{INU, KNU, IDN\}$.
This yields $3 \times 15 = 45$ category-specific dimension pairs.

\paragraph{Level 3: Cross-Category-Dimension Complementarity.}
How does a dimension in one category interact with the same (or different) 
dimension in another category?

\begin{equation}
\gamma_{d^c, d'^{c'}} \text{ where } c \neq c'
\end{equation}

For example: How does \emph{my} financial utility ($F^{INU}$) interact with 
the \emph{group's} financial utility ($F^{KNU}$)?

%==============================================================================
\subsubsection{Inter-Category Complementarity: INU $\times$ KNU $\times$ IDN}
\label{sec:inter-category}

\paragraph{The Three Inter-Category Coefficients.}

\begin{center}
\renewcommand{\arraystretch}{1.4}
\begin{tabular}{|l|c|c|c|}
\hline
& \textbf{INU} & \textbf{KNU} & \textbf{IDN} \\
\hline
\textbf{INU} & -- & $\gamma_{INU,KNU}$ & $\gamma_{INU,IDN}$ \\
\textbf{KNU} & $\gamma_{INU,KNU}$ & -- & $\gamma_{KNU,IDN}$ \\
\textbf{IDN} & $\gamma_{INU,IDN}$ & $\gamma_{KNU,IDN}$ & -- \\
\hline
\end{tabular}
\end{center}

%--- INU x KNU ---

\paragraph{$\gamma_{INU,KNU}$: Individual $\times$ Collective Complementarity.}

\begin{tcolorbox}[colback=yellow!5, colframe=yellow!50!black]
\textbf{Sign:} Context-dependent (critical for intervention design) \\
\textbf{Typical Range:} $\gamma_{INU,KNU} \in [-0.5, +0.5]$
\end{tcolorbox}

\textbf{The Question:} Does individual gain enhance or diminish collective welfare?

\textbf{When Complementary ($\gamma_{INU,KNU} > 0$):}
\begin{itemize}
    \item \textit{Team bonus structures}: My success contributes to team success
    \item \textit{Family prosperity}: My income supports family welfare
    \item \textit{Cooperative ventures}: Individual effort benefits the group
    \item \textit{High-trust public goods}: My contribution is valued by the collective
    \item \textit{Aligned incentives}: What's good for me is good for us
\end{itemize}

\textbf{When Substitutes ($\gamma_{INU,KNU} < 0$):}
\begin{itemize}
    \item \textit{Zero-sum competition}: My gain is your loss
    \item \textit{Free-riding}: I benefit while the group pays
    \item \textit{Rent-seeking}: Individual extraction from collective
    \item \textit{Positional goods}: My status depends on your lower status
    \item \textit{Tragedy of the commons}: Individual optimization harms collective
\end{itemize}

\textbf{Key Insight.}
The sign of $\gamma_{INU,KNU}$ determines whether markets and competition 
are beneficial or harmful. When $\gamma_{INU,KNU} > 0$, Adam Smith's 
``invisible hand'' operates; when $\gamma_{INU,KNU} < 0$, collective action 
problems arise.

\textbf{Context Dependence.}
\begin{equation}
\gamma_{INU,KNU}(\Psi) = f(\text{institutions}, \text{culture}, \text{incentive design})
\end{equation}

\textbf{Axiom.}
\begin{axiom}[Individual-Collective Complementarity Variability]
\label{ax:B-IK}
\textbf{(B-IK)} The sign of $\gamma_{INU,KNU}$ is not fixed but depends on 
institutional design:
\begin{equation}
    \text{sign}(\gamma_{INU,KNU}) = f(\text{rules of the game})
\end{equation}
Good institutions make $\gamma_{INU,KNU} > 0$; poor institutions make 
$\gamma_{INU,KNU} < 0$.
\end{axiom}

%--- INU x IDN ---

\paragraph{$\gamma_{INU,IDN}$: Individual $\times$ Identity Complementarity.}

\begin{tcolorbox}[colback=yellow!5, colframe=yellow!50!black]
\textbf{Sign:} Context-dependent (crucial for wellbeing) \\
\textbf{Typical Range:} $\gamma_{INU,IDN} \in [-0.4, +0.6]$
\end{tcolorbox}

\textbf{The Question:} Does individual success confirm or violate my identity?

\textbf{When Complementary ($\gamma_{INU,IDN} > 0$):}
\begin{itemize}
    \item \textit{Self-congruent success}: Achievement aligns with self-image
    \item \textit{Authentic pursuit}: Goals match values
    \item \textit{Identity-confirming consumption}: Purchases express who I am
    \item \textit{Vocation}: Work as expression of identity
\end{itemize}

\textbf{When Substitutes ($\gamma_{INU,IDN} < 0$):}
\begin{itemize}
    \item \textit{``Selling out''}: Material gain at cost of authenticity
    \item \textit{Identity violation}: Success requires becoming someone else
    \item \textit{Moral compromise}: Individual gain through unethical means
    \item \textit{Status-identity conflict}: Social climbing conflicts with origins
\end{itemize}

\textbf{Key Insight.}
High income combined with identity violation can reduce overall wellbeing.
This explains why some people reject higher-paying jobs that conflict with 
their self-image \citep{akerlofkranton}.

\textbf{Axiom.}
\begin{axiom}[Individual-Identity Alignment]
\label{ax:B-II}
\textbf{(B-II)} Sustainable wellbeing requires $\gamma_{INU,IDN} > 0$:
\begin{equation}
    \frac{\partial Q}{\partial t} > 0 \text{ (long-run)} \Rightarrow \gamma_{INU,IDN} > 0
\end{equation}
Individual success that violates identity is ultimately self-undermining.
\end{axiom}

%--- KNU x IDN ---

\paragraph{$\gamma_{KNU,IDN}$: Collective $\times$ Identity Complementarity.}

\begin{tcolorbox}[colback=yellow!5, colframe=yellow!50!black]
\textbf{Sign:} Context-dependent (critical for group dynamics) \\
\textbf{Typical Range:} $\gamma_{KNU,IDN} \in [-0.3, +0.7]$
\end{tcolorbox}

\textbf{The Question:} Does group welfare align with my identity?

\textbf{When Complementary ($\gamma_{KNU,IDN} > 0$):}
\begin{itemize}
    \item \textit{In-group identification}: I am proud when my group succeeds
    \item \textit{Collective identity}: My identity is defined by group membership
    \item \textit{Shared fate}: What happens to us happens to me
    \item \textit{Team sports}: Vicarious joy from team victory
\end{itemize}

\textbf{When Substitutes ($\gamma_{KNU,IDN} < 0$):}
\begin{itemize}
    \item \textit{Out-group identity}: I define myself against this group
    \item \textit{Counter-identity}: Group success threatens my distinctiveness
    \item \textit{Oppositional culture}: Mainstream success = identity betrayal
    \item \textit{Negative reference group}: I am what they are not
\end{itemize}

\textbf{Key Insight.}
Social identity theory \citep{tajfel1979integrative} explains why 
$\gamma_{KNU,IDN}$ can be strongly positive for in-groups but negative 
for out-groups. This has major implications for polarization.

%==============================================================================
\subsubsection{Category-Specific FEPSDE Matrices}
\label{sec:category-specific-fepsde}

\paragraph{The Key Question.}
Is the FEPSDE complementarity matrix the same across categories, or does 
each category have its own structure?

\begin{equation}
\gamma^{INU}_{dd'} \stackrel{?}{=} \gamma^{KNU}_{dd'} \stackrel{?}{=} \gamma^{IDN}_{dd'}
\end{equation}

\paragraph{Hypothesis: Similarity with Systematic Differences.}
We hypothesize that the \emph{sign} of most complementarities is stable 
across categories, but the \emph{magnitude} varies:

\begin{equation}
\text{sign}(\gamma^{INU}_{dd'}) = \text{sign}(\gamma^{KNU}_{dd'}) = \text{sign}(\gamma^{IDN}_{dd'})
\quad \text{(for most pairs)}
\end{equation}

but:

\begin{equation}
|\gamma^{INU}_{dd'}| \neq |\gamma^{KNU}_{dd'}| \neq |\gamma^{IDN}_{dd'}|
\quad \text{(magnitudes differ)}
\end{equation}

\paragraph{Example: Emotional $\times$ Physical ($\gamma_{EP}$).}

\begin{center}
\renewcommand{\arraystretch}{1.3}
\begin{tabular}{|l|c|l|}
\hline
\textbf{Category} & \textbf{$\gamma_{EP}$} & \textbf{Interpretation} \\
\hline
INU & $+0.6$ & Direct mind-body connection in individual \\
KNU & $+0.4$ & Group emotional \& physical health (less direct) \\
IDN & $+0.5$ & Identity as ``healthy person'' (self-image mediated) \\
\hline
\end{tabular}
\end{center}

All three are positive (sign stability), but INU shows the strongest 
coupling because the individual mind-body connection is most direct.

\paragraph{Default Assumption.}
For tractability, we adopt the FEPSDE matrix developed in the previous 
section as the \emph{default} for all categories:

\begin{equation}
\gamma_{dd'} \approx \gamma^{INU}_{dd'} \approx \gamma^{KNU}_{dd'} \approx \gamma^{IDN}_{dd'}
\end{equation}

Category-specific deviations should be noted where empirically important.

%==============================================================================
\subsubsection{Cross-Category-Dimension Complementarity}
\label{sec:cross-category-dimension}

\paragraph{The Additional Complexity.}
Beyond inter-category and intra-category complementarity, there is a third 
type: how does dimension $d$ in category $c$ interact with dimension $d'$ 
in category $c'$?

\paragraph{Same-Dimension Cross-Category.}
The most important case: $\gamma_{d^{INU}, d^{KNU}}$ for the same dimension $d$.

\textbf{Example: Financial.}
\begin{itemize}
    \item $\gamma_{F^{INU}, F^{KNU}}$: My financial utility $\times$ Group financial utility
    \item When $> 0$: My wealth and group wealth reinforce (shared prosperity)
    \item When $< 0$: My wealth and group wealth compete (zero-sum)
\end{itemize}

\textbf{Example: Social.}
\begin{itemize}
    \item $\gamma_{S^{INU}, S^{KNU}}$: My social connections $\times$ Group social cohesion
    \item Typically $> 0$: Individual social capital contributes to collective social capital
\end{itemize}

\paragraph{Simplifying Assumption.}
For most applications, we can assume that same-dimension cross-category 
complementarity is captured by the inter-category coefficient:

\begin{equation}
\gamma_{d^{INU}, d^{KNU}} \approx \gamma_{INU,KNU} \cdot \delta_{dd'}
\end{equation}

This reduces the 108 potential cross-category-dimension pairs to the 3 
inter-category coefficients.




%==============================================================================
% PART II EXTENSION: STRATEGIC HETEROGENEITY AND MIXED-SIGN COMPLEMENTARITIES
%==============================================================================

\subsubsection{Beyond Pure Supermodularity: Strategic Heterogeneity}
\label{sec:strategic-heterogeneity}

The FEPSDE complementarity matrix developed above reveals a crucial empirical 
fact: not all interactions are positive. Some dimension pairs show complements 
($\gamma > 0$), others show substitutes ($\gamma < 0$), and many are context-dependent 
(sign varies). This section provides the theoretical foundation for analyzing 
such \emph{mixed} systems.

\paragraph{The Problem with Pure Supermodularity.}
Classical supermodular game theory \citep{topkis1998} assumes that \emph{all} 
cross-partial derivatives are non-negative:

\begin{equation}
\frac{\partial^2 f}{\partial x_i \partial x_j} \geq 0 \quad \forall i \neq j
\end{equation}

This yields powerful results: existence of equilibrium via Tarski's theorem, 
monotone comparative statics, and predictable policy responses.

But real utility functions rarely satisfy this condition. In the EBF framework:
\begin{itemize}
    \item $\gamma_{FEc} < 0$ (Financial and Ecological are typically substitutes)
    \item $\gamma_{ED}$ varies by context (Digital-Emotional relationship is ambiguous)
    \item $\gamma_{INU,KNU}$ depends on institutional design
\end{itemize}

We need a theory that accommodates \emph{mixed} complementarities.

%------------------------------------------------------------------------------
\paragraph{Games of Strategic Heterogeneity (Barthel and Hoffman 2019).}

\citet{barthel2019rationalizability} introduce a general framework where players 
may exhibit \emph{either increasing or decreasing} best responses:

\begin{definition}[Strategic Heterogeneity]
\label{def:strategic-heterogeneity}
A game exhibits \textbf{strategic heterogeneity} if, for different pairs of 
players $(i,j)$, the best response of player $i$ to player $j$'s action may be:
\begin{itemize}
    \item \textbf{Increasing} (strategic complements): $\partial BR_i / \partial a_j > 0$
    \item \textbf{Decreasing} (strategic substitutes): $\partial BR_i / \partial a_j < 0$
\end{itemize}
\end{definition}

\textbf{Key Result.}
Even with mixed-monotone best responses, equilibrium existence and learning 
convergence results extend, provided the game can be ``monotonically embedded'' 
into a larger game with consistent structure.

\textbf{EBF Implication.}
The FEPSDE utility system is a \emph{game of strategic heterogeneity} at the 
individual level. Different dimension pairs interact differently:

\begin{center}
\renewcommand{\arraystretch}{1.3}
\begin{tabular}{|l|c|l|}
\hline
\textbf{Dimension Pair} & \textbf{Sign} & \textbf{Type} \\
\hline
Emotional $\times$ Physical & $+$ & Strategic Complements \\
Emotional $\times$ Social & $+$ & Strategic Complements \\
Financial $\times$ Ecological & $-$ & Strategic Substitutes \\
Financial $\times$ Emotional & $\pm$ & Heterogeneous \\
Digital $\times$ Social & $\pm$ & Heterogeneous \\
\hline
\end{tabular}
\end{center}

%------------------------------------------------------------------------------
\paragraph{Christensen's Condition: MCS with Heterogeneous Interactions.}

\citet{christensen2019comparative} provides a crucial result: monotone comparative 
statics (MCS) are possible even when the system is \emph{not} supermodular, 
under a specific averaging condition.

\begin{theorem}[Christensen's Condition for MCS]
\label{thm:christensen}
Consider a system of $n$ first-order conditions $\nabla f(x,t) = 0$ where $t$ 
is a parameter complementary to the endogenous variables. Monotone comparative 
statics hold if the Jacobian $D_x \nabla f(x,t)$ satisfies:

\begin{equation}
\boxed{
\frac{1}{n} \sum_{j=1}^{n} \frac{\partial^2 f_i}{\partial x_i \partial x_j} 
\leq \min\left\{ 0, \frac{\partial^2 f_i}{\partial x_i \partial x_j} \Big| j \neq i \right\}
\quad \text{for } i = 1, \ldots, n
}
\label{eq:christensen-condition}
\end{equation}
\end{theorem}

\textbf{Interpretation.}
The condition says: the \emph{average} interaction effect must be bounded by 
the \emph{minimum} interaction effect. This allows for:
\begin{itemize}
    \item Some positive interactions ($\gamma_{ij} > 0$)
    \item Some negative interactions ($\gamma_{ij} < 0$)
    \item As long as the average is appropriately bounded
\end{itemize}

\textbf{The Trade-off.}
Christensen identifies a fundamental trade-off:

\begin{quote}
``A trade-off between the \emph{heterogeneity} and \emph{magnitude} of 
endogenous interactions in maintaining MCS.''
\end{quote}

\begin{itemize}
    \item High heterogeneity (mix of $+$ and $-$) $\Rightarrow$ requires smaller magnitudes
    \item Low heterogeneity (all same sign) $\Rightarrow$ allows larger magnitudes
\end{itemize}

%------------------------------------------------------------------------------
\paragraph{Application to EBF: The Coherence-Heterogeneity Trade-off.}

This result has profound implications for EBF. Define:

\begin{definition}[Interaction Heterogeneity]
\label{def:heterogeneity}
The \textbf{heterogeneity} of a $\gamma$-matrix is:
\begin{equation}
H(\Gamma) = \text{Var}(\{\gamma_{ij}\}_{i<j}) = \frac{1}{|\{i<j\}|} \sum_{i<j} (\gamma_{ij} - \bar{\gamma})^2
\end{equation}
where $\bar{\gamma}$ is the mean complementarity coefficient.
\end{definition}

\begin{proposition}[Coherence-Heterogeneity Trade-off]
\label{prop:coherence-heterogeneity}
For a utility system with interaction heterogeneity $H(\Gamma)$:
\begin{equation}
K_{max} = K_{max}(H) \quad \text{where} \quad \frac{\partial K_{max}}{\partial H} < 0
\end{equation}
Higher heterogeneity in the $\gamma$-matrix reduces the maximum achievable 
coherence $K$.
\end{proposition}

\textbf{Intuition.}
When some dimensions are complements and others substitutes, perfect alignment 
becomes impossible. Optimizing for one pair may de-optimize another.

\textbf{Policy Implication.}
Interventions should:
\begin{enumerate}
    \item Identify the dominant sign of $\gamma$ in the relevant subsystem
    \item Design for coherence within that subsystem
    \item Accept trade-offs with dimensions of opposite sign
\end{enumerate}

%------------------------------------------------------------------------------
\paragraph{The Mixed-System Equilibrium.}

For systems with strategic heterogeneity, we cannot expect corner solutions 
(all high or all low). Instead, equilibrium involves \emph{interior} solutions 
that balance competing complementarities.

\begin{equation}
x^* = \arg\max_x \left[ \sum_i u_i(x_i) + \sum_{i<j} \gamma_{ij} x_i x_j \right]
\end{equation}

When $\gamma_{ij}$ have mixed signs, the solution satisfies:

\begin{equation}
\frac{\partial u_i}{\partial x_i} + \sum_{j \neq i} \gamma_{ij} x_j = 0
\quad \forall i
\end{equation}

This system generally has interior solutions, not the corner solutions 
(``all high'' or ``all low'') typical of pure supermodular systems.

%------------------------------------------------------------------------------
\paragraph{Axioms for Strategic Heterogeneity.}

\begin{axiom}[Mixed Complementarity]
\label{ax:B-MIX}
\textbf{(B-MIX)} The FEPSDE utility system exhibits strategic heterogeneity:
\begin{equation}
\exists\, (d, d') : \gamma_{dd'} > 0 \quad \text{and} \quad 
\exists\, (d'', d''') : \gamma_{d''d'''} < 0
\end{equation}
Not all dimension pairs are complements; not all are substitutes.
\end{axiom}

\begin{axiom}[Bounded Heterogeneity]
\label{ax:B-BND}
\textbf{(B-BND)} The heterogeneity of the $\gamma$-matrix is bounded such that 
Christensen's condition holds:
\begin{equation}
H(\Gamma) < H_{crit} \quad \text{where } H_{crit} \text{ ensures MCS}
\end{equation}
\end{axiom}

\begin{axiom}[Sign Stability]
\label{ax:B-SIGN}
\textbf{(B-SIGN)} While magnitudes vary with context, the \emph{sign} of most 
$\gamma_{dd'}$ is stable:
\begin{equation}
\text{sign}(\gamma_{dd'}(\Psi_1)) = \text{sign}(\gamma_{dd'}(\Psi_2)) 
\quad \text{for most } \Psi_1, \Psi_2, d, d'
\end{equation}
Exceptions (sign-switching pairs) are identified and modeled separately.
\end{axiom}

%------------------------------------------------------------------------------
\paragraph{Literature on Mixed Systems.}

\begin{center}
\renewcommand{\arraystretch}{1.3}
\begin{tabular}{|l|l|l|}
\hline
\textbf{Reference} & \textbf{Contribution} & \textbf{EBF Relevance} \\
\hline
\citet{barthel2019rationalizability} & Strategic heterogeneity & Mixed $\gamma$ signs \\
\citet{christensen2019comparative} & MCS with heterogeneity & Condition for $K$ \\
\citet{roy2012characterizing} & Strategic substitutes & Negative $\gamma$ \\
\citet{monaco2016games} & Complements + substitutes & General framework \\
\citet{amir2019supermodularity} & Special issue overview & State of the art \\
\hline
\end{tabular}
\end{center}



\subsubsection{The FEPSDE Dimensions (Recap from CORE-WHAT)}

\begin{center}
\renewcommand{\arraystretch}{1.3}
\begin{tabular}{|c|l|l|l|}
\hline
\textbf{Code} & \textbf{Dimension} & \textbf{Category} & \textbf{Core Concern} \\
\hline
F & Financial & INU (Individual) & Material resources, wealth \\
E & Emotional & INU (Individual) & Psychological wellbeing, affect \\
P & Physical & INU (Individual) & Health, bodily functioning \\
S & Social & KNU (Collective) & Relationships, belonging \\
D & Digital & KNU (Collective) & Information, connectivity \\
Ec & Ecological & IDN (Identity) & Environmental sustainability \\
\hline
\end{tabular}
\end{center}

\subsubsection{The Complete 6×6 $\gamma$-Matrix}

\begin{equation}
\Gamma_{FEPSDE} = \begin{pmatrix}
-- & \gamma_{FE} & \gamma_{FP} & \gamma_{FS} & \gamma_{FD} & \gamma_{FEc} \\
\gamma_{FE} & -- & \gamma_{EP} & \gamma_{ES} & \gamma_{ED} & \gamma_{EEc} \\
\gamma_{FP} & \gamma_{EP} & -- & \gamma_{PS} & \gamma_{PD} & \gamma_{PEc} \\
\gamma_{FS} & \gamma_{ES} & \gamma_{PS} & -- & \gamma_{SD} & \gamma_{SEc} \\
\gamma_{FD} & \gamma_{ED} & \gamma_{PD} & \gamma_{SD} & -- & \gamma_{DEc} \\
\gamma_{FEc} & \gamma_{EEc} & \gamma_{PEc} & \gamma_{SEc} & \gamma_{DEc} & -- \\
\end{pmatrix}
\label{eq:gamma-fepsde-matrix}
\end{equation}

The diagonal is undefined (a dimension's complementarity with itself is not 
meaningful in this framework). By symmetry, we need to specify 15 coefficients.

%------------------------------------------------------------------------------
% FINANCIAL INTERACTIONS (5 pairs)
%------------------------------------------------------------------------------

\subsubsection{Financial × Other Dimensions}

\paragraph{$\gamma_{FE}$: Financial × Emotional}

\begin{tcolorbox}[colback=yellow!5, colframe=yellow!50!black]
\textbf{Sign:} Ambiguous / Context-dependent \\
\textbf{Typical Range:} $\gamma_{FE} \in [-0.3, +0.3]$
\end{tcolorbox}

\textbf{Theory.}
The relationship between money and happiness is one of the most studied in 
behavioral economics. Key findings:
\begin{itemize}
    \item \textit{Easterlin Paradox} \citep{easterlin1974}: Beyond a threshold, 
          income gains do not increase happiness proportionally
    \item \textit{Relative Income Hypothesis}: Emotional utility depends on 
          income relative to reference group, not absolute level
    \item \textit{Hedonic Adaptation}: People adapt to income changes, 
          reducing long-term emotional impact
\end{itemize}

\textbf{When Complementary ($\gamma_{FE} > 0$):}
\begin{itemize}
    \item Financial security reduces anxiety, enabling positive emotions
    \item Resources enable experiences that generate happiness (travel, leisure)
    \item Below subsistence threshold, money directly enables emotional stability
\end{itemize}

\textbf{When Substitutes ($\gamma_{FE} < 0$):}
\begin{itemize}
    \item Time-money tradeoff: earning more may reduce time for relationships
    \item Materialistic pursuit crowds out intrinsic satisfaction
    \item High income jobs often carry stress, reducing wellbeing
\end{itemize}

\textbf{Context Dependence.}
\begin{equation}
\gamma_{FE}(\Psi) = \begin{cases}
    > 0 & \text{if } F < F_{threshold} \text{ (poverty)} \\
    \approx 0 & \text{if } F \in [F_{threshold}, F_{satiation}] \\
    < 0 & \text{if } F > F_{satiation} \text{ and time-poor}
\end{cases}
\end{equation}

\textbf{Key References.}
\citet{kahneman2010high} find satiation around \$75,000/year (2010 USD).
\citet{killingsworth2021experienced} find continued increases but at 
diminishing rate. \citet{stevenson2008economic} challenge the Easterlin paradox.

%---

\paragraph{$\gamma_{FP}$: Financial × Physical}

\begin{tcolorbox}[colback=green!5, colframe=green!50!black]
\textbf{Sign:} Positive (generally) \\
\textbf{Typical Range:} $\gamma_{FP} \in [+0.2, +0.6]$
\end{tcolorbox}

\textbf{Theory.}
Financial resources enable health investments and health enables earning capacity.

\textbf{Mechanisms for Complementarity:}
\begin{itemize}
    \item Healthcare access: Money buys medical care, insurance, preventive services
    \item Nutrition: Higher income enables better food quality
    \item Living conditions: Better housing, lower pollution exposure
    \item Leisure time: Resources enable rest, exercise, recovery
    \item Reverse causation: Health enables work, earning, wealth accumulation
\end{itemize}

\textbf{Empirical Evidence.}
The income-health gradient is one of the most robust findings in social 
epidemiology \citep{marmot2004status}. Each step up the socioeconomic ladder 
corresponds to better health outcomes.

\textbf{Context Modulation.}
\begin{itemize}
    \item Healthcare system: $\gamma_{FP}$ higher in US (private) than UK (NHS)
    \item Age: $\gamma_{FP}$ increases with age (health costs rise)
    \item Occupation: Manual labor shows different $\gamma_{FP}$ than knowledge work
\end{itemize}

\textbf{Axiom.}
\begin{axiom}[Financial-Physical Complementarity]
\label{ax:B-FP}
\textbf{(B-FP)} Financial and Physical utility are complements at the individual level:
\begin{equation}
    \gamma^{L=0}_{FP} > 0 \quad \text{(default assumption)}
\end{equation}
\end{axiom}

%---

\paragraph{$\gamma_{FS}$: Financial × Social}

\begin{tcolorbox}[colback=yellow!5, colframe=yellow!50!black]
\textbf{Sign:} Ambiguous / Context-dependent \\
\textbf{Typical Range:} $\gamma_{FS} \in [-0.2, +0.4]$
\end{tcolorbox}

\textbf{Theory.}
Money can both enable and undermine social relationships.

\textbf{When Complementary ($\gamma_{FS} > 0$):}
\begin{itemize}
    \item Resources enable social participation (clubs, events, hospitality)
    \item Financial stability reduces relationship stress
    \item Shared resources strengthen family/household bonds
    \item Philanthropy creates social connections
\end{itemize}

\textbf{When Substitutes ($\gamma_{FS} < 0$):}
\begin{itemize}
    \item Time-money tradeoff reduces time for relationships
    \item Wealth inequality creates social distance
    \item Monetizing relationships (paying for care) may erode bonds
    \item Competition for resources strains relationships
\end{itemize}

\textbf{Level Dependence.}
\begin{equation}
\gamma^L_{FS} = \begin{cases}
    \text{ambiguous} & L = 0 \text{ (individual)} \\
    > 0 & L = 1 \text{ (household: shared resources)} \\
    > 0 & L = 2 \text{ (organization: compensation \& culture)}
\end{cases}
\end{equation}

\textbf{Key References.}
\citet{vohs2006psychological} on money priming reducing prosociality.
\citet{aknin2013prosocial} on spending on others increasing happiness.

%---

\paragraph{$\gamma_{FD}$: Financial × Digital}

\begin{tcolorbox}[colback=green!5, colframe=green!50!black]
\textbf{Sign:} Positive (strong) \\
\textbf{Typical Range:} $\gamma_{FD} \in [+0.3, +0.7]$
\end{tcolorbox}

\textbf{Theory.}
In the modern economy, digital capability and financial resources are 
strongly complementary.

\textbf{Mechanisms:}
\begin{itemize}
    \item Digital access requires investment (devices, connectivity)
    \item Digital skills enhance earning capacity (labor market premium)
    \item FinTech: Digital tools enable financial management
    \item E-commerce: Digital access expands consumption options
    \item Information asymmetry: Digital literacy improves financial decisions
\end{itemize}

\textbf{Digital Divide Implication.}
Low $F$ $\rightarrow$ low $D$ access $\rightarrow$ lower returns to $F$ 
investment $\rightarrow$ poverty trap.

\textbf{Axiom.}
\begin{axiom}[Financial-Digital Complementarity]
\label{ax:B-FD}
\textbf{(B-FD)} Financial and Digital utility are strong complements in 
the contemporary economy:
\begin{equation}
    \gamma_{FD} > 0 \quad \text{and increasing over time}
\end{equation}
\end{axiom}

%---

\paragraph{$\gamma_{FEc}$: Financial × Ecological}

\begin{tcolorbox}[colback=red!5, colframe=red!50!black]
\textbf{Sign:} Negative (typically) / Context-dependent \\
\textbf{Typical Range:} $\gamma_{FEc} \in [-0.4, +0.2]$
\end{tcolorbox}

\textbf{Theory.}
The relationship between economic activity and environmental sustainability 
is contentious.

\textbf{When Substitutes ($\gamma_{FEc} < 0$):}
\begin{itemize}
    \item Production externalities: Economic output generates pollution
    \item Consumption footprint: Higher income $\rightarrow$ higher resource use
    \item Growth imperative conflicts with planetary boundaries
    \item Discount rate: Financial optimization undervalues future environment
\end{itemize}

\textbf{When Complementary ($\gamma_{FEc} > 0$):}
\begin{itemize}
    \item Environmental Kuznets Curve: Beyond threshold, wealth enables cleanup
    \item Green investment: Resources enable renewable transition
    \item Ecosystem services: Environmental quality supports economic activity
    \item Regulation capacity: Wealth enables environmental governance
\end{itemize}

\textbf{The Decoupling Question.}
Can economic growth be ``decoupled'' from environmental impact? If yes, 
$\gamma_{FEc}$ could become positive. Current evidence is mixed 
\citep{hickel2020decoupling}.

\textbf{Context Dependence.}
\begin{equation}
\gamma_{FEc}(\Psi) = f(\text{technology}, \text{regulation}, \text{values})
\end{equation}

%------------------------------------------------------------------------------
% EMOTIONAL INTERACTIONS (4 remaining pairs)
%------------------------------------------------------------------------------

\subsubsection{Emotional × Other Dimensions}

\paragraph{$\gamma_{EP}$: Emotional × Physical}

\begin{tcolorbox}[colback=green!5, colframe=green!50!black]
\textbf{Sign:} Positive (strong) \\
\textbf{Typical Range:} $\gamma_{EP} \in [+0.4, +0.7]$
\end{tcolorbox}

\textbf{Theory.}
The mind-body connection is one of the strongest complementarities in the 
FEPSDE system.

\textbf{Mechanisms:}
\begin{itemize}
    \item Psychosomatic effects: Emotional states affect physical health
    \item Exercise-mood link: Physical activity improves emotional wellbeing
    \item Stress-health pathway: Chronic stress impairs immune function
    \item Pain-affect interaction: Physical pain causes emotional distress
    \item Sleep: Bidirectional relationship with both E and P
\end{itemize}

\textbf{Empirical Evidence.}
\begin{itemize}
    \item Depression increases cardiovascular disease risk by 80\% \citep{nicholson2006depression}
    \item Exercise is as effective as antidepressants for mild-moderate depression \citep{blumenthal2007exercise}
    \item Positive emotions improve immune function \citep{cohen2003emotional}
\end{itemize}

\textbf{Axiom.}
\begin{axiom}[Emotional-Physical Complementarity]
\label{ax:B-EP}
\textbf{(B-EP)} Emotional and Physical utility are strong complements:
\begin{equation}
    \gamma_{EP} > 0 \quad \text{(robust across contexts)}
\end{equation}
This is one of the most stable complementarities in the framework.
\end{axiom}

%---

\paragraph{$\gamma_{ES}$: Emotional × Social}

\begin{tcolorbox}[colback=green!5, colframe=green!50!black]
\textbf{Sign:} Positive (strong) \\
\textbf{Typical Range:} $\gamma_{ES} \in [+0.5, +0.8]$
\end{tcolorbox}

\textbf{Theory.}
Humans are social beings; emotional wellbeing is deeply tied to social connection.

\textbf{Mechanisms:}
\begin{itemize}
    \item Attachment theory: Social bonds provide emotional security
    \item Social support: Relationships buffer against emotional distress
    \item Loneliness: Social isolation is a major risk factor for depression
    \item Emotional contagion: Emotions spread through social networks
    \item Identity: Social belonging contributes to self-worth
\end{itemize}

\textbf{Empirical Evidence.}
\begin{itemize}
    \item Loneliness is as harmful as smoking 15 cigarettes/day \citep{holt2010social}
    \item Social relationships are the strongest predictor of happiness \citep{helliwell2012world}
    \item Social support predicts recovery from depression \citep{george1989social}
\end{itemize}

\textbf{Axiom.}
\begin{axiom}[Emotional-Social Complementarity]
\label{ax:B-ES}
\textbf{(B-ES)} Emotional and Social utility are strong complements:
\begin{equation}
    \gamma_{ES} > 0 \quad \text{(one of the strongest in FEPSDE)}
\end{equation}
\end{axiom}

%---

\paragraph{$\gamma_{ED}$: Emotional × Digital}

\begin{tcolorbox}[colback=yellow!5, colframe=yellow!50!black]
\textbf{Sign:} Ambiguous / Context-dependent \\
\textbf{Typical Range:} $\gamma_{ED} \in [-0.3, +0.3]$
\end{tcolorbox}

\textbf{Theory.}
Digital technology has complex effects on emotional wellbeing.

\textbf{When Complementary ($\gamma_{ED} > 0$):}
\begin{itemize}
    \item Connection: Digital tools enable social contact (especially during isolation)
    \item Information: Access to mental health resources
    \item Entertainment: Digital media provides emotional experiences
    \item Expression: Platforms for emotional sharing
\end{itemize}

\textbf{When Substitutes ($\gamma_{ED} < 0$):}
\begin{itemize}
    \item Social comparison: Social media induces envy, inadequacy
    \item Addiction: Compulsive use undermines emotional regulation
    \item Displacement: Screen time crowds out in-person connection
    \item Misinformation: Exposure to distressing content
    \item Sleep disruption: Blue light, late-night use
\end{itemize}

\textbf{Age Moderation.}
\begin{equation}
\gamma_{ED}(\text{age}) = \begin{cases}
    < 0 & \text{adolescents (vulnerable)} \\
    \approx 0 & \text{adults (mixed)} \\
    > 0 & \text{elderly (connection value)}
\end{cases}
\end{equation}

\textbf{Key References.}
\citet{twenge2018increases} on social media and teen depression.
\citet{orben2019association} finding small effect sizes.

%---

\paragraph{$\gamma_{EEc}$: Emotional × Ecological}

\begin{tcolorbox}[colback=green!5, colframe=green!50!black]
\textbf{Sign:} Positive \\
\textbf{Typical Range:} $\gamma_{EEc} \in [+0.2, +0.5]$
\end{tcolorbox}

\textbf{Theory.}
Connection to nature supports emotional wellbeing.

\textbf{Mechanisms:}
\begin{itemize}
    \item Biophilia: Innate human affinity for natural environments
    \item Stress reduction: Nature exposure lowers cortisol
    \item Eco-anxiety: Environmental degradation causes emotional distress
    \item Meaning: Environmental stewardship provides purpose
    \item Aesthetics: Natural beauty generates positive emotions
\end{itemize}

\textbf{Empirical Evidence.}
\begin{itemize}
    \item 20 minutes in nature significantly reduces cortisol \citep{hunter2019urban}
    \item ``Forest bathing'' improves mood and reduces anxiety \citep{li2010effect}
    \item Climate anxiety is rising, especially among youth \citep{hickman2021climate}
\end{itemize}

%------------------------------------------------------------------------------
% PHYSICAL INTERACTIONS (3 remaining pairs)
%------------------------------------------------------------------------------

\subsubsection{Physical × Other Dimensions}

\paragraph{$\gamma_{PS}$: Physical × Social}

\begin{tcolorbox}[colback=green!5, colframe=green!50!black]
\textbf{Sign:} Positive \\
\textbf{Typical Range:} $\gamma_{PS} \in [+0.2, +0.5]$
\end{tcolorbox}

\textbf{Theory.}
Physical health and social connection are mutually reinforcing.

\textbf{Mechanisms:}
\begin{itemize}
    \item Social activities often involve physical activity (sports, walking)
    \item Social support improves health behaviors and outcomes
    \item Illness can lead to social isolation
    \item Caregiving: Physical needs require social support
    \item Public health: Social norms influence health behaviors
\end{itemize}

\textbf{Empirical Evidence.}
Social integration predicts mortality risk, controlling for health behaviors 
\citep{holt2010social}.

%---

\paragraph{$\gamma_{PD}$: Physical × Digital}

\begin{tcolorbox}[colback=yellow!5, colframe=yellow!50!black]
\textbf{Sign:} Ambiguous \\
\textbf{Typical Range:} $\gamma_{PD} \in [-0.2, +0.3]$
\end{tcolorbox}

\textbf{Theory.}
Digital technology has mixed effects on physical health.

\textbf{When Complementary:}
\begin{itemize}
    \item Health tracking: Wearables, apps for fitness and health monitoring
    \item Telemedicine: Digital access to healthcare
    \item Health information: Online resources for self-care
    \item Remote work: Reduced commuting stress (for some)
\end{itemize}

\textbf{When Substitutes:}
\begin{itemize}
    \item Sedentary behavior: Screen time displaces physical activity
    \item Sleep disruption: Device use before bed
    \item Ergonomic issues: Posture, eye strain, repetitive stress
\end{itemize}

%---

\paragraph{$\gamma_{PEc}$: Physical × Ecological}

\begin{tcolorbox}[colback=green!5, colframe=green!50!black]
\textbf{Sign:} Positive (strong) \\
\textbf{Typical Range:} $\gamma_{PEc} \in [+0.3, +0.6]$
\end{tcolorbox}

\textbf{Theory.}
Environmental quality directly affects physical health.

\textbf{Mechanisms:}
\begin{itemize}
    \item Air quality: Pollution causes respiratory and cardiovascular disease
    \item Water quality: Clean water essential for health
    \item Climate: Extreme heat, disease vectors, food security
    \item Green space: Access to nature enables physical activity
    \item Toxins: Environmental contamination causes illness
\end{itemize}

\textbf{Axiom.}
\begin{axiom}[Physical-Ecological Complementarity]
\label{ax:B-PEc}
\textbf{(B-PEc)} Physical and Ecological utility are complements:
\begin{equation}
    \gamma_{PEc} > 0 \quad \text{(environmental health = human health)}
\end{equation}
\end{axiom}

%------------------------------------------------------------------------------
% SOCIAL INTERACTIONS (2 remaining pairs)
%------------------------------------------------------------------------------

\subsubsection{Social × Other Dimensions}

\paragraph{$\gamma_{SD}$: Social × Digital}

\begin{tcolorbox}[colback=yellow!5, colframe=yellow!50!black]
\textbf{Sign:} Ambiguous / Rapidly evolving \\
\textbf{Typical Range:} $\gamma_{SD} \in [-0.2, +0.5]$
\end{tcolorbox}

\textbf{Theory.}
Digital technology transforms social interaction.

\textbf{When Complementary:}
\begin{itemize}
    \item Connectivity: Maintain relationships across distance
    \item Community: Online groups for niche interests
    \item Coordination: Digital tools enable collective action
    \item Information sharing: Social learning at scale
\end{itemize}

\textbf{When Substitutes:}
\begin{itemize}
    \item Displacement: Online interaction crowds out in-person
    \item Quality: Digital interactions may be shallower
    \item Polarization: Filter bubbles fragment social cohesion
    \item Anonymity: Online disinhibition, harassment
\end{itemize}

\textbf{Generational Shift.}
For digital natives, $\gamma_{SD}$ may be structurally different than for 
older cohorts.

%---

\paragraph{$\gamma_{SEc}$: Social × Ecological}

\begin{tcolorbox}[colback=green!5, colframe=green!50!black]
\textbf{Sign:} Positive \\
\textbf{Typical Range:} $\gamma_{SEc} \in [+0.2, +0.4]$
\end{tcolorbox}

\textbf{Theory.}
Environmental stewardship is fundamentally a collective action problem.

\textbf{Mechanisms:}
\begin{itemize}
    \item Commons management: Social norms govern resource use
    \item Collective action: Environmental protection requires coordination
    \item Community resilience: Social capital enables climate adaptation
    \item Environmental justice: Social structures determine exposure
    \item Shared identity: ``We'' framing enables sustainability action
\end{itemize}

\textbf{Key Reference.}
\citet{ostrom1990governing} on social solutions to environmental commons.

%------------------------------------------------------------------------------
% DIGITAL × ECOLOGICAL (final pair)
%------------------------------------------------------------------------------

\subsubsection{Digital × Ecological}

\paragraph{$\gamma_{DEc}$: Digital × Ecological}

\begin{tcolorbox}[colback=yellow!5, colframe=yellow!50!black]
\textbf{Sign:} Ambiguous \\
\textbf{Typical Range:} $\gamma_{DEc} \in [-0.3, +0.3]$
\end{tcolorbox}

\textbf{Theory.}
Digital technology has complex environmental implications.

\textbf{When Complementary:}
\begin{itemize}
    \item Efficiency: Smart systems reduce resource use
    \item Monitoring: Sensors enable environmental tracking
    \item Dematerialization: Digital goods replace physical
    \item Information: Awareness of environmental issues
    \item Coordination: Digital tools for environmental action
\end{itemize}

\textbf{When Substitutes:}
\begin{itemize}
    \item Energy use: Data centers, device manufacturing
    \item E-waste: Toxic disposal of electronics
    \item Rebound effects: Efficiency gains spent on more consumption
    \item Planned obsolescence: Short device lifespans
\end{itemize}

%------------------------------------------------------------------------------
% SUMMARY TABLE
%------------------------------------------------------------------------------

\subsubsection{Summary: The Complete $\gamma_{dd'}$ Matrix}

\begin{center}
\renewcommand{\arraystretch}{1.4}
\begin{tabular}{|c|c|c|c|c|c|c|}
\hline
$\gamma_{dd'}$ & \textbf{F} & \textbf{E} & \textbf{P} & \textbf{S} & \textbf{D} & \textbf{Ec} \\
\hline
\textbf{F} & -- & $\pm$ & \cellcolor{green!20}+ & $\pm$ & \cellcolor{green!20}+ & \cellcolor{red!20}-- \\
\textbf{E} & $\pm$ & -- & \cellcolor{green!30}++ & \cellcolor{green!30}++ & $\pm$ & \cellcolor{green!20}+ \\
\textbf{P} & \cellcolor{green!20}+ & \cellcolor{green!30}++ & -- & \cellcolor{green!20}+ & $\pm$ & \cellcolor{green!20}+ \\
\textbf{S} & $\pm$ & \cellcolor{green!30}++ & \cellcolor{green!20}+ & -- & $\pm$ & \cellcolor{green!20}+ \\
\textbf{D} & \cellcolor{green!20}+ & $\pm$ & $\pm$ & $\pm$ & -- & $\pm$ \\
\textbf{Ec} & \cellcolor{red!20}-- & \cellcolor{green!20}+ & \cellcolor{green!20}+ & \cellcolor{green!20}+ & $\pm$ & -- \\
\hline
\end{tabular}
\end{center}

\textbf{Legend:}
\begin{itemize}
    \item \colorbox{green!30}{++} Strong complement (robust across contexts)
    \item \colorbox{green!20}{+} Complement (typical)
    \item $\pm$ Ambiguous / Context-dependent
    \item \colorbox{red!20}{--} Substitute (typical)
\end{itemize}

\subsubsection{Part II Axiom Summary}

\begin{center}
\renewcommand{\arraystretch}{1.3}
\begin{tabular}{|l|l|l|}
\hline
\textbf{Code} & \textbf{Axiom} & \textbf{Statement} \\
\hline
B-FP & Financial-Physical & $\gamma_{FP} > 0$ \\
B-FD & Financial-Digital & $\gamma_{FD} > 0$ (increasing) \\
B-EP & Emotional-Physical & $\gamma_{EP} > 0$ (robust) \\
B-ES & Emotional-Social & $\gamma_{ES} > 0$ (strongest) \\
B-PEc & Physical-Ecological & $\gamma_{PEc} > 0$ \\
\hline
\end{tabular}
\end{center}

The five axioms above represent the most stable complementarities in the 
FEPSDE system. The remaining 10 pairs are context-dependent and require 
calibration to specific populations and situations.



%%%%%%%%%%%%%%%%%%%%%%%%%%%%%%%%%%%%%%%%%%%%%%%%%%%%%%%%%%%%%%%%%%%%%%%%%%%%%%%
% END OF PART I
% Parts II-VIII to follow in subsequent sections
%%%%%%%%%%%%%%%%%%%%%%%%%%%%%%%%%%%%%%%%%%%%%%%%%%%%%%%%%%%%%%%%%%%%%%%%%%%%%%%


%------------------------------------------------------------------------------
\subsection{Teil III: Utility $\times$ Context---The Awareness Function}
\label{sec:utility-context}

While Teil II examined complementarities \emph{within} the utility function (how 
different welfare dimensions interact), Teil III addresses a fundamentally different 
question: how does \emph{context} interact with utility? This is the domain of the 
EBF Awareness function $A(\Psi)$.

\subsubsection{The Conceptual Framework}

\paragraph{Context as a Lattice.}
The 8$\Psi$-Dimensions framework defines context as an 8-dimensional vector:
\begin{equation}
\Psi = (\psi_1, \ldots, \psi_8) \in [0,1]^8
\end{equation}
representing Temporal Orientation, Uncertainty Tolerance, Social Embeddedness, 
Cognitive Load, Emotional Valence, Goal Accessibility, Resource Scarcity, and 
Identity Salience.

The context space $[0,1]^8$ forms a \emph{complete lattice} under the componentwise 
order: $\Psi \leq \Psi'$ iff $\psi_i \leq \psi'_i$ for all $i$. This lattice 
structure enables the application of supermodularity theory.

\paragraph{The Awareness Function.}
EBF defines Awareness as a mapping from context to ``effective utility weights'':
\begin{equation}
A: [0,1]^8 \to [0,1]^6, \quad \Psi \mapsto A(\Psi) = (a_F, a_E, a_P, a_S, a_D, a_{Ec})
\end{equation}
where $a_d(\Psi) \in [0,1]$ represents the degree to which dimension $d$ is 
``accessible'' or ``salient'' in context $\Psi$.

The effective utility becomes:
\begin{equation}
U^{eff}(\mathbf{d}; \Psi) = \sum_d a_d(\Psi) \cdot u_d + 
\sum_{d < d'} \gamma_{dd'} \cdot a_d(\Psi) \cdot a_{d'}(\Psi) \cdot u_d \cdot u_{d'}
\end{equation}

\paragraph{Supermodularity in Context.}
The key question: Under what conditions is $A(\Psi)$ supermodular in $\Psi$?

\begin{definition}[Context Supermodularity]
\label{def:context-super}
The awareness function $A$ exhibits \textbf{context supermodularity} if for each 
dimension $d$:
\begin{equation}
a_d(\Psi \vee \Psi') + a_d(\Psi \wedge \Psi') \geq a_d(\Psi) + a_d(\Psi')
\end{equation}
for all context pairs $\Psi, \Psi' \in [0,1]^8$.
\end{definition}

\textbf{Interpretation}: Context factors are \emph{complements} in producing awareness. 
High temporal orientation ($\psi_1$) and low cognitive load ($\psi_4$) together 
produce \emph{more} than additive awareness of future-oriented dimensions.

%------------------------------------------------------------------------------
\subsubsection{Theoretical Foundations: Monotone Comparative Statics}

The interaction between utility and context falls squarely within the domain of 
\emph{monotone comparative statics}---how optimal choices change with parameters.

\paragraph{The Milgrom-Shannon Framework Applied.}

Consider an agent choosing action $x$ from set $X$ to maximize:
\begin{equation}
\max_{x \in X} f(x, \Psi) = \max_{x \in X} U^{eff}(x; \Psi)
\end{equation}
where $\Psi$ is the context parameter.

\begin{theorem}[Monotone Comparative Statics for Awareness]
\label{thm:mcs-awareness}
If $(X, \leq_X)$ and $([0,1]^8, \leq)$ are lattices and $f(x, \Psi)$ satisfies:
\begin{enumerate}
    \item $f$ is supermodular in $x$ for each fixed $\Psi$
    \item $f$ has increasing differences in $(x, \Psi)$
\end{enumerate}
then the set of optimal actions $X^*(\Psi) = \argmax_{x \in X} f(x, \Psi)$ is 
\emph{increasing} in $\Psi$ (in the strong set order).
\end{theorem}

\textbf{EBF Application}: If awareness $a_d(\Psi)$ is increasing in context 
factors that enhance salience, then optimal behavior shifts monotonically with 
context improvements.

\paragraph{Increasing Differences: The Cross-Partial Condition.}

For smooth functions, increasing differences is equivalent to:
\begin{equation}
\frac{\partial^2 f}{\partial x_i \partial \psi_j} \geq 0 \quad \text{for all } i, j
\end{equation}

In the EBF framework:
\begin{equation}
\frac{\partial^2 U^{eff}}{\partial d_i \partial \psi_j} = 
\frac{\partial a_{d_i}}{\partial \psi_j} + 
\sum_{d' \neq d_i} \gamma_{d_i d'} \cdot \frac{\partial a_{d_i}}{\partial \psi_j} \cdot a_{d'}
\end{equation}

This is positive when awareness is increasing in context ($\partial a / \partial \psi > 0$) 
and complementarities are positive ($\gamma > 0$).

%------------------------------------------------------------------------------
\subsubsection{Context-Dependent Reference Points}

A crucial application of Utility $\times$ Context complementarity is 
\emph{reference point formation}.

\paragraph{Reference Points as Context-Dependent.}

Following \citet{PAP-koszegi2006model}, the reference point $r$ is not fixed but 
depends on context:
\begin{equation}
r = r(\Psi, \mathbb{E}[x])
\end{equation}
where expectations $\mathbb{E}[x]$ are themselves context-dependent.

\begin{axiom}[Context-Dependent Reference]
\label{ax:B-REF}
\textbf{(B-REF)} The reference point function $r: [0,1]^8 \times \mathbb{R}^n \to \mathbb{R}^n$ 
satisfies:
\begin{enumerate}
    \item \textbf{Monotonicity}: $r$ is increasing in expectations
    \item \textbf{Context sensitivity}: $\partial r / \partial \psi_j \neq 0$ for 
    relevant context dimensions
    \item \textbf{Supermodularity}: $r(\Psi \vee \Psi', \cdot) + r(\Psi \wedge \Psi', \cdot) 
    \geq r(\Psi, \cdot) + r(\Psi', \cdot)$
\end{enumerate}
\end{axiom}

\paragraph{Gain-Loss Utility with Context.}

The value function becomes:
\begin{equation}
v(x | r, \Psi) = \begin{cases}
a_+(\Psi) \cdot (x - r)^\alpha & \text{if } x \geq r \\
-a_-(\Psi) \cdot \lambda(\Psi) \cdot (r - x)^\beta & \text{if } x < r
\end{cases}
\end{equation}
where:
\begin{itemize}
    \item $a_+(\Psi), a_-(\Psi)$: Context-dependent gain/loss sensitivity
    \item $\lambda(\Psi)$: Context-dependent loss aversion coefficient
\end{itemize}

\textbf{Empirical prediction}: Loss aversion $\lambda(\Psi)$ increases with:
\begin{itemize}
    \item Resource scarcity ($\psi_7 \uparrow \Rightarrow \lambda \uparrow$)
    \item Emotional arousal ($|\psi_5 - 0.5| \uparrow \Rightarrow \lambda \uparrow$)
    \item Identity threat ($\psi_8 \uparrow \Rightarrow \lambda \uparrow$)
\end{itemize}

%------------------------------------------------------------------------------
\subsubsection{Visceral Factors and Belief Filters}

The Awareness function incorporates two additional mechanisms that create 
Utility $\times$ Context interactions.

\paragraph{Visceral Factors.}

Following \citet{PAP-loewenstein2000emotions}, visceral states (hunger, pain, arousal) 
systematically distort utility weights:
\begin{equation}
a_d^{visceral}(\Psi, v) = a_d(\Psi) \cdot \phi_d(v)
\end{equation}
where $v \in [0,1]$ is visceral intensity and $\phi_d: [0,1] \to [0, \infty)$ is 
a dimension-specific distortion function.

\begin{example}[Hunger and Financial Decisions]
When $v_{hunger}$ is high:
\begin{itemize}
    \item $\phi_F(v) < 1$: Financial awareness \emph{decreases}
    \item $\phi_P(v) > 1$: Physical (food) awareness \emph{increases}
\end{itemize}
This explains why hungry shoppers overspend on food.
\end{example}

\paragraph{Belief Filters.}

Motivated beliefs create systematic distortions in awareness:
\begin{equation}
a_d^{filtered}(\Psi, \theta) = a_d(\Psi) \cdot B_d(\theta)
\end{equation}
where $\theta$ represents identity-relevant beliefs and $B_d(\theta)$ is a 
belief-dependent filter.

\begin{axiom}[Belief Filter Supermodularity]
\label{ax:B-BELIEF}
\textbf{(B-BELIEF)} The belief filter $B: \Theta \times [0,1]^8 \to [0,1]^6$ satisfies:
\begin{equation}
B_d(\theta, \Psi) \text{ is supermodular in } (\theta, \Psi)
\end{equation}
for identity-consistent dimensions $d$.
\end{axiom}

\textbf{Interpretation}: Strong identity beliefs ($\theta$ high) and supportive 
context ($\Psi$ favorable) together produce \emph{more than additive} filtering 
of identity-threatening information.




%------------------------------------------------------------------------------
\subsection{Teil III: Utility × Context Complementarity}
\label{sec:utility-context}

The distinctive contribution of EBF lies not merely in modeling complementarities 
\emph{among} utility dimensions (Teil II), but in how \textbf{context variables} 
systematically modulate these complementarities. This section develops the formal 
foundations for Utility × Context interactions.

%------------------------------------------------------------------------------
\subsubsection{The Context-Modulated Complementarity Structure}

\paragraph{From Static to Dynamic Complementarity.}

Classical supermodularity theory treats the complementarity matrix $\Gamma = [\gamma_{dd'}]$ 
as fixed. EBF introduces a fundamental innovation: complementarity parameters 
are \emph{functions} of context:
\begin{equation}
\gamma_{dd'}(\boldsymbol{\psi}) : \Psi \to \mathbb{R}
\label{eq:context-gamma}
\end{equation}
where $\boldsymbol{\psi} = (\psi_1, \ldots, \psi_8)$ represents the 8 Ψ-Dimensions 
of behavioral context (Temporal Orientation, Risk Attitude, Social Orientation, 
Cognitive Load, Emotional State, Identity Salience, Autonomy, Trust).

\begin{definition}[Context-Dependent Supermodularity]
\label{def:context-supermodular}
A utility function $U: \mathcal{D} \times \Psi \to \mathbb{R}$ exhibits 
\textbf{context-dependent supermodularity} if for each fixed context 
$\boldsymbol{\psi} \in \Psi$, the function $U(\cdot, \boldsymbol{\psi})$ 
is supermodular on $\mathcal{D}$, but the complementarity structure 
$\gamma_{dd'}(\boldsymbol{\psi})$ varies with $\boldsymbol{\psi}$.
\end{definition}

This is a strict generalization: if $\gamma_{dd'}(\boldsymbol{\psi}) = \gamma_{dd'}$ 
(constant), we recover classical supermodularity.

\paragraph{The Modulation Function.}

Following the EBF architecture, context modulates complementarity through:
\begin{equation}
\gamma_{dd'}(\boldsymbol{\psi}) = \bar{\gamma}_{dd'} \cdot 
\prod_{k=1}^{8} m_k(\psi_k; d, d')
\label{eq:modulation}
\end{equation}
where:
\begin{itemize}
    \item $\bar{\gamma}_{dd'}$ is the \emph{baseline} complementarity (context-neutral)
    \item $m_k(\psi_k; d, d') \in (0, \infty)$ is the modulation factor from Ψ-dimension $k$
    \item $m_k = 1$ indicates no modulation; $m_k > 1$ amplifies; $m_k < 1$ attenuates
\end{itemize}

\begin{example}[Temporal Orientation Modulates F-E Complementarity]
Consider the complementarity between Financial (F) and Emotional (E) utility:
\begin{itemize}
    \item High future orientation ($\psi_1 \uparrow$): Planning enables F-E synergies 
    (saving reduces anxiety) $\Rightarrow m_1 > 1$
    \item High present orientation ($\psi_1 \downarrow$): Immediate gratification 
    dominates $\Rightarrow m_1 < 1$
\end{itemize}
\end{example}

%------------------------------------------------------------------------------
\subsubsection{Monotone Comparative Statics with Context}

\paragraph{Context as Parameter.}

The MCS framework from \citet{milgrom1994monotone} extends naturally when context 
enters as an exogenous parameter. Let $\boldsymbol{\psi}$ be ordered component-wise.

\begin{theorem}[Context-Dependent MCS]
\label{thm:context-mcs}
Suppose $U(d, \boldsymbol{\psi})$ is supermodular in $d$ for each $\boldsymbol{\psi}$, 
and satisfies increasing differences in $(d, \boldsymbol{\psi})$. Then optimal choices 
$d^*(\boldsymbol{\psi})$ are increasing in $\boldsymbol{\psi}$ (in the strong set order).
\end{theorem}

\textbf{EBF Interpretation}: As context shifts (e.g., trust increases), not only 
do optimal utility allocations change, but the \emph{pattern} of complementarities 
that drives these allocations also shifts.

\paragraph{Single-Crossing and Context Heterogeneity.}

The single-crossing property from \citet{milgrom1994monotone} has a natural 
context interpretation:

\begin{definition}[Context Single-Crossing]
\label{def:context-scp}
$U$ satisfies \textbf{context single-crossing} if for all $d' > d$ and 
$\boldsymbol{\psi}' > \boldsymbol{\psi}$:
\begin{equation}
U(d', \boldsymbol{\psi}) \geq U(d, \boldsymbol{\psi}) \Rightarrow 
U(d', \boldsymbol{\psi}') \geq U(d, \boldsymbol{\psi}')
\end{equation}
\end{definition}

This formalizes: if a higher utility profile is preferred in one context, it remains 
preferred in ``higher'' contexts---capturing systematic context effects on preferences.

%------------------------------------------------------------------------------
\subsubsection{The Awareness-Complementarity Nexus}

\paragraph{Awareness as Context-Utility Bridge.}

The EBF Awareness function $A(\boldsymbol{\psi})$ serves as the critical bridge 
between context and utility. From Appendix WAT, awareness determines which utility 
dimensions are ``activated'':
\begin{equation}
\tilde{U}(\mathbf{d}, \boldsymbol{\psi}) = \sum_{d \in \mathcal{D}} 
A_d(\boldsymbol{\psi}) \cdot u_d + 
\sum_{d < d'} A_{dd'}(\boldsymbol{\psi}) \cdot \gamma_{dd'} \cdot u_d \cdot u_{d'}
\label{eq:awareness-utility}
\end{equation}
where $A_d(\boldsymbol{\psi}) \in [0,1]$ is the awareness weight for dimension $d$.

\begin{proposition}[Awareness-Modulated Supermodularity]
\label{prop:awareness-super}
The awareness-weighted utility $\tilde{U}$ is supermodular if and only if:
\begin{equation}
A_{dd'}(\boldsymbol{\psi}) \cdot \gamma_{dd'} \geq 0 \quad \forall d \neq d'
\end{equation}
\end{proposition}

This reveals a key insight: even if underlying complementarity is positive 
($\gamma_{dd'} > 0$), \emph{lack of awareness} ($A_{dd'} \approx 0$) effectively 
``turns off'' the complementarity in decision-making.

\paragraph{Rational Inattention and Complementarity.}

Building on \citet{sims2003implications}, awareness allocation itself can be 
modeled as an optimization problem:
\begin{equation}
\max_{\{A_d\}} \mathbb{E}[\tilde{U}] - \lambda \cdot I(\mathbf{A})
\end{equation}
where $I(\mathbf{A})$ is the information cost (mutual information) of the 
awareness allocation and $\lambda$ is the shadow cost of attention.

\begin{axiom}[Awareness-Complementarity Interaction]
\label{ax:B-AWX}
\textbf{(B-AWX)} Context modulates complementarity through two channels:
\begin{enumerate}
    \item \textbf{Direct}: $\gamma_{dd'}(\boldsymbol{\psi})$ varies with context
    \item \textbf{Indirect}: $A_{dd'}(\boldsymbol{\psi})$ determines which 
    complementarities are ``seen''
\end{enumerate}
Both channels exhibit supermodular structure in $(\boldsymbol{\psi}, \mathbf{d})$.
\end{axiom}

%------------------------------------------------------------------------------
\subsubsection{Visceral Factors and State-Dependent Complementarity}

\paragraph{Loewenstein's Visceral Factors.}

\citet{PAP-loewenstein1996out} identifies visceral factors---hunger, pain, arousal, 
emotion---that systematically distort utility weighting. In the EBF framework, 
visceral states $v$ modulate complementarity:
\begin{equation}
\gamma_{dd'}(v) = \bar{\gamma}_{dd'} \cdot \phi_v(d, d')
\end{equation}
where $\phi_v$ captures visceral distortion.

\begin{example}[Hunger and F-P Complementarity]
Under hunger ($v = $ high), the complementarity between Financial (F) and Physical (P) 
utility shifts:
\begin{itemize}
    \item Normally: $\gamma_{FP} > 0$ (money enables health investments)
    \item Under visceral hunger: Immediate food dominates; $\gamma_{FP}(v) \downarrow$
\end{itemize}
\end{example}

\paragraph{Hot-Cold Empathy Gap.}

The ``empathy gap'' from \citet{loewenstein2005hot} manifests as context-dependent 
prediction errors:
\begin{equation}
\hat{\gamma}_{dd'}(\psi_{\text{cold}}) \neq \gamma_{dd'}(\psi_{\text{hot}})
\end{equation}
Individuals in ``cold'' states systematically mis-predict their complementarity 
structure in ``hot'' states.

\begin{axiom}[State-Dependent Complementarity]
\label{ax:B-VIS}
\textbf{(B-VIS)} Visceral and emotional states $v$ enter as additional context 
dimensions that modulate $\gamma_{dd'}$. Supermodularity is preserved within 
each state, but complementarity magnitudes vary across states.
\end{axiom}

%------------------------------------------------------------------------------
\subsubsection{Belief-Dependent Complementarity}

\paragraph{Motivated Beliefs and Utility Structure.}

Following \citet{benabou2016mindful}, beliefs about complementarity can be 
endogenously motivated:
\begin{equation}
\hat{\gamma}_{dd'} = \arg\max_{\gamma} \left[ 
\mathbb{E}[U | \gamma] - c(\gamma - \gamma^*) 
\right]
\end{equation}
where $\gamma^*$ is the true complementarity and $c(\cdot)$ is the cost of 
belief distortion.

\begin{proposition}[Optimistic Complementarity Beliefs]
\label{prop:motivated-gamma}
Under motivated beliefs, individuals overestimate complementarities that 
support desired self-images:
\begin{equation}
\hat{\gamma}_{dd'} > \gamma^*_{dd'} \text{ when } d, d' \text{ are identity-relevant}
\end{equation}
\end{proposition}

This connects to Identity utility (IDN): complementarities involving identity-relevant 
dimensions are subject to motivated distortion.

%------------------------------------------------------------------------------
\subsubsection{Literature Integration}

\begin{table}[h]
\centering
\small
\begin{tabular}{p{3.5cm}p{4cm}p{5cm}}
\toprule
\textbf{Source} & \textbf{Contribution} & \textbf{EBF Application} \\
\midrule
Milgrom \& Shannon (1994) & Single-crossing, MCS & Context as ordered parameter \\
Sims (2003) & Rational inattention & Awareness allocation costs \\
Loewenstein (1996, 2005) & Visceral factors, empathy gap & State-dependent $\gamma$ \\
Bénabou \& Tirole (2016) & Motivated beliefs & Endogenous $\hat{\gamma}$ \\
Bordalo et al. (2012) & Salience theory & Awareness weights $A_d(\psi)$ \\
Kőszegi \& Rabin (2006) & Reference-dependent utility & Reference-modulated $\gamma$ \\
\bottomrule
\end{tabular}
\caption{Literature foundations for Utility × Context complementarity}
\label{tab:utility-context-lit}
\end{table}

\paragraph{Axiom Summary for Teil III.}

\begin{itemize}
    \item \textbf{B-AWX}: Awareness-Complementarity Interaction (dual channels)
    \item \textbf{B-VIS}: Visceral State-Dependent Complementarity
\end{itemize}




%------------------------------------------------------------------------------
\subsection{Teil IV: Context × Context Complementarity}
\label{sec:context-context}

The 8 Ψ-Dimensions of behavioral context are not independent. This section develops 
the formal structure of \textbf{complementarities among context variables themselves}---a 
second-order effect that shapes how context configurations influence behavior.

%------------------------------------------------------------------------------
\subsubsection{The Ψ-Interaction Structure}

\paragraph{Context as Lattice.}

Let $\Psi = \Psi_1 \times \cdots \times \Psi_8$ be the context space, where each 
$\Psi_k$ is a bounded interval representing one Ψ-dimension:
\begin{align*}
\Psi_1 &: \text{Temporal Orientation} & \Psi_5 &: \text{Emotional State} \\
\Psi_2 &: \text{Risk Attitude} & \Psi_6 &: \text{Identity Salience} \\
\Psi_3 &: \text{Social Orientation} & \Psi_7 &: \text{Autonomy} \\
\Psi_4 &: \text{Cognitive Load} & \Psi_8 &: \text{Trust}
\end{align*}

With component-wise ordering, $\Psi$ forms a \textbf{complete lattice}. This enables 
application of the full supermodularity apparatus to context interactions.

\begin{definition}[Ψ-Complementarity Matrix]
\label{def:psi-matrix}
The \textbf{Ψ-complementarity matrix} $\Lambda = [\lambda_{kl}]_{k,l=1}^{8}$ captures 
pairwise interactions among context dimensions:
\begin{equation}
\lambda_{kl} := \frac{\partial^2 V}{\partial \psi_k \partial \psi_l}
\end{equation}
where $V(\boldsymbol{\psi})$ is the ``context value function''---the indirect utility 
achievable given context configuration $\boldsymbol{\psi}$.
\end{definition}

\paragraph{Supermodularity in Context Space.}

\begin{proposition}[Context Supermodularity]
\label{prop:context-super}
$V(\boldsymbol{\psi})$ is supermodular on $\Psi$ if and only if $\lambda_{kl} \geq 0$ 
for all $k \neq l$. Equivalently, Ψ-dimensions are complements in determining 
achievable welfare.
\end{proposition}

\textbf{Interpretation}: Context supermodularity means that improving one Ψ-dimension 
(e.g., increasing trust) raises the marginal value of improving others (e.g., 
increasing autonomy).

%------------------------------------------------------------------------------
\subsubsection{Empirically Relevant Ψ-Interactions}

\paragraph{The Λ-Matrix for EBF.}

Based on behavioral economics literature, we identify the sign structure of key 
Ψ-interactions:

\begin{table}[h]
\centering
\small
\begin{tabular}{l|cccccccc}
\toprule
$\Lambda$ & $\psi_1$ & $\psi_2$ & $\psi_3$ & $\psi_4$ & $\psi_5$ & $\psi_6$ & $\psi_7$ & $\psi_8$ \\
 & Temp & Risk & Soc & Cog & Emo & Iden & Auto & Trust \\
\midrule
Temporal ($\psi_1$) & -- & $+$ & $?$ & $-$ & $-$ & $+$ & $+$ & $+$ \\
Risk ($\psi_2$) & & -- & $+$ & $-$ & $-$ & $?$ & $+$ & $+$ \\
Social ($\psi_3$) & & & -- & $-$ & $+$ & $+$ & $-$ & $+$ \\
Cognitive ($\psi_4$) & & & & -- & $-$ & $-$ & $-$ & $-$ \\
Emotional ($\psi_5$) & & & & & -- & $+$ & $?$ & $+$ \\
Identity ($\psi_6$) & & & & & & -- & $+$ & $+$ \\
Autonomy ($\psi_7$) & & & & & & & -- & $+$ \\
Trust ($\psi_8$) & & & & & & & & -- \\
\bottomrule
\end{tabular}
\caption{Hypothesized sign structure of Ψ-complementarity matrix $\Lambda$}
\label{tab:lambda-matrix}
\end{table}

\paragraph{Key Interaction Patterns.}

\begin{enumerate}
    \item \textbf{Temporal-Risk Complementarity} ($\lambda_{12} > 0$): Future orientation 
    and risk tolerance reinforce each other. Patient individuals can afford to take 
    calculated risks; risk-tolerant individuals benefit more from long-term planning.
    
    \item \textbf{Cognitive Load as Universal Substitute} ($\lambda_{4k} < 0$ for most $k$): 
    High cognitive load undermines the effectiveness of all other context dimensions. 
    This formalizes the ``bandwidth tax'' of \citet{PAP-mullainathan2013scarcity}.
    
    \item \textbf{Trust-Autonomy Complementarity} ($\lambda_{78} > 0$): Trust enables 
    effective exercise of autonomy; autonomy makes trust more valuable. This captures 
    the institutional foundations of economic freedom.
    
    \item \textbf{Social-Autonomy Substitutability} ($\lambda_{37} < 0$): Strong social 
    orientation may conflict with individual autonomy---the tension between collective 
    and individual decision-making.
    
    \item \textbf{Identity-Emotional Complementarity} ($\lambda_{56} > 0$): Emotional 
    engagement amplifies identity concerns; salient identity intensifies emotional 
    responses. This drives ``identity-protective cognition'' \citep{kahan2017cultural}.
\end{enumerate}

%------------------------------------------------------------------------------
\subsubsection{Dynamic Context Complementarity}

\paragraph{Intertemporal Ψ-Interactions.}

Context dimensions exhibit \emph{intertemporal} complementarities---current context 
affects future context through accumulation and habit formation:
\begin{equation}
\psi_k(t+1) = g_k(\psi_k(t), \boldsymbol{\psi}_{-k}(t), a(t))
\end{equation}
where $a(t)$ is the action taken at time $t$.

\begin{definition}[Dynamic Context Complementarity]
\label{def:dynamic-psi}
Ψ-dimensions $k$ and $l$ exhibit \textbf{dynamic complementarity} if:
\begin{equation}
\frac{\partial^2 \psi_k(t+1)}{\partial \psi_l(t) \partial a(t)} > 0
\end{equation}
Current investment in $l$ raises the productivity of actions in building $k$.
\end{definition}

\begin{example}[Trust-Building Dynamics]
Trust ($\psi_8$) accumulates through repeated positive interactions. The marginal 
trust-building effect of cooperation is higher when:
\begin{itemize}
    \item Initial trust is moderate (not too low or too high)
    \item Social orientation ($\psi_3$) is high (cooperation is valued)
    \item Identity alignment ($\psi_6$) is present (in-group interactions)
\end{itemize}
This illustrates $\lambda_{38} > 0$ and $\lambda_{68} > 0$ in the dynamic setting.
\end{example}

\paragraph{Habit Formation and Context Lock-In.}

\citet{PAP-BeckerMurphy1988} model habit formation as intertemporal complementarity 
in consumption. Extending to context:
\begin{equation}
V(\boldsymbol{\psi}(t), \boldsymbol{\psi}(t-1)) \text{ exhibits increasing differences}
\end{equation}
Past context configurations create ``grooves'' that shape current context effectiveness.

\begin{axiom}[Dynamic Context Complementarity]
\label{ax:B-DYN}
\textbf{(B-DYN)} The context transition function $g: \Psi \times \mathcal{A} \to \Psi$ 
preserves supermodularity: if $\boldsymbol{\psi}' \geq \boldsymbol{\psi}$, then 
$g(\boldsymbol{\psi}', a) \geq g(\boldsymbol{\psi}, a)$ in the lattice order, and 
$g$ exhibits increasing differences in $(\boldsymbol{\psi}, a)$.
\end{axiom}

%------------------------------------------------------------------------------
\subsubsection{Context Clustering and Behavioral Types}

\paragraph{Supermodular Context Configurations.}

Given Ψ-complementarities, certain context configurations are ``mutually reinforcing'' 
while others are ``internally conflicted'':

\begin{definition}[Coherent Context Configuration]
\label{def:coherent-context}
A context configuration $\boldsymbol{\psi}^*$ is \textbf{coherent} if it is a 
fixed point of the best-response correspondence in context space:
\begin{equation}
\boldsymbol{\psi}^* \in \arg\max_{\boldsymbol{\psi}} V(\boldsymbol{\psi}) 
\text{ s.t. complementarity constraints}
\end{equation}
\end{definition}

\begin{proposition}[Context Clustering]
\label{prop:context-cluster}
Under Ψ-supermodularity, coherent configurations cluster at lattice extremes. 
The set of stable behavioral types is characterized by:
\begin{equation}
\mathcal{T} = \{\boldsymbol{\psi}^L, \boldsymbol{\psi}^H\} \cup \mathcal{T}_{\text{mixed}}
\end{equation}
where $\boldsymbol{\psi}^L$ and $\boldsymbol{\psi}^H$ are the minimal and maximal 
coherent types.
\end{proposition}

\textbf{EBF Interpretation}: This explains the emergence of distinct ``behavioral 
personas''---internally consistent combinations of context dimensions that form 
stable types (e.g., the ``patient, trusting, autonomous'' type vs. the ``present-focused, 
skeptical, conformist'' type).

\paragraph{The Topology of Behavioral Types.}

The structure of $\Lambda$ determines which behavioral types can coexist:

\begin{itemize}
    \item \textbf{Fully complementary} ($\lambda_{kl} \geq 0$ for all $k,l$): Two extreme 
    types plus continuum of intermediate types
    
    \item \textbf{Block-diagonal}: Clustering within subsets of Ψ-dimensions; types 
    defined by combinations across blocks
    
    \item \textbf{Sparse}: Few significant interactions; near-independence of dimensions 
    allows factorial combination of types
\end{itemize}

%------------------------------------------------------------------------------
\subsubsection{Intervention Design with Context Complementarity}

\paragraph{Targeting Context Bundles.}

Single-dimension interventions may be ineffective if Ψ-complementarities are strong:

\begin{proposition}[Bundle Effectiveness]
\label{prop:bundle}
Let $\Delta \boldsymbol{\psi} = (\Delta \psi_k, \Delta \psi_l, 0, \ldots, 0)$ be a 
two-dimensional context intervention. The welfare gain is:
\begin{equation}
\Delta V \approx \frac{\partial V}{\partial \psi_k} \Delta \psi_k + 
\frac{\partial V}{\partial \psi_l} \Delta \psi_l + 
\lambda_{kl} \Delta \psi_k \Delta \psi_l
\end{equation}
When $\lambda_{kl} > 0$, bundled interventions exhibit \emph{superadditive} returns.
\end{proposition}

\begin{example}[Financial Capability Intervention]
To improve financial behavior, target both:
\begin{itemize}
    \item Temporal orientation ($\psi_1$): Future-self connection exercises
    \item Trust in institutions ($\psi_8$): Information about deposit insurance
\end{itemize}
Since $\lambda_{18} > 0$, the combined intervention exceeds the sum of separate effects.
\end{example}

\paragraph{Sequencing Complementary Interventions.}

Dynamic complementarity implies that intervention \emph{order} matters:

\begin{axiom}[Intervention Sequencing]
\label{ax:B-SEQ}
\textbf{(B-SEQ)} Under dynamic Ψ-complementarity, interventions on dimension $k$ 
should precede interventions on dimension $l$ if:
\begin{equation}
\frac{\partial^2 \psi_l(t+1)}{\partial \psi_k(t) \partial a_l(t)} > 
\frac{\partial^2 \psi_k(t+1)}{\partial \psi_l(t) \partial a_k(t)}
\end{equation}
Building $k$ first creates better conditions for building $l$.
\end{axiom}

%------------------------------------------------------------------------------
\subsubsection{Literature Foundations}

\begin{table}[h]
\centering
\small
\begin{tabular}{p{3.5cm}p{4cm}p{5cm}}
\toprule
\textbf{Source} & \textbf{Contribution} & \textbf{EBF Application} \\
\midrule
Becker \& Murphy (1988) & Rational addiction, habits & Dynamic context complementarity \\
Mullainathan \& Shafir (2013) & Scarcity, bandwidth tax & Cognitive load as substitute \\
Kahan (2017) & Cultural cognition & Identity-emotional interaction \\
Bowles (1998) & Endogenous preferences & Context dynamics \\
Heckman et al. (2006) & Skill complementarity & Ψ-dimension accumulation \\
Acemoglu \& Robinson (2012) & Institutional complementarity & Trust-autonomy nexus \\
\bottomrule
\end{tabular}
\caption{Literature foundations for Context × Context complementarity}
\label{tab:context-context-lit}
\end{table}

\paragraph{Axiom Summary for Teil IV.}

\begin{itemize}
    \item \textbf{B-DYN}: Dynamic Context Complementarity (transition supermodularity)
    \item \textbf{B-SEQ}: Intervention Sequencing (optimal ordering under dynamic Ψ-complementarity)
\end{itemize}




%------------------------------------------------------------------------------
\subsection{Teil V: Mechanism Design under Complementarity}
\label{sec:mechanism-design}

Behavioral interventions---nudges, incentive schemes, information provision---must 
account for complementarity structure. This section develops mechanism design 
principles when agents exhibit supermodular preferences and strategic interactions.

%------------------------------------------------------------------------------
\subsubsection{Incentive Compatibility with Complementarities}

\paragraph{The Revelation Principle under Supermodularity.}

Classical mechanism design assumes additively separable utilities. With complementarities, 
the design space expands but so do the constraints:

\begin{definition}[Supermodular Mechanism]
\label{def:super-mechanism}
A mechanism $\mathcal{M} = (\mathcal{A}, g, t)$ is \textbf{supermodular} if:
\begin{enumerate}
    \item The action space $\mathcal{A}$ forms a lattice
    \item The outcome function $g: \mathcal{A}^n \to X$ preserves lattice structure
    \item Transfers $t_i(a)$ exhibit increasing differences in $(a_i, a_{-i})$
\end{enumerate}
\end{definition}

\begin{theorem}[Kushnir-Liu Extension]
\label{thm:kushnir-liu}
\textup{(\citet{kushnir2019equivalence})} In supermodular environments, Bayesian 
incentive compatibility is equivalent to dominant strategy incentive compatibility 
when utilities satisfy:
\begin{equation}
\frac{\partial^2 u_i}{\partial a_i \partial \theta_i} \geq 0 \quad \text{and} \quad
\frac{\partial^2 u_i}{\partial a_i \partial a_j} \geq 0
\end{equation}
where $\theta_i$ is agent $i$'s type and $a_j$ are others' actions.
\end{theorem}

\textbf{EBF Implication}: When designing behavioral interventions, if the target 
behavior exhibits complementarity with context ($\psi$) and with others' behaviors, 
mechanisms that work under dominant strategies also work under Bayesian equilibrium---a 
significant robustness property.

\paragraph{Monotone Mechanisms.}

\begin{definition}[Monotone Allocation Rule]
\label{def:monotone-allocation}
An allocation rule $q: \Theta \to X$ is \textbf{monotone} if for types 
$\theta' \geq \theta$ (in the lattice order):
\begin{equation}
q(\theta') \geq q(\theta)
\end{equation}
Higher types receive (weakly) higher allocations.
\end{definition}

\begin{proposition}[Implementability under Complementarity]
\label{prop:implement-comp}
Under supermodular preferences, an allocation rule is implementable if and only if 
it is monotone. The payment rule is uniquely determined (up to a constant) by:
\begin{equation}
t(\theta) = t(\underline{\theta}) + \int_{\underline{\theta}}^{\theta} 
\frac{\partial u}{\partial q}(q(s), s) \cdot q'(s) \, ds
\end{equation}
\end{proposition}

%------------------------------------------------------------------------------
\subsubsection{Nudge Design with Complementarity}

\paragraph{The Anatomy of a Nudge.}

In the EBF framework, a nudge operates by shifting context $\boldsymbol{\psi}$:
\begin{equation}
\text{Nudge}: \boldsymbol{\psi} \mapsto \boldsymbol{\psi}' = \boldsymbol{\psi} + \Delta\boldsymbol{\psi}
\end{equation}

The effectiveness depends on how this context shift interacts with existing 
complementarity structure:

\begin{theorem}[Nudge Effectiveness under Complementarity]
\label{thm:nudge-effect}
Let $\Delta V$ be the welfare change from nudge $\Delta\boldsymbol{\psi}$. Then:
\begin{equation}
\Delta V = \underbrace{\nabla_{\psi} V \cdot \Delta\boldsymbol{\psi}}_{\text{Direct effect}} + 
\underbrace{\frac{1}{2} \Delta\boldsymbol{\psi}^T \Lambda \Delta\boldsymbol{\psi}}_{\text{Complementarity effect}} + 
\underbrace{O(\|\Delta\boldsymbol{\psi}\|^3)}_{\text{Higher order}}
\end{equation}
where $\Lambda$ is the Ψ-complementarity matrix from Teil IV.
\end{theorem}

\textbf{Design Principle}: Nudges should target complementary context dimensions 
jointly. A nudge on $\psi_k$ alone captures only the direct effect; bundling with 
complementary $\psi_l$ (where $\lambda_{kl} > 0$) captures additional gains.

\paragraph{Optimal Nudge Portfolios.}

Given budget constraint $C(\Delta\boldsymbol{\psi}) \leq B$, the optimal nudge 
portfolio solves:
\begin{equation}
\max_{\Delta\boldsymbol{\psi}} \left\{ \nabla_{\psi} V \cdot \Delta\boldsymbol{\psi} + 
\frac{1}{2} \Delta\boldsymbol{\psi}^T \Lambda \Delta\boldsymbol{\psi} \right\}
\quad \text{s.t.} \quad C(\Delta\boldsymbol{\psi}) \leq B
\end{equation}

\begin{proposition}[Complementarity Premium]
\label{prop:comp-premium}
The shadow value of the budget constraint is higher when $\Lambda$ has more 
positive entries. Complementary nudge portfolios generate greater ``bang per buck.''
\end{proposition}

%------------------------------------------------------------------------------
\subsubsection{Information Design and Attention Allocation}

\paragraph{Bayesian Persuasion under Complementarity.}

Following \citet{kamenica2011bayesian}, an information designer chooses signal 
structure $\pi$ to influence beliefs and thus behavior. With complementarities:

\begin{definition}[Complementarity-Aware Persuasion]
\label{def:comp-persuasion}
Signal $\pi$ is \textbf{complementarity-exploiting} if it induces beliefs that 
activate complementary dimensions:
\begin{equation}
\mathbb{E}_{\pi}[A_d(\boldsymbol{\psi}) \cdot A_{d'}(\boldsymbol{\psi})] > 
\mathbb{E}_{\pi}[A_d(\boldsymbol{\psi})] \cdot \mathbb{E}_{\pi}[A_{d'}(\boldsymbol{\psi})]
\end{equation}
for pairs $(d, d')$ with $\gamma_{dd'} > 0$.
\end{definition}

\textbf{Interpretation}: Effective persuasion doesn't just shift beliefs about 
individual dimensions---it creates \emph{correlated awareness} across complementary 
dimensions.

\paragraph{Rational Inattention Meets Mechanism Design.}

When agents face attention costs \citep{sims2003implications}, the mechanism 
designer must account for endogenous information acquisition:

\begin{equation}
\max_{a} \mathbb{E}[u(a, \theta) | s] - \lambda \cdot I(s; \theta)
\end{equation}
where $I(s; \theta)$ is mutual information between signal $s$ and state $\theta$.

\begin{axiom}[Attention-Compatible Mechanisms]
\label{ax:B-ATT}
\textbf{(B-ATT)} Optimal mechanisms under rational inattention satisfy:
\begin{enumerate}
    \item \textbf{Simplicity}: Reduce dimensionality of decision-relevant information
    \item \textbf{Salience}: Highlight complementary attributes jointly
    \item \textbf{Chunking}: Bundle complementary dimensions into single ``attention units''
\end{enumerate}
\end{axiom}

%------------------------------------------------------------------------------
\subsubsection{Dynamic Mechanism Design}

\paragraph{Intertemporal Incentive Compatibility.}

With dynamic complementarities (Teil IV), mechanisms must maintain incentive 
compatibility across time:

\begin{definition}[Dynamically Supermodular Mechanism]
\label{def:dynamic-mechanism}
A dynamic mechanism $\{\mathcal{M}_t\}_{t=1}^{T}$ is \textbf{dynamically supermodular} 
if:
\begin{equation}
V_t(a_t, \boldsymbol{\psi}_t; \{a_s\}_{s<t}) \text{ exhibits increasing differences in } 
(a_t, \{a_s\}_{s<t})
\end{equation}
Current incentive compatibility strengthens (not weakens) with past compliance.
\end{definition}

\begin{example}[Commitment Devices]
A savings commitment device is dynamically supermodular if:
\begin{itemize}
    \item Past deposits ($a_{t-1}$) increase marginal value of current deposits ($a_t$)
    \item Context ($\psi_t$: trust, future orientation) amplifies this effect
    \item The mechanism ``locks in'' complementarity gains over time
\end{itemize}
\end{example}

\paragraph{Ratchet Effects and Mechanism Persistence.}

\begin{proposition}[Complementarity Ratchet]
\label{prop:ratchet}
Under dynamic supermodularity, welfare gains from early mechanism adoption are 
\emph{persistent}:
\begin{equation}
\frac{\partial V_T}{\partial a_1} > 0 \quad \text{(early adoption raises terminal welfare)}
\end{equation}
This justifies front-loading behavioral interventions.
\end{proposition}

%------------------------------------------------------------------------------
\subsubsection{Robust Mechanism Design}

\paragraph{Uncertainty about Complementarity Structure.}

In practice, the exact $\Gamma$ and $\Lambda$ matrices are unknown. Robust design 
hedges against this uncertainty:

\begin{definition}[Γ-Robust Mechanism]
\label{def:gamma-robust}
A mechanism is \textbf{Γ-robust} if it achieves performance guarantee:
\begin{equation}
\min_{\Gamma \in \mathcal{G}} V(\mathcal{M}; \Gamma) \geq \underline{V}
\end{equation}
where $\mathcal{G}$ is the uncertainty set for the complementarity matrix.
\end{definition}

\begin{proposition}[Robustness-Efficiency Tradeoff]
\label{prop:robust-tradeoff}
Let $\mathcal{M}^*(\Gamma)$ be the optimal mechanism given known $\Gamma$, and 
$\mathcal{M}^R$ be the Γ-robust mechanism. Then:
\begin{equation}
\mathbb{E}_{\Gamma}[V(\mathcal{M}^R; \Gamma)] \leq \mathbb{E}_{\Gamma}[V(\mathcal{M}^*(\Gamma); \Gamma)]
\end{equation}
with equality only when $|\mathcal{G}| = 1$.
\end{proposition}

\paragraph{Learning Complementarities.}

Adaptive mechanisms can learn $\Gamma$ over time:
\begin{equation}
\hat{\Gamma}_{t+1} = \hat{\Gamma}_t + \alpha_t \cdot (r_t - \hat{r}_t) \cdot \nabla_{\Gamma} \hat{r}_t
\end{equation}
where $r_t$ is observed response and $\hat{r}_t$ is predicted response.

\begin{axiom}[Adaptive Mechanism Design]
\label{ax:B-ADP}
\textbf{(B-ADP)} Mechanisms should incorporate learning about complementarity structure:
\begin{enumerate}
    \item \textbf{Exploration}: Occasionally test non-obvious dimension combinations
    \item \textbf{Exploitation}: Concentrate on estimated high-$\gamma$ pairs
    \item \textbf{Updating}: Bayesian revision of $\hat{\Gamma}$ based on outcomes
\end{enumerate}
\end{axiom}

%------------------------------------------------------------------------------
\subsubsection{Applications to Behavioral Policy}

\paragraph{Default Design.}

Defaults are powerful nudges. Under complementarity:

\begin{proposition}[Optimal Default under Complementarity]
\label{prop:optimal-default}
The optimal default $d^*$ maximizes:
\begin{equation}
d^* = \arg\max_d \left\{ \sum_{i} \pi_i \cdot U_i(d) + (1-\pi_i) \cdot U_i(a_i^*) \right\}
\end{equation}
where $\pi_i$ is the probability agent $i$ sticks with default, and this probability 
itself depends on context: $\pi_i = \pi(\boldsymbol{\psi}_i, d)$.
\end{proposition}

With complementarities in $\boldsymbol{\psi}$, default stickiness varies systematically 
across populations---high cognitive load ($\psi_4$ low) increases $\pi_i$.

\paragraph{Incentive Design.}

\begin{table}[h]
\centering
\small
\begin{tabular}{p{3cm}p{4cm}p{5cm}}
\toprule
\textbf{Intervention} & \textbf{Complementarity Used} & \textbf{Design Implication} \\
\midrule
Savings nudge & $\gamma_{FP} > 0$ (F-P) & Bundle with health framing \\
Exercise app & $\lambda_{15} > 0$ (Temp-Emo) & Add future-self visualization \\
Voting reminder & $\lambda_{36} > 0$ (Soc-Iden) & Emphasize group identity \\
Tax compliance & $\lambda_{38} > 0$ (Soc-Trust) & Social norm + trust messaging \\
Retirement planning & $\lambda_{12} > 0$ (Temp-Risk) & Address both time and risk \\
\bottomrule
\end{tabular}
\caption{Complementarity-informed behavioral intervention design}
\label{tab:intervention-design}
\end{table}

%------------------------------------------------------------------------------
\subsubsection{Literature Integration}

\begin{table}[h]
\centering
\small
\begin{tabular}{p{3.5cm}p{4cm}p{5cm}}
\toprule
\textbf{Source} & \textbf{Contribution} & \textbf{EBF Application} \\
\midrule
Kushnir \& Liu (2019) & BIC $\Leftrightarrow$ DSIC equivalence & Robust mechanism design \\
Kamenica \& Gentzkow (2011) & Bayesian persuasion & Information design \\
Sims (2003) & Rational inattention & Attention-compatible mechanisms \\
Bergemann \& Morris (2019) & Information design survey & Signal structure optimization \\
Pavan et al. (2014) & Dynamic mechanism design & Intertemporal IC \\
Carroll (2015) & Robustness and contracts & Γ-robust mechanisms \\
\bottomrule
\end{tabular}
\caption{Literature foundations for mechanism design under complementarity}
\label{tab:mechanism-lit}
\end{table}

\paragraph{Axiom Summary for Teil V.}

\begin{itemize}
    \item \textbf{B-ATT}: Attention-Compatible Mechanisms (simplicity, salience, chunking)
    \item \textbf{B-ADP}: Adaptive Mechanism Design (explore, exploit, update $\hat{\Gamma}$)
\end{itemize}




%------------------------------------------------------------------------------
\subsection{Teil VI: Large Supermodular Games}
\label{sec:large-games}

Behavioral interventions at scale require understanding how complementarities 
operate in populations. This section develops the theory of supermodular games 
with many players---the foundation for modeling social dynamics under complementarity.

%------------------------------------------------------------------------------
\subsubsection{Supermodular Games: Foundations}

\paragraph{Strategic Complementarity in Games.}

\begin{definition}[Supermodular Game]
\label{def:super-game}
A game $\mathcal{G} = (N, \{A_i\}, \{u_i\})$ is \textbf{supermodular} if for each 
player $i \in N$:
\begin{enumerate}
    \item The action space $A_i$ is a complete lattice
    \item The payoff $u_i(a_i, a_{-i})$ is supermodular in $a_i$ for each $a_{-i}$
    \item $u_i$ exhibits \textbf{increasing differences} in $(a_i, a_{-i})$:
    \begin{equation}
    u_i(a_i', a_{-i}') - u_i(a_i, a_{-i}') \geq u_i(a_i', a_{-i}) - u_i(a_i, a_{-i})
    \end{equation}
    for all $a_i' \geq a_i$ and $a_{-i}' \geq a_{-i}$
\end{enumerate}
\end{definition}

The increasing differences condition captures \emph{strategic complementarity}: 
higher actions by others raise the marginal benefit of one's own higher action.

\begin{theorem}[Topkis-Vives]
\label{thm:topkis-vives}
\textup{(\citet{topkis1998supermodularity}, \citet{vives1990nash})}
In a supermodular game:
\begin{enumerate}
    \item The set of Nash equilibria is nonempty
    \item There exist \textbf{extremal equilibria} $\underline{a}^*$ and $\bar{a}^*$
    \item Best response correspondences are increasing (in the lattice order)
    \item Iterating best responses from any initial point converges monotonically
\end{enumerate}
\end{theorem}

\textbf{EBF Implication}: When behavioral interventions exhibit strategic 
complementarity across individuals, equilibrium existence is guaranteed, and the 
system converges to either a ``high'' or ``low'' equilibrium---explaining why 
populations can ``tip'' between stable behavioral patterns.

%------------------------------------------------------------------------------
\subsubsection{Games with a Continuum of Players}

\paragraph{Mean-Field Supermodular Games.}

Following \citet{balbus2019supermodular}, large populations are modeled with a 
continuum of players $i \in [0,1]$:

\begin{definition}[Mean-Field Supermodular Game]
\label{def:mean-field}
A mean-field game is \textbf{supermodular} if payoffs depend on own action $a_i$ 
and the distribution of others' actions $\mu \in \mathcal{P}(A)$:
\begin{equation}
u_i(a_i, \mu) \text{ exhibits increasing differences in } (a_i, \mu)
\end{equation}
where $\mu' \geq \mu$ in the first-order stochastic dominance order.
\end{definition}

\begin{theorem}[Equilibrium in Mean-Field Supermodular Games]
\label{thm:mean-field-eq}
\textup{(\citet{balbus2019supermodular})}
In mean-field supermodular games with compact action spaces:
\begin{enumerate}
    \item Equilibrium distributions exist
    \item Extremal equilibria $\underline{\mu}^*$ and $\bar{\mu}^*$ exist
    \item The equilibrium set forms a complete lattice under FOSD ordering
\end{enumerate}
\end{theorem}

\paragraph{Population Complementarity.}

Define the \textbf{population complementarity coefficient}:
\begin{equation}
\gamma^{\text{pop}} := \frac{\partial^2 u_i}{\partial a_i \partial \bar{a}}
\end{equation}
where $\bar{a} = \int a \, d\mu(a)$ is the mean action.

\begin{proposition}[Multiplier Effect]
\label{prop:multiplier}
Under mean-field supermodularity, a uniform exogenous shock $\Delta \theta$ to all 
agents' incentives produces equilibrium response:
\begin{equation}
\Delta \bar{a}^* = \frac{1}{1 - \gamma^{\text{pop}}} \cdot \frac{\partial a^*}{\partial \theta} \cdot \Delta\theta
\end{equation}
When $\gamma^{\text{pop}} > 0$, the social multiplier $\frac{1}{1-\gamma^{\text{pop}}} > 1$ 
amplifies individual responses.
\end{proposition}

\textbf{EBF Application}: Behavioral interventions in populations with positive 
$\gamma^{\text{pop}}$ exhibit ``contagion''---individual behavior change spreads 
through complementarity, producing aggregate effects larger than the sum of 
individual effects.

%------------------------------------------------------------------------------
\subsubsection{Multiple Equilibria and Tipping Points}

\paragraph{Equilibrium Multiplicity.}

Supermodular games generically exhibit multiple equilibria:

\begin{proposition}[Equilibrium Structure]
\label{prop:eq-structure}
In symmetric supermodular games with scalar actions, the set of symmetric Nash 
equilibria has odd cardinality. Generically, there are either 1 or 3 symmetric 
equilibria.
\end{proposition}

With three equilibria, the middle one is unstable, leaving two stable equilibria---a 
``high'' and ``low'' behavioral norm.

\paragraph{Tipping Dynamics.}

\begin{definition}[Tipping Point]
\label{def:tipping}
A \textbf{tipping point} occurs at context $\boldsymbol{\psi}^*$ where:
\begin{equation}
\bar{a}^*(\boldsymbol{\psi}^* - \varepsilon) \neq \lim_{\varepsilon \to 0^+} 
\bar{a}^*(\boldsymbol{\psi}^* + \varepsilon)
\end{equation}
The equilibrium jumps discontinuously as context crosses a threshold.
\end{definition}

\begin{theorem}[Tipping under Complementarity]
\label{thm:tipping}
In mean-field supermodular games, tipping points exist if and only if there 
exists $\boldsymbol{\psi}$ with multiple equilibria and the equilibrium 
correspondence is not lower hemicontinuous at $\boldsymbol{\psi}$.
\end{theorem}

\textbf{Policy Implication}: Near tipping points, small interventions can trigger 
large behavioral shifts. Identifying proximity to tipping points is crucial for 
intervention timing.

\paragraph{Hysteresis.}

\begin{proposition}[Behavioral Hysteresis]
\label{prop:hysteresis}
If the population is at high equilibrium $\bar{a}^H$ and context deteriorates to 
$\boldsymbol{\psi}'$, the population may remain at $\bar{a}^H$ even if 
$\boldsymbol{\psi}'$ also supports low equilibrium $\bar{a}^L$. The \textbf{basin of 
attraction} depends on history.
\end{proposition}

This explains why behavioral norms are ``sticky''---complementarity creates 
path-dependence that resists small shocks.

%------------------------------------------------------------------------------
\subsubsection{Heterogeneous Populations}

\paragraph{Type Heterogeneity.}

Following Teil II's strategic heterogeneity framework, let agents have types 
$\theta_i \in \Theta$ that affect both preferences and complementarity:
\begin{equation}
u_i(a_i, \mu; \theta_i) = v(a_i; \theta_i) + \gamma(\theta_i) \cdot a_i \cdot \bar{a}
\end{equation}

\begin{definition}[Type-Dependent Complementarity]
\label{def:type-comp}
The population exhibits \textbf{type-dependent complementarity} if:
\begin{equation}
\gamma(\theta') > \gamma(\theta) \quad \text{for } \theta' > \theta
\end{equation}
Higher types have stronger strategic complementarity with others.
\end{definition}

\begin{proposition}[Sorting by Complementarity]
\label{prop:sorting}
Under type-dependent complementarity, equilibrium features \textbf{assortative 
behavior}: high-$\gamma$ types cluster at extremal actions.
\end{proposition}

\paragraph{Ψ-Heterogeneity in Populations.}

In EBF, heterogeneity arises from variation in context vectors:
\begin{equation}
\boldsymbol{\psi}_i \sim F_{\Psi} \quad \text{(population distribution of context)}
\end{equation}

\begin{proposition}[Context-Mediated Complementarity]
\label{prop:context-mediated}
Population complementarity depends on context distribution:
\begin{equation}
\gamma^{\text{pop}}(F_{\Psi}) = \int \gamma(\boldsymbol{\psi}) \, dF_{\Psi}(\boldsymbol{\psi})
\end{equation}
Interventions that shift $F_{\Psi}$ toward high-$\gamma$ configurations amplify 
social multiplier effects.
\end{proposition}

%------------------------------------------------------------------------------
\subsubsection{Network Games and Local Complementarity}

\paragraph{Games on Networks.}

When interactions are structured by a network $G = (N, E)$:
\begin{equation}
u_i(a_i, a_{-i}) = v(a_i) + \gamma \cdot a_i \cdot \sum_{j \in N_i(G)} w_{ij} a_j
\end{equation}
where $N_i(G)$ is $i$'s neighborhood and $w_{ij}$ are edge weights.

\begin{theorem}[Network Supermodular Games]
\label{thm:network-super}
\textup{(\citet{jackson2015games})}
In network supermodular games:
\begin{enumerate}
    \item Equilibrium actions are increasing in network centrality
    \item The key player (whose removal maximally reduces total activity) has 
    highest \textbf{intercentrality}: $\sum_j (I - \gamma G)^{-1}_{ij}$
    \item Optimal targeting follows eigenvector centrality of complementarity-weighted 
    adjacency matrix
\end{enumerate}
\end{theorem}

\paragraph{Local vs. Global Complementarity.}

\begin{definition}[Complementarity Scope]
\label{def:comp-scope}
\begin{itemize}
    \item \textbf{Local}: $\gamma_{ij} > 0$ only for $(i,j) \in E$ (network neighbors)
    \item \textbf{Global}: $\gamma_{ij} > 0$ for all $(i,j)$ (mean-field)
    \item \textbf{Hybrid}: $\gamma_{ij} = \gamma^L \cdot \mathbf{1}_{(i,j) \in E} + 
    \gamma^G$ (local + global components)
\end{itemize}
\end{definition}

\begin{axiom}[Network Complementarity Structure]
\label{ax:B-NET}
\textbf{(B-NET)} Behavioral complementarity in populations has both local (network) 
and global (mean-field) components. Optimal interventions account for network 
structure when $\gamma^L / \gamma^G$ is large.
\end{axiom}

%------------------------------------------------------------------------------
\subsubsection{Dynamic Population Games}

\paragraph{Evolutionary Dynamics under Complementarity.}

Population behavior evolves according to:
\begin{equation}
\dot{\mu}_t = \Phi(\mu_t, \boldsymbol{\psi}_t)
\end{equation}
where $\Phi$ is the evolutionary operator (e.g., replicator dynamics, best-response 
dynamics).

\begin{theorem}[Monotone Dynamics]
\label{thm:monotone-dynamics}
Under supermodularity, if $\mu_0' \geq \mu_0$ (FOSD), then $\mu_t' \geq \mu_t$ for all 
$t \geq 0$. The dynamics preserve the lattice order.
\end{theorem}

\paragraph{Convergence to Extremal Equilibria.}

\begin{proposition}[Basin Characterization]
\label{prop:basin}
In symmetric supermodular games with best-response dynamics:
\begin{enumerate}
    \item Starting from $\mu_0 = \delta_{\underline{a}}$ (everyone at lowest action), 
    dynamics converge to $\underline{\mu}^*$
    \item Starting from $\mu_0 = \delta_{\bar{a}}$ (everyone at highest action), 
    dynamics converge to $\bar{\mu}^*$
    \item Intermediate starting points may converge to either, depending on 
    complementarity strength
\end{enumerate}
\end{proposition}

\begin{axiom}[Dynamic Complementarity Propagation]
\label{ax:B-PROP}
\textbf{(B-PROP)} In populations with supermodular interactions, behavioral change 
propagates monotonically: early adopters of higher actions pull others upward through 
complementarity, and vice versa for lower actions.
\end{axiom}

%------------------------------------------------------------------------------
\subsubsection{Policy Implications for Population Interventions}

\paragraph{Targeting Strategies.}

\begin{table}[h]
\centering
\small
\begin{tabular}{p{3cm}p{4cm}p{5cm}}
\toprule
\textbf{Strategy} & \textbf{When Optimal} & \textbf{Mechanism} \\
\midrule
Uniform nudge & Low $\gamma^{\text{pop}}$ & Direct effect dominates \\
Target high-$\gamma$ types & Moderate $\gamma^{\text{pop}}$ & Amplify multiplier \\
Target central nodes & Network structure matters & Maximize spillovers \\
Critical mass & Near tipping point & Trigger equilibrium switch \\
Seed high equilibrium & Multiple equilibria & Establish favorable norm \\
\bottomrule
\end{tabular}
\caption{Population intervention strategies under complementarity}
\label{tab:population-strategies}
\end{table}

\paragraph{The Critical Mass Problem.}

\begin{proposition}[Critical Mass Threshold]
\label{prop:critical-mass}
Let $\pi^*$ be the fraction of population that must adopt high action for the 
population to tip to high equilibrium. Then:
\begin{equation}
\pi^* = \frac{\bar{a}^M - \bar{a}^L}{\bar{a}^H - \bar{a}^L}
\end{equation}
where $\bar{a}^M$ is the unstable middle equilibrium.
\end{proposition}

Under strong complementarity ($\gamma^{\text{pop}}$ high), $\pi^*$ is lower---fewer 
early adopters needed to tip the population.

%------------------------------------------------------------------------------
\subsubsection{Literature Integration}

\begin{table}[h]
\centering
\small
\begin{tabular}{p{3.5cm}p{4cm}p{5cm}}
\toprule
\textbf{Source} & \textbf{Contribution} & \textbf{EBF Application} \\
\midrule
Topkis (1998) & Supermodular game theory & Equilibrium existence \\
Vives (1990, 1999) & Strategic complementarity & Monotone comparative statics \\
Balbus et al. (2019) & Continuum player games & Mean-field analysis \\
Jackson (2015) & Network games & Targeting, centrality \\
Morris \& Shin (2002) & Global games, coordination & Tipping points \\
Acemoglu \& Ozdaglar (2011) & Network economics & Propagation dynamics \\
\bottomrule
\end{tabular}
\caption{Literature foundations for large supermodular games}
\label{tab:large-games-lit}
\end{table}

\paragraph{Axiom Summary for Teil VI.}

\begin{itemize}
    \item \textbf{B-NET}: Network Complementarity Structure (local + global)
    \item \textbf{B-PROP}: Dynamic Complementarity Propagation (monotone spread)
\end{itemize}




%------------------------------------------------------------------------------
\subsection{Teil VII: Matching and Assignment under Complementarity}
\label{sec:matching}

Allocation problems---matching workers to firms, customers to products, patients to 
treatments---fundamentally depend on complementarity structure. This final section 
develops matching theory under supermodularity, providing foundations for EBF 
assignment applications.

%------------------------------------------------------------------------------
\subsubsection{Becker's Assortative Matching}

\paragraph{The Classic Framework.}

\citet{PAP-becker1973theory} established that complementarity in production drives 
assortative matching:

\begin{theorem}[Becker's Assortative Matching]
\label{thm:becker}
Let $f(x, y)$ be a production function where $x \in X$ is worker type and $y \in Y$ 
is firm type. If $f$ is \textbf{supermodular}:
\begin{equation}
\frac{\partial^2 f}{\partial x \partial y} > 0
\end{equation}
then the efficient assignment is \textbf{positively assortative}: high-type workers 
match with high-type firms.
\end{theorem}

Conversely, if $f$ is submodular ($\frac{\partial^2 f}{\partial x \partial y} < 0$), 
efficient matching is \emph{negatively assortative}.

\paragraph{EBF Interpretation.}

In behavioral contexts, ``types'' are context configurations $\boldsymbol{\psi}$:
\begin{equation}
f(\boldsymbol{\psi}_{\text{agent}}, \boldsymbol{\psi}_{\text{environment}})
\end{equation}
Supermodularity implies that agents with favorable contexts benefit more from 
favorable environments---a ``Matthew effect'' in behavioral outcomes.

%------------------------------------------------------------------------------
\subsubsection{Matching with Incomplete Information}

\paragraph{Two-Sided Uncertainty.}

Following \citet{johnson2019matching}, consider matching when types are privately known:

\begin{definition}[Bayesian Matching Game]
\label{def:bayesian-matching}
A \textbf{Bayesian matching game} consists of:
\begin{itemize}
    \item Two populations with types $\theta_i \sim F$ and $\eta_j \sim G$
    \item Match value $v(\theta_i, \eta_j)$ supermodular in $(\theta, \eta)$
    \item Signals $s_i = \theta_i + \varepsilon_i$ observed by potential partners
\end{itemize}
\end{definition}

\begin{theorem}[Assortative Matching under Uncertainty]
\label{thm:uncertain-matching}
\textup{(\citet{johnson2019matching})}
In Bayesian matching games with supermodular match values:
\begin{enumerate}
    \item Equilibrium matching is assortative in \textbf{signals} (not true types)
    \item Information revelation improves matching efficiency
    \item The efficiency loss from incomplete information is bounded by signal noise
\end{enumerate}
\end{theorem}

\textbf{EBF Application}: When matching interventions to individuals, incomplete 
knowledge of $\boldsymbol{\psi}$ leads to matching on observable proxies. Better 
context assessment (Appendix WAT) improves allocation efficiency.

%------------------------------------------------------------------------------
\subsubsection{Multi-Dimensional Matching}

\paragraph{Beyond Scalar Types.}

Real matching involves multi-dimensional characteristics. Let types be vectors 
$\boldsymbol{\theta} \in \mathbb{R}^K$:

\begin{definition}[Multi-Attribute Supermodularity]
\label{def:multi-super}
Match value $v(\boldsymbol{\theta}, \boldsymbol{\eta})$ exhibits 
\textbf{multi-attribute supermodularity} if:
\begin{equation}
\frac{\partial^2 v}{\partial \theta_k \partial \eta_l} \geq 0 \quad 
\forall k, l \in \{1, \ldots, K\}
\end{equation}
All cross-partials between agent and partner attributes are nonnegative.
\end{definition}

\begin{proposition}[Dimension-wise Assortativity]
\label{prop:dimension-assort}
Under multi-attribute supermodularity, efficient matching is assortative in 
\textbf{each} dimension separately:
\begin{equation}
\text{Corr}(\theta_k^{\text{agent}}, \eta_k^{\text{partner}}) > 0 \quad \forall k
\end{equation}
\end{proposition}

\paragraph{FEPSDE Matching.}

In EBF, matching individuals to interventions involves 6 utility dimensions 
and 8 context dimensions. The match value:
\begin{equation}
v(\boldsymbol{\psi}_{\text{person}}, \boldsymbol{\phi}_{\text{intervention}}) = 
\sum_{d,d'} \gamma_{dd'}(\boldsymbol{\psi}) \cdot u_d(\boldsymbol{\phi}) \cdot u_{d'}(\boldsymbol{\phi})
\end{equation}
where $\boldsymbol{\phi}$ characterizes the intervention's impact on utility dimensions.

%------------------------------------------------------------------------------
\subsubsection{Stable Matching with Complementarities}

\paragraph{Two-Sided Markets.}

In markets where both sides have preferences (marriage, labor markets), stability 
is the key solution concept:

\begin{definition}[Stable Matching]
\label{def:stable}
A matching $\mu: M \to W$ is \textbf{stable} if:
\begin{enumerate}
    \item No individual prefers being unmatched
    \item No pair $(m, w)$ with $\mu(m) \neq w$ both prefer each other to current matches
\end{enumerate}
\end{definition}

\begin{theorem}[Existence under Complementarity]
\label{thm:stable-exist}
\textup{(\citet{gale1962college})}
Stable matchings exist in two-sided markets. Under supermodular preferences, the 
set of stable matchings forms a \textbf{complete lattice} with man-optimal and 
woman-optimal extremes.
\end{theorem}

\paragraph{Complementarity in Partner Preferences.}

\begin{definition}[Preference Complementarity]
\label{def:pref-comp}
Agent $i$'s preferences exhibit \textbf{complementarity} if:
\begin{equation}
u_i(j | \text{high context}) - u_i(j' | \text{high context}) > 
u_i(j | \text{low context}) - u_i(j' | \text{low context})
\end{equation}
for partners $j > j'$. Better contexts amplify preference for better partners.
\end{definition}

\begin{proposition}[Context-Dependent Assortativity]
\label{prop:context-assort}
Under preference complementarity, positive context shocks increase assortativity 
in equilibrium matching.
\end{proposition}

%------------------------------------------------------------------------------
\subsubsection{Optimal Assignment Mechanisms}

\paragraph{Mechanism Design for Matching.}

When a planner assigns agents to treatments/products:

\begin{definition}[Assignment Mechanism]
\label{def:assignment}
An \textbf{assignment mechanism} $(\mathcal{M}, \mu, t)$ specifies:
\begin{itemize}
    \item Message space $\mathcal{M}$ for agents
    \item Assignment rule $\mu: \mathcal{M}^n \to \Pi$ (set of matchings)
    \item Transfer rule $t: \mathcal{M}^n \to \mathbb{R}^n$
\end{itemize}
\end{definition}

\begin{theorem}[Optimal Assignment under Complementarity]
\label{thm:optimal-assign}
When match values are supermodular and types are private:
\begin{enumerate}
    \item The VCG mechanism implements efficient assignment
    \item Assortative assignment is incentive compatible
    \item Revenue is maximized by serial dictatorship with priority to high types
\end{enumerate}
\end{theorem}

\paragraph{Behavioral Assignment.}

In EBF intervention assignment:
\begin{equation}
\max_{\mu} \sum_{i} v(\boldsymbol{\psi}_i, \boldsymbol{\phi}_{\mu(i)})
\quad \text{s.t. capacity constraints}
\end{equation}

\begin{axiom}[Complementarity-Aware Assignment]
\label{ax:B-ASSIGN}
\textbf{(B-ASSIGN)} Optimal intervention assignment accounts for person-intervention 
complementarity. Individuals with high $\boldsymbol{\psi}$ should receive 
high-impact interventions when $v$ is supermodular; low-impact ``maintenance'' 
interventions when $v$ is submodular.
\end{axiom}

%------------------------------------------------------------------------------
\subsubsection{Dynamic Matching and Learning}

\paragraph{Sequential Matching.}

In practice, matching unfolds over time with learning:

\begin{definition}[Dynamic Matching Problem]
\label{def:dynamic-match}
At each $t$:
\begin{enumerate}
    \item Observe signal $s_t$ about agent types
    \item Update beliefs $\hat{\boldsymbol{\psi}}_t$
    \item Make assignment $\mu_t$
    \item Observe outcome, update for $t+1$
\end{enumerate}
\end{definition}

\begin{proposition}[Value of Information in Matching]
\label{prop:voi-matching}
Under supermodular match values, the value of information about types is:
\begin{equation}
\text{VOI} = \mathbb{E}[v(\theta, \eta^*(\theta))] - \mathbb{E}[v(\theta, \eta^*(\mathbb{E}[\theta]))]
\end{equation}
This is positive and increasing in the degree of supermodularity.
\end{proposition}

\paragraph{Explore-Exploit in Assignment.}

\begin{proposition}[Optimal Learning-Assignment Tradeoff]
\label{prop:explore-exploit}
In repeated assignment with unknown complementarity structure:
\begin{enumerate}
    \item Early periods: Explore diverse assignments to learn $\gamma$
    \item Later periods: Exploit learned structure with assortative assignment
    \item Transition timing depends on complementarity strength and horizon
\end{enumerate}
\end{proposition}

%------------------------------------------------------------------------------
\subsubsection{Applications to EBF}

\paragraph{Intervention-Person Matching.}

\begin{table}[h]
\centering
\small
\begin{tabular}{p{3cm}p{4cm}p{5cm}}
\toprule
\textbf{Context Profile} & \textbf{Optimal Intervention} & \textbf{Complementarity Basis} \\
\midrule
High $\psi_1$ (future) & Long-term commitment devices & Patience × lock-in synergy \\
High $\psi_4$ (cog load) & Simple, automatic nudges & Complexity reduction critical \\
High $\psi_6$ (identity) & Identity-framed messaging & Identity × message resonance \\
High $\psi_8$ (trust) & Information provision & Trust × credibility synergy \\
Low $\psi_2$ (risk averse) & Guaranteed outcomes & Safety × certainty value \\
\bottomrule
\end{tabular}
\caption{Complementarity-based intervention matching}
\label{tab:intervention-matching}
\end{table}

\paragraph{Team Formation.}

When forming teams, complementarity determines optimal composition:

\begin{proposition}[Team Complementarity]
\label{prop:team-comp}
Let team output be $f(\boldsymbol{\psi}_1, \ldots, \boldsymbol{\psi}_n)$. 
\begin{itemize}
    \item \textbf{Supermodular} $f$: Assemble homogeneous high-$\psi$ teams
    \item \textbf{Submodular} $f$: Assemble diverse teams for complementary skills
    \item \textbf{Mixed}: Match complementary dimensions (e.g., high-$\psi_1$ with 
    high-$\psi_2$ if $\lambda_{12} > 0$ in team production)
\end{itemize}
\end{proposition}

\paragraph{Market Segmentation.}

\begin{axiom}[Complementarity-Based Segmentation]
\label{ax:B-SEG}
\textbf{(B-SEG)} Market segments should be defined by complementarity clusters:
\begin{enumerate}
    \item Identify $\boldsymbol{\psi}$-profiles with similar $\gamma$-structures
    \item Design segment-specific intervention bundles
    \item Assign individuals to segments based on estimated $\hat{\boldsymbol{\psi}}$
\end{enumerate}
\end{axiom}

%------------------------------------------------------------------------------
\subsubsection{Literature Integration}

\begin{table}[h]
\centering
\small
\begin{tabular}{p{3.5cm}p{4cm}p{5cm}}
\toprule
\textbf{Source} & \textbf{Contribution} & \textbf{EBF Application} \\
\midrule
Becker (1973) & Assortative matching theory & Person-intervention matching \\
Gale \& Shapley (1962) & Stable matching algorithm & Two-sided market design \\
Johnson (2019) & Matching under uncertainty & Incomplete $\psi$ information \\
Roth (2002) & Market design & Intervention allocation \\
Chiappori et al. (2017) & Multi-dimensional matching & FEPSDE matching \\
Pycia \& Ünver (2017) & Incentive compatible assignment & Mechanism design \\
\bottomrule
\end{tabular}
\caption{Literature foundations for matching under complementarity}
\label{tab:matching-lit}
\end{table}

\paragraph{Axiom Summary for Teil VII.}

\begin{itemize}
    \item \textbf{B-ASSIGN}: Complementarity-Aware Assignment (match high-$\psi$ to 
    high-impact when supermodular)
    \item \textbf{B-SEG}: Complementarity-Based Segmentation (cluster by $\gamma$-structure)
\end{itemize}

%------------------------------------------------------------------------------
\subsubsection{Synthesis: The Complete Complementarity Framework}

This appendix has developed a comprehensive theory of complementarity for EBF:

\begin{enumerate}
    \item \textbf{Teil I}: Mathematical foundations---supermodularity, lattice theory, 
    ordinal equivalence, $\infty$-supermodularity, $r$-decomposability
    
    \item \textbf{Teil II}: Utility × Utility---FEPSDE complementarities, $\Gamma$-matrix, 
    strategic heterogeneity across populations
    
    \item \textbf{Teil III}: Utility × Context---how $\boldsymbol{\psi}$ modulates 
    $\gamma_{dd'}$, awareness-complementarity nexus, motivated beliefs
    
    \item \textbf{Teil IV}: Context × Context---$\Lambda$-matrix of Ψ-interactions, 
    behavioral types, intervention sequencing
    
    \item \textbf{Teil V}: Mechanism Design---nudge portfolios, information design, 
    attention-compatible mechanisms
    
    \item \textbf{Teil VI}: Large Games---social multipliers, tipping points, 
    network complementarity, population dynamics
    
    \item \textbf{Teil VII}: Matching---assortative assignment, intervention targeting, 
    team formation, market segmentation
\end{enumerate}

Together, these provide the theoretical foundation for EBF's approach to 
behavioral intervention: understanding complementarity structure enables more 
effective, efficient, and welfare-enhancing behavioral change.




%------------------------------------------------------------------------------
\subsection{Teil VII: Matching and Assignment under Complementarity}
\label{sec:matching}

Allocation problems---matching workers to firms, consumers to products, interventions 
to populations---depend critically on complementarity structure. This section develops 
matching theory when attributes exhibit supermodular interactions.

%------------------------------------------------------------------------------
\subsubsection{Becker's Assortative Matching}

\paragraph{The Classical Result.}

\citet{PAP-becker1973theory} established the foundational connection between complementarity 
and matching patterns:

\begin{theorem}[Becker's Assortative Matching]
\label{thm:becker}
Consider a two-sided matching market with types $x \in X$ and $y \in Y$. If the 
match surplus function $f(x, y)$ is \textbf{supermodular}:
\begin{equation}
f(x', y') + f(x, y) \geq f(x', y) + f(x, y') \quad \forall x' > x, \, y' > y
\end{equation}
then the efficient matching is \textbf{positively assortative}: high-$x$ matches 
with high-$y$.
\end{theorem}

\textbf{Intuition}: Under complementarity, ``likes match with likes'' because the 
surplus from combining similar types exceeds the surplus from mixing.

\paragraph{Negative Assortative Matching.}

\begin{corollary}[Submodular $\Rightarrow$ Negative Assortment]
\label{cor:negative-assort}
If $f(x, y)$ is \textbf{submodular} (reverse inequality), efficient matching is 
\textbf{negatively assortative}: high-$x$ matches with low-$y$.
\end{corollary}

\textbf{EBF Application}: The sign of complementarity determines whether 
behavioral interventions should target \emph{similar} or \emph{dissimilar} individuals 
for joint treatment.

%------------------------------------------------------------------------------
\subsubsection{Matching with Incomplete Information}

\paragraph{Synchronized Matching.}

\citet{johnson2019synchronized} extends Becker to settings with incomplete information:

\begin{definition}[Synchronized Matching]
\label{def:synchronized}
A matching is \textbf{synchronized} if agents match based on signals $s_x, s_y$ 
that are monotonically related to true types $x, y$:
\begin{equation}
s_x = g_x(x) + \varepsilon_x, \quad s_y = g_y(y) + \varepsilon_y
\end{equation}
with $g_x, g_y$ increasing.
\end{definition}

\begin{theorem}[Assortment under Incomplete Information]
\label{thm:incomplete-assort}
Under supermodular surplus and synchronized signals:
\begin{enumerate}
    \item Equilibrium matching is assortative \emph{in signals}
    \item Expected surplus is lower than complete information benchmark
    \item The \textbf{mismatch cost} is proportional to signal noise and 
    complementarity strength:
    \begin{equation}
    \text{Loss} \propto \frac{\partial^2 f}{\partial x \partial y} \cdot 
    \text{Var}(\varepsilon_x) \cdot \text{Var}(\varepsilon_y)
    \end{equation}
\end{enumerate}
\end{theorem}

\textbf{Implication}: When matching individuals to interventions, measurement error 
in context assessment ($\boldsymbol{\psi}$) is costlier when complementarities are 
stronger---precision matters more under supermodularity.

%------------------------------------------------------------------------------
\subsubsection{Multi-Attribute Matching}

\paragraph{Matching on Multiple Dimensions.}

Real matching problems involve multiple attributes. Let $\mathbf{x} = (x_1, \ldots, x_K)$ 
and $\mathbf{y} = (y_1, \ldots, y_L)$:

\begin{definition}[Multi-Attribute Complementarity]
\label{def:multi-attr}
The surplus function $f(\mathbf{x}, \mathbf{y})$ exhibits:
\begin{itemize}
    \item \textbf{Within-side complementarity}: $\frac{\partial^2 f}{\partial x_k \partial x_l} \geq 0$
    \item \textbf{Cross-side complementarity}: $\frac{\partial^2 f}{\partial x_k \partial y_l} \geq 0$
    \item \textbf{Full supermodularity}: All mixed partials nonnegative
\end{itemize}
\end{definition}

\begin{proposition}[Multi-Dimensional Assortment]
\label{prop:multi-assort}
Under full supermodularity in $(\mathbf{x}, \mathbf{y})$:
\begin{enumerate}
    \item Efficient matching exhibits positive assortment on \emph{all} attribute pairs
    \item The matching can be computed via linear programming on the lattice
    \item Decentralized markets achieve efficiency if transfers are unrestricted
\end{enumerate}
\end{proposition}

\paragraph{EBF: FEPSDE × Ψ Matching.}

Consider matching individuals (characterized by $\boldsymbol{\psi}$) to intervention 
bundles (characterized by target utility dimensions $\mathbf{d}$):

\begin{equation}
f(\boldsymbol{\psi}, \mathbf{d}) = \sum_{k,l} \gamma_{kl}^{\text{match}} \cdot \psi_k \cdot d_l
\end{equation}

\begin{axiom}[Intervention-Context Matching]
\label{ax:B-MATCH}
\textbf{(B-MATCH)} Optimal intervention assignment matches individuals to treatments 
based on complementarity between context profile $\boldsymbol{\psi}$ and intervention 
characteristics $\mathbf{d}$:
\begin{equation}
\text{Assign } i \to j \text{ if } \boldsymbol{\psi}_i \cdot \Gamma^{\text{match}} \cdot \mathbf{d}_j 
\text{ is maximized}
\end{equation}
\end{axiom}

%------------------------------------------------------------------------------
\subsubsection{Stable Matching under Complementarity}

\paragraph{Stability and Supermodularity.}

\begin{definition}[Stable Matching]
\label{def:stable}
A matching $\mu$ is \textbf{stable} if there exists no blocking pair $(x, y)$ 
preferring each other over their current matches.
\end{definition}

\begin{theorem}[Lattice Structure of Stable Matchings]
\label{thm:stable-lattice}
In one-to-one matching markets with supermodular preferences:
\begin{enumerate}
    \item The set of stable matchings forms a \textbf{complete lattice}
    \item There exist $X$-optimal and $Y$-optimal stable matchings
    \item Deferred acceptance algorithms converge to extremal stable matchings
\end{enumerate}
\end{theorem}

\paragraph{Many-to-One Matching.}

For assignment of multiple individuals to treatment groups:

\begin{definition}[Substitutes Condition]
\label{def:substitutes}
Treatment $j$'s preferences over individual sets $S$ satisfy \textbf{substitutes} 
if for $S' \supseteq S$:
\begin{equation}
i \in C_j(S) \text{ and } i \in S' \Rightarrow i \in C_j(S')
\end{equation}
where $C_j(S)$ is $j$'s choice from $S$.
\end{definition}

\begin{theorem}[Kelso-Crawford]
\label{thm:kelso-crawford}
\textup{(\citet{kelso1982job})}
If all treatment preferences satisfy substitutes, stable matchings exist and can 
be found via generalized deferred acceptance.
\end{theorem}

\textbf{Tension}: Supermodularity among individuals (they are complements in treatment 
effectiveness) violates substitutes, potentially destabilizing matching markets.

%------------------------------------------------------------------------------
\subsubsection{Optimal Transport and Matching}

\paragraph{Matching as Optimal Transport.}

Matching problems can be formulated as optimal transport \citep{villani2009optimal}:
\begin{equation}
\min_{\pi \in \Pi(\mu_X, \mu_Y)} \int c(x, y) \, d\pi(x, y)
\end{equation}
where $\Pi(\mu_X, \mu_Y)$ is the set of couplings with marginals $\mu_X, \mu_Y$.

\begin{theorem}[Brenier-McCann]
\label{thm:brenier}
If cost $c(x, y) = -f(x, y)$ (negative surplus) satisfies the \textbf{twist condition} 
(supermodularity implies this), the optimal transport plan is:
\begin{equation}
\pi^* = (\text{id}, T)_{\#} \mu_X
\end{equation}
for some increasing transport map $T: X \to Y$.
\end{theorem}

\textbf{Computational Advantage}: Supermodularity enables efficient algorithms for 
large-scale matching problems (Sinkhorn iterations, auction algorithms).

%------------------------------------------------------------------------------
\subsubsection{Dynamic Matching}

\paragraph{Matching Over Time.}

When individuals arrive sequentially and must be assigned immediately:

\begin{definition}[Online Matching]
\label{def:online}
In \textbf{online matching}, one side (individuals) arrives over time, and 
assignments are irrevocable.
\end{definition}

\begin{proposition}[Competitive Ratio under Complementarity]
\label{prop:competitive}
For online bipartite matching with supermodular surplus:
\begin{enumerate}
    \item Greedy assignment achieves $(1 - 1/e)$-approximation
    \item With stochastic arrivals, competitive ratio improves to $1 - O(1/\sqrt{T})$
    \item Complementarity structure enables better predictions, improving online performance
\end{enumerate}
\end{proposition}

\paragraph{Re-Matching and Adjustment.}

\begin{axiom}[Dynamic Re-Matching]
\label{ax:B-REMATCH}
\textbf{(B-REMATCH)} Under changing context $\boldsymbol{\psi}_t$, optimal intervention 
assignment should be periodically re-evaluated:
\begin{equation}
\mu_t^* = \arg\max_{\mu} \sum_{(i,j) \in \mu} f(\boldsymbol{\psi}_{i,t}, \mathbf{d}_j)
\end{equation}
Re-matching frequency increases with context volatility and complementarity strength.
\end{axiom}

%------------------------------------------------------------------------------
\subsubsection{Applications to Behavioral Intervention Assignment}

\paragraph{Personalized Intervention Matching.}

\begin{table}[h]
\centering
\small
\begin{tabular}{p{3cm}p{4cm}p{5cm}}
\toprule
\textbf{Matching Problem} & \textbf{Complementarity} & \textbf{Assignment Rule} \\
\midrule
Clients to advisors & $\gamma(\psi_{\text{trust}}, \text{expertise})$ & High-trust clients to expert advisors \\
Employees to teams & $\gamma(\psi_{\text{social}}, \text{cohesion})$ & Social individuals to cohesive teams \\
Patients to treatments & $\gamma(\psi_{\text{compliance}}, \text{complexity})$ & Compliant patients to complex regimens \\
Students to programs & $\gamma(\psi_{\text{cognitive}}, \text{rigor})$ & High-capacity to rigorous programs \\
Nudges to segments & $\gamma(\boldsymbol{\psi}, \mathbf{d})$ & Context-intervention alignment \\
\bottomrule
\end{tabular}
\caption{Behavioral intervention matching applications}
\label{tab:matching-applications}
\end{table}

\paragraph{The Personalization-Scalability Tradeoff.}

\begin{proposition}[Matching Granularity]
\label{prop:granularity}
Let $n$ be the number of distinct intervention types and $m$ the number of context 
segments. The welfare gain from finer matching is:
\begin{equation}
\Delta W(n, m) \propto \bar{\gamma} \cdot \sqrt{n \cdot m}
\end{equation}
where $\bar{\gamma}$ is average complementarity. Higher complementarity justifies 
more granular personalization.
\end{proposition}

%------------------------------------------------------------------------------
\subsubsection{Literature Integration}

\begin{table}[h]
\centering
\small
\begin{tabular}{p{3.5cm}p{4cm}p{5cm}}
\toprule
\textbf{Source} & \textbf{Contribution} & \textbf{EBF Application} \\
\midrule
Becker (1973) & Assortative matching & Positive/negative assortment \\
Gale \& Shapley (1962) & Stable matching & Intervention stability \\
Kelso \& Crawford (1982) & Substitutes condition & Many-to-one assignment \\
Johnson (2019) & Incomplete information & Context measurement error \\
Villani (2009) & Optimal transport & Large-scale computation \\
Chiappori et al. (2017) & Multi-dimensional matching & FEPSDE × Ψ alignment \\
\bottomrule
\end{tabular}
\caption{Literature foundations for matching under complementarity}
\label{tab:matching-lit}
\end{table}

\paragraph{Axiom Summary for Teil VII.}

\begin{itemize}
    \item \textbf{B-MATCH}: Intervention-Context Matching (optimize $\boldsymbol{\psi} \cdot \Gamma \cdot \mathbf{d}$)
    \item \textbf{B-REMATCH}: Dynamic Re-Matching (adjust to context evolution)
\end{itemize}

%------------------------------------------------------------------------------
\subsubsection{Synthesis: The Complete Complementarity Framework}

This appendix has developed a comprehensive theory of complementarity at multiple 
levels:

\begin{center}
\begin{tabular}{lll}
\toprule
\textbf{Teil} & \textbf{Domain} & \textbf{Key Structure} \\
\midrule
I & Foundations & Ordinal $\Leftrightarrow$ Cardinal equivalence \\
II & Utility × Utility & $\Gamma = [\gamma_{dd'}]$ matrix \\
III & Utility × Context & $\gamma_{dd'}(\boldsymbol{\psi})$ modulation \\
IV & Context × Context & $\Lambda = [\lambda_{kl}]$ Ψ-interactions \\
V & Mechanism Design & Nudge portfolios, information design \\
VI & Large Games & Social multiplier, tipping points \\
VII & Matching & Assortative assignment \\
\bottomrule
\end{tabular}
\end{center}

Together, these provide the mathematical foundation for EBF's treatment of 
behavioral complementarity across all analytical levels.


%==============================================================================
% BIBLIOGRAPHY
%==============================================================================


%------------------------------------------------------------------------------
\subsection{Teil VIII: Organizational Implications of Complementarity}
\label{sec:organizational}

The preceding sections developed complementarity theory for individual decisions 
and populations. This section extends the framework to \textbf{organizational design}---how 
complementarity structure determines firm boundaries, governance modes, and ecosystem 
architecture. We draw heavily on \citet{baldwin2024design}.

%------------------------------------------------------------------------------
\subsubsection{Strong vs. Supermodular Complementarity: Governance Implications}

\paragraph{The Fundamental Distinction.}

Not all complementarities have the same organizational consequences. Baldwin 
distinguishes two forms:

\begin{definition}[Strong Complementarity]
\label{def:strong-comp}
Components $X$ and $Y$ exhibit \textbf{strong complementarity} if neither has 
value without the other:
\begin{equation}
V(X, 0) = V(0, Y) = 0, \quad V(X, Y) > 0
\label{eq:strong-comp}
\end{equation}
Using binary indicators $x, y \in \{0, 1\}$:
\begin{equation}
V(x, y) = V_0 \cdot x \cdot y
\end{equation}
This multiplicative structure implies complete value destruction if either 
component is missing.
\end{definition}

\begin{definition}[Supermodular Complementarity]
\label{def:supermod-comp-org}
Components exhibit \textbf{supermodular complementarity} if increasing one raises 
the marginal return to increasing the other (``increasing differences''):
\begin{equation}
V(x + \Delta x, y + \Delta y) - V(x, y + \Delta y) \geq V(x + \Delta x, y) - V(x, y)
\label{eq:increasing-diff}
\end{equation}
Equivalently, for smooth functions: $\frac{\partial^2 V}{\partial x \partial y} \geq 0$.
\end{definition}

\begin{proposition}[Inclusion Relationship]
\label{prop:strong-subset}
Strong complementarity $\Rightarrow$ Supermodular complementarity, but not conversely.
\end{proposition}

\begin{proof}
For strong complements with $V(x,y) = V_0 \cdot x \cdot y$, setting $x=y=0$ and 
$\Delta x = \Delta y = 1$:
\begin{align*}
\text{LHS} &= V(1,1) - V(0,1) = V_0 - 0 = V_0 \\
\text{RHS} &= V(1,0) - V(0,0) = 0 - 0 = 0
\end{align*}
Since $V_0 > 0 = 0$, the supermodularity condition holds.
\end{proof}

\paragraph{The Critical Organizational Implication.}

\begin{theorem}[Governance Mode Selection]
\label{thm:governance}
\textup{(Baldwin 2024)}
\begin{enumerate}
    \item \textbf{Strong complementarity} $\Rightarrow$ \textbf{Unified Governance} (integration)
    \item \textbf{Supermodular complementarity} (without strong) $\Rightarrow$ 
    \textbf{Distributed Governance} possible (ecosystems, platforms)
\end{enumerate}
\end{theorem}

The economic logic: Strong complementarity creates \textbf{holdup problems}. If 
$V(X, 0) = 0$, the owner of $Y$ can threaten to ``set $y = 0$,'' destroying all 
value created by $X$. This bargaining hazard necessitates unified control.

%------------------------------------------------------------------------------
\subsubsection{Distributed Supermodular Complementarity (DSMC)}

\paragraph{The Three Conditions for Distributed Governance.}

\citet{baldwin2024design} identifies three mathematical conditions that enable 
independent actors to create complementary value without integration:

\begin{axiom}[DSMC Conditions]
\label{ax:B-DSMC}
\textbf{(B-DSMC)} Distributed Supermodular Complementarity requires:

\textbf{Condition 1: Value Separability}
\begin{equation}
V(x, y) = V_X(x; y) + V_Y(y; x)
\label{eq:value-sep}
\end{equation}
Each actor can compute their own value function, treating the other's choice 
as a parameter (semicolon notation).

\textbf{Condition 2: Integration Costs}
\begin{equation}
V^{FB}(x^*, y^*) - C_{int} < V_X(x^{ind}; y^{ind}) + V_Y(y^{ind}; x^{ind})
\label{eq:integration-cost}
\end{equation}
The efficiency loss from distributed decisions is smaller than integration costs 
(bureaucracy, lost incentives, merger transaction costs).

\textbf{Condition 3: Incentive Alignment}
\begin{equation}
V_X(x + \Delta x; y) - C_X(\Delta x) \geq 0 \quad \text{(and similarly for } Y\text{)}
\label{eq:incentive-align}
\end{equation}
Private returns exceed private costs for each actor independently.
\end{axiom}

\paragraph{Simultaneous vs. Sequential Investment.}

The timing of decisions critically affects whether DSMC can be sustained:

\begin{proposition}[Sequential Investment Advantage]
\label{prop:sequential}
Under supermodularity, sequential investment relaxes incentive constraints:
\begin{enumerate}
    \item If $X$ invests first: $V_X(x + \Delta x; y_0) - C_X \geq 0$
    \item $Y$ observes $x + \Delta x$, then faces: $V_Y(y + \Delta y; x + \Delta x) - C_Y \geq 0$
\end{enumerate}
By increasing differences, $Y$'s hurdle is \emph{lower} than in simultaneous case:
\begin{equation}
V_Y(y + \Delta y; x + \Delta x) > V_Y(y + \Delta y; x_0)
\end{equation}
\end{proposition}

This explains \textbf{ecosystem dynamics}: A platform provider invests first (creating 
infrastructure), then complementors observe and invest (creating apps/services), 
each ``climbing the supermodular value staircase.''

%------------------------------------------------------------------------------
\subsubsection{Network Effects as Supermodularity}

\paragraph{Direct Network Effects.}

Let $n$ be network size. The number of potential connections scales as:
\begin{equation}
\text{Connections}(n) = \frac{n(n-1)}{2} \approx \frac{n^2}{2}
\end{equation}

The marginal value of the $n$-th user:
\begin{equation}
\Delta V(n) = V(n) - V(n-1) \propto n - 1
\end{equation}

Since $\Delta V(n)$ increases in $n$, network effects satisfy increasing differences:
\begin{equation}
\frac{\partial^2 V}{\partial n^2} > 0 \quad \Rightarrow \quad \text{Supermodular in } n
\end{equation}

\paragraph{Indirect Network Effects (Two-Sided Markets).}

For platforms with developers $D$ and users $U$:
\begin{align}
V_U(u; d) &= f_U(u) + g_U(d) \quad \text{with } g_U'(d) > 0 \\
V_D(d; u) &= f_D(d) + g_D(u) \quad \text{with } g_D'(u) > 0
\end{align}

Cross-derivatives are positive:
\begin{equation}
\frac{\partial V_U}{\partial d} > 0, \quad \frac{\partial V_D}{\partial u} > 0
\end{equation}

This creates \textbf{feedback loops}: More users attract more developers, which 
attracts more users---a supermodular cascade.

%------------------------------------------------------------------------------
\subsubsection{Modularity and Thin Crossing Points}

\paragraph{The Role of Modularity.}

Modularity---decomposing systems into loosely coupled components with well-defined 
interfaces---is the architectural prerequisite for DSMC.

\begin{definition}[Thin Crossing Point]
\label{def:thin-crossing}
A \textbf{thin crossing point} is an interface between modules where:
\begin{enumerate}
    \item Information exchange is standardized and limited
    \item Transaction costs are low
    \item Contractual solutions (including side payments) are feasible
\end{enumerate}
\end{definition}

\begin{proposition}[Modularity Enables DSMC]
\label{prop:modularity-dsmc}
At thin crossing points, the Coase theorem applies: If $X$'s investment creates 
positive externality for $Y$ but $V_X < C_X$, side payment $T$ from $Y$ to $X$ 
can restore efficiency:
\begin{equation}
V_X(x + \Delta x; y) + T - C_X \geq 0 \quad \text{where} \quad T \leq V_Y(y; x + \Delta x) - V_Y(y; x)
\end{equation}
Both parties gain if $V_X + V_Y > C_X + C_Y$.
\end{proposition}

Without modularity (thick interfaces, high transaction costs), side payments 
become infeasible, and integration is required.

%------------------------------------------------------------------------------
\subsubsection{EBF Organizational Applications}

\paragraph{The Γ-Matrix Determines Governance.}

The complementarity structure $\Gamma = [\gamma_{dd'}]$ from Teil II has direct 
organizational implications:

\begin{table}[h]
\centering
\small
\begin{tabular}{p{3.5cm}p{4cm}p{5cm}}
\toprule
\textbf{Γ Structure} & \textbf{Complementarity Type} & \textbf{Governance Implication} \\
\midrule
$V(d, 0) = 0$ for some $d$ & Strong & Unified (single firm controls all) \\
$\gamma_{dd'} > 0$, separable & Supermodular, DSMC & Distributed (ecosystem feasible) \\
$\gamma_{dd'} \approx 0$ & Near-independent & Modular (market transactions) \\
$\gamma_{dd'} < 0$ & Substitutes & Competition, not collaboration \\
\bottomrule
\end{tabular}
\caption{Mapping complementarity structure to governance modes}
\label{tab:gamma-governance}
\end{table}

\paragraph{Intervention Design at the Organizational Level.}

\begin{enumerate}
    \item \textbf{Identify complementarity type}: Is the behavioral intervention 
    ``strong'' (worthless without full adoption) or ``supermodular'' (valuable 
    even with partial adoption)?
    
    \item \textbf{Check DSMC conditions}: Can different actors (firms, agencies, 
    individuals) independently create value? Are integration costs high?
    
    \item \textbf{Design accordingly}:
    \begin{itemize}
        \item Strong $\Rightarrow$ Centralized program, single implementing agency
        \item DSMC $\Rightarrow$ Platform model, enable independent complementors
        \item Sequential DSMC $\Rightarrow$ Anchor investor strategy, then followers
    \end{itemize}
\end{enumerate}

\begin{example}[Health Behavior Ecosystem]
Consider a wellness intervention targeting Financial (F), Physical (P), and 
Emotional (E) utility:
\begin{itemize}
    \item If $\gamma_{FP}$, $\gamma_{FE}$, $\gamma_{PE}$ are supermodular but 
    not strong: Enable distributed providers (fitness apps, financial advisors, 
    therapists) with common platform/data standards
    \item If $\gamma$ involves strong complementarity (all-or-nothing treatment): 
    Integrated care model with single provider
\end{itemize}
\end{example}

%------------------------------------------------------------------------------
\subsubsection{Externalities and the Coase Theorem}

\paragraph{When Side Payments Fail.}

Even with DSMC, market failure can occur if:
\begin{equation}
V_X(x + \Delta x; y) < C_X \quad \text{but} \quad V_X + V_Y > C_X
\end{equation}
The investment creates value, but the returns accrue to $Y$ (externality).

\begin{proposition}[Coasean Resolution]
\label{prop:coase}
If transaction costs at crossing points are sufficiently low, $Y$ can pay $X$ 
to invest:
\begin{equation}
T^* \in [C_X - V_X, \, V_Y^{\text{incremental}}]
\end{equation}
This interval is non-empty when total value exceeds total cost.
\end{proposition}

\textbf{EBF Implication}: Behavioral interventions often create externalities 
(individual health behavior benefits employers, insurers, society). Designing 
mechanisms for ``side payments'' (subsidies, tax incentives, insurance premium 
reductions) can align private and social incentives under DSMC.

%------------------------------------------------------------------------------
\subsubsection{The Governance Decision Tree}

Baldwin's framework implies a systematic decision process:

\begin{enumerate}
    \item \textbf{Test for Strong Complementarity}: Does $V(X, 0) = 0$?
    \begin{itemize}
        \item Yes $\Rightarrow$ Holdup risk $\Rightarrow$ \textbf{Unified Governance}
        \item No $\Rightarrow$ Proceed to DSMC test
    \end{itemize}
    
    \item \textbf{Test DSMC Conditions}:
    \begin{itemize}
        \item Separability? (Given for modular designs)
        \item Integration costs high? (Given for specialized knowledge/incentives)
        \item Incentive alignment? (Check private returns vs. costs)
    \end{itemize}
    
    \item \textbf{If DSMC holds}:
    \begin{itemize}
        \item Simultaneous: Independent firms invest in parallel
        \item Sequential: Anchor investor, then ecosystem growth
    \end{itemize}
    
    \item \textbf{If DSMC fails}: Consider side payments or integration
\end{enumerate}

\begin{axiom}[Governance Mode Selection]
\label{ax:B-GOV}
\textbf{(B-GOV)} The optimal governance mode for behavioral interventions depends 
on complementarity structure:
\begin{equation}
\text{Governance}^* = \begin{cases}
\text{Unified} & \text{if strong complementarity (holdup risk)} \\
\text{Platform/Ecosystem} & \text{if DSMC conditions hold} \\
\text{Market + Contracts} & \text{if separable with low transaction costs} \\
\text{Hybrid} & \text{otherwise}
\end{cases}
\end{equation}
\end{axiom}

%------------------------------------------------------------------------------
\subsubsection{Literature Integration}

\begin{table}[h]
\centering
\small
\begin{tabular}{p{3.5cm}p{4cm}p{5cm}}
\toprule
\textbf{Source} & \textbf{Contribution} & \textbf{EBF Application} \\
\midrule
Baldwin (2024) & DSMC, thin crossing points & Ecosystem design \\
Williamson (1985) & Transaction cost economics & Integration vs. market \\
Hart \& Moore (1990) & Property rights, holdup & Strong complementarity \\
Milgrom \& Roberts (1990) & Complementarities in organizations & Supermodular firms \\
Rochet \& Tirole (2003) & Two-sided markets & Platform economics \\
Coase (1960) & Side payments, externalities & Contractual solutions \\
\bottomrule
\end{tabular}
\caption{Literature foundations for organizational complementarity}
\label{tab:org-lit}
\end{table}

\paragraph{Axiom Summary for Teil VIII.}

\begin{itemize}
    \item \textbf{B-DSMC}: Distributed Supermodular Complementarity (three conditions)
    \item \textbf{B-GOV}: Governance Mode Selection (strong vs. supermodular)
\end{itemize}




%------------------------------------------------------------------------------
\subsubsection{Empirical Methods for Complementarity Detection}
\label{sec:empirical-methods}

The theoretical elegance of supermodularity faces significant empirical challenges. 
Following \citet{brynjolfsson2013complementarities}, we review methods for detecting 
and measuring complementarities in organizational and behavioral data.

\paragraph{The Fundamental Measurement Problem.}

Two approaches exist for testing complementarity, each with distinct biases:

\begin{definition}[Performance Test]
\label{def:perf-test}
Estimate the production/utility function directly:
\begin{equation}
Y = \alpha + \beta_1 X_1 + \beta_2 X_2 + \gamma_{12} X_1 X_2 + \varepsilon
\label{eq:perf-test}
\end{equation}
Complementarity is supported if $\hat{\gamma}_{12} > 0$ (interaction term positive).
\end{definition}

\begin{definition}[Correlation Test]
\label{def:corr-test}
Test whether practices cluster in adoption:
\begin{equation}
\text{Corr}(X_1, X_2) > 0 \quad \text{conditional on controls}
\label{eq:corr-test}
\end{equation}
Under rational choice, complements should be adopted together.
\end{definition}

\begin{proposition}[Opposite Biases]
\label{prop:opposite-bias}
\textup{(\citet{brynjolfsson2013complementarities}, Section 3)}
\begin{enumerate}
    \item The \textbf{performance test} is biased \emph{downward}: Firms adopting 
    only one practice self-select because they have unobserved advantages, making 
    isolated adoption appear more profitable than it truly is.
    \item The \textbf{correlation test} is biased \emph{upward}: ``Good managers'' 
    adopt both practices regardless of complementarity, creating spurious correlation.
\end{enumerate}
\end{proposition}

\paragraph{Solutions to the Identification Problem.}

\begin{enumerate}
    \item \textbf{Homogeneous populations}: Study narrow industries where external 
    confounds are minimized (``insider econometrics'')
    
    \item \textbf{Panel data with fixed effects}: Remove time-constant firm-specific 
    factors (e.g., management quality)
    
    \item \textbf{Natural experiments}: Use exogenous shocks (regulatory mandates, 
    technology discontinuities) that affect one practice
    
    \item \textbf{Randomized controlled trials}: The gold standard---randomly assign 
    practices to treatment groups
\end{enumerate}

\begin{example}[Ichniowski et al. 1997]
Studying 36 homogeneous steel finishing lines, researchers found that bundled 
HR practices (teams, incentives, flexible jobs) increased productivity 6.7\%, 
while individual practices had no significant effect---direct evidence for 
$\gamma_{12} > 0$ when properly controlling for industry homogeneity.
\end{example}

%------------------------------------------------------------------------------
\subsubsection{The Ten Theorems of Organizational Complementarity}

\citet{brynjolfsson2013complementarities} derive ten theorems from supermodularity 
that hold without assumptions of convexity or differentiability---crucial for 
discrete organizational decisions.

\begin{theorem}[Pairwise Testing]
\label{thm:pairwise}
\textbf{(Theorem 1)} Complementarity is a pairwise property. For $n$ practices, 
if every pair $(x_i, x_j)$ satisfies supermodularity, the entire system is 
supermodular. No hidden three-way interactions can destroy system complementarity.
\end{theorem}

\textbf{Implication}: To verify EBF's $\Gamma$-matrix, test each $\gamma_{dd'}$ 
pairwise---no need for complex multi-way interaction tests.

\begin{theorem}[Clustering of Optima]
\label{thm:clustering}
\textbf{(Theorem 2)} The set of optimal solutions forms a sublattice. Optimal 
practices cluster: if high-performance organizations use practice $A$, they 
likely also use complementary practice $B$.
\end{theorem}

\begin{theorem}[Coherent Search]
\label{thm:coherent-search}
\textbf{(Theorem 3)} Starting from any point, more than half the potential gains 
can be realized by considering only changes that move \emph{all} variables in 
the same direction (all increase or all decrease).
\end{theorem}

\textbf{Strategic Implication}: On the ``rugged landscape'' of organizational 
possibilities, coherent search dramatically simplifies optimization.

\begin{theorem}[Monotone Comparative Statics]
\label{thm:mcs-org}
\textbf{(Theorem 4)} If an exogenous parameter $\theta$ (e.g., falling IT prices) 
increases the return to one activity, optimal choices of \emph{all} complementary 
activities increase. The optimal solution $x^*(\theta)$ is isotone (non-decreasing) 
in $\theta$.
\end{theorem}

\textbf{EBF Application}: A context shift in one $\psi_k$ dimension propagates 
to all complementary dimensions and utility investments.

\begin{theorem}[Momentum and Inertia]
\label{thm:momentum}
\textbf{(Theorem 8)} In dynamic systems, change exhibits momentum. Once an 
organization starts moving in one direction, marginal returns to further steps 
increase. Conversely, without initial impulse, inertia dominates.
\end{theorem}

\begin{theorem}[Le Chatelier Principle]
\label{thm:lechatelier}
\textbf{(Theorem 7)} Long-run responses exceed short-run responses. When some 
variables are fixed short-term, adjustment is constrained. Long-term, complementarity 
amplifies the response of initially flexible variables.
\end{theorem}

%------------------------------------------------------------------------------
\subsubsection{The Local Optima Problem}

\paragraph{Why Organizational Change Fails.}

Firms often occupy local optima---internally coherent but globally suboptimal 
configurations.

\begin{proposition}[Worse-Before-Better]
\label{prop:worse-before}
Changing only one practice in a complementary system destroys the coherence of 
the old system without achieving the benefits of the new. Performance initially 
\emph{declines} before potentially improving.
\end{proposition}

This explains why incremental change often fails: managers observe declining 
performance and abandon the transition before complementary adjustments are made.

\begin{example}[Johnson \& Johnson]
A factory purchased flexible manufacturing equipment but retained old production 
rules (long runs to minimize setup time). The expensive technology delivered no 
benefit because it was designed for short, flexible runs. Only a ``greenfield'' 
restart with new workers and new mental models unlocked the complementary system.
\end{example}

\paragraph{The Coordination Challenge.}

\begin{proposition}[Big Bang vs. Incremental Change]
\label{prop:big-bang}
Under strong complementarity, simultaneous change of all practices (``Big Bang'') 
dominates incremental change. But Big Bang requires:
\begin{itemize}
    \item Central coordination across departments
    \item Strong leadership as focal point
    \item Shared culture/vision to align expectations
\end{itemize}
Without these, organizations remain trapped in local optima.
\end{proposition}

%------------------------------------------------------------------------------
\subsubsection{The Matrix of Change}

\citet{brynjolfsson2013complementarities} propose a practical tool for managing 
transitions in complementary systems.

\begin{definition}[Matrix of Change]
\label{def:matrix-change}
A visualization tool mapping:
\begin{itemize}
    \item Rows: Current practices
    \item Columns: Target practices
    \item Cell entries: $+$ (reinforcing), $-$ (conflicting), $0$ (neutral)
\end{itemize}
\end{definition}

\begin{table}[h]
\centering
\small
\begin{tabular}{l|ccc}
\toprule
& \textbf{New Practice A} & \textbf{New Practice B} & \textbf{New Practice C} \\
\midrule
\textbf{Old Practice X} & $-$ & $-$ & $0$ \\
\textbf{Old Practice Y} & $-$ & $+$ & $-$ \\
\textbf{Old Practice Z} & $0$ & $-$ & $+$ \\
\bottomrule
\end{tabular}
\caption{Matrix of Change: Identifying transition conflicts}
\label{tab:matrix-change}
\end{table}

\paragraph{Diagnostic Use.}

\begin{enumerate}
    \item \textbf{Diagonal conflicts} (Old$_i$ vs New$_i$): Direct substitution required
    \item \textbf{Many $-$ entries}: Incremental change will be painful; consider 
    Big Bang or spin-off
    \item \textbf{Few $-$ entries}: Gradual transition feasible
\end{enumerate}

\begin{example}[Merrill Lynch]
Transitioning from traditional to online brokerage, nearly all old practices 
(high fees, personal advice, branch network) conflicted with new practices 
(low fees, automation, digital-first). The Matrix revealed that incremental 
change was infeasible; a separate online unit was required.
\end{example}

%------------------------------------------------------------------------------
\subsubsection{Implications for EBF Interventions}

\paragraph{The Cherry-Picking Trap.}

\begin{axiom}[Anti-Cherry-Picking]
\label{ax:B-CHERRY}
\textbf{(B-CHERRY)} Behavioral interventions should not be adopted in isolation 
from complementary practices. A nudge that works in Organization A may fail in 
Organization B if complementary elements (culture, incentives, information systems) 
are absent.
\end{axiom}

\paragraph{Intervention Design Process.}

\begin{enumerate}
    \item \textbf{Map the Γ-matrix}: Identify which behavioral dimensions are 
    complementary in the target population
    
    \item \textbf{Test pairwise}: Use Theorem 1---if all pairs show positive 
    interaction, the system is complementary
    
    \item \textbf{Plan coherent bundles}: Use Theorem 3---design interventions 
    that move all complementary dimensions together
    
    \item \textbf{Anticipate worse-before-better}: Prepare stakeholders for 
    initial performance dips during transition
    
    \item \textbf{Build the Matrix of Change}: Identify conflicts between current 
    and target practices before implementation
\end{enumerate}

\paragraph{M\&A and Cultural Integration.}

\begin{proposition}[Fragility of Organizational Capital]
\label{prop:fragility}
When two organizations with different (internally coherent) complementary systems 
merge, mixing practices can destroy both systems. This explains why many mergers 
destroy rather than create value.
\end{proposition}

\textbf{EBF Implication}: When integrating behavioral programs across 
organizations, assess complementarity structure compatibility before merging 
approaches.

%------------------------------------------------------------------------------
\subsubsection{Empirical Evidence on IT-Complementarity}

The IT productivity paradox---decades of IT investment with little measurable 
productivity gain---was resolved by complementarity theory.

\begin{theorem}[IT Complementarity Resolution]
\label{thm:it-paradox}
\textup{(\citet{brynjolfsson2003computing})}
\begin{enumerate}
    \item Short-run IT returns: Normal (comparable to other capital)
    \item Long-run IT returns (5-7 years): Significantly higher
    \item Explanation: Firms require years to implement complementary organizational 
    innovations (decentralization, process redesign, skill upgrading)
\end{enumerate}
\end{theorem}

\begin{table}[h]
\centering
\small
\begin{tabular}{p{4cm}p{4cm}p{4cm}}
\toprule
\textbf{Study} & \textbf{Finding} & \textbf{Complementarity Evidence} \\
\midrule
Brynjolfsson \& Hitt (2003) & Long-run IT returns 5$\times$ short-run & Organizational co-investment \\
Bloom et al. (2007) & US firms in UK more productive & Imported complementary practices \\
Ichniowski et al. (1997) & Bundled HR $>$ isolated HR & Direct $\gamma_{12} > 0$ \\
Bartel et al. (2007) & IT + skills + incentives & Three-way complementarity \\
\bottomrule
\end{tabular}
\caption{Empirical evidence for organizational complementarity}
\label{tab:empirical-evidence}
\end{table}

%------------------------------------------------------------------------------
\subsubsection{Literature Integration}

\begin{table}[h]
\centering
\small
\begin{tabular}{p{3.5cm}p{4cm}p{5cm}}
\toprule
\textbf{Source} & \textbf{Contribution} & \textbf{EBF Application} \\
\midrule
Brynjolfsson \& Milgrom (2013) & Ten theorems, measurement & Empirical methodology \\
Ichniowski et al. (1997) & Steel finishing lines & Homogeneous population test \\
Brynjolfsson \& Hitt (2003) & IT productivity paradox & Long-run complementarity \\
Bloom et al. (2007) & Cross-country evidence & Cultural transmission \\
Milgrom \& Roberts (1990) & Modern manufacturing & Original theory \\
\bottomrule
\end{tabular}
\caption{Empirical literature on organizational complementarity}
\label{tab:empirical-lit}
\end{table}




%==============================================================================
% APPENDIX B - ADDENDA: NOTATION, GLOSSARY, AXIOM INDEX
%==============================================================================

%------------------------------------------------------------------------------

%==============================================================================
% TEIL IX: PARAMETER ESTIMATION METHODOLOGY
%==============================================================================

\subsection{Parameter Estimation}
\label{sec:B-estimation-ref}

The estimation of complementarity parameters ($\Gamma$, $\Lambda$, $\delta$) is 
treated comprehensively in \textbf{Appendix EST: The Epistemics of Parameter 
Estimation}. That appendix provides:

\begin{itemize}
    \item The three-tier hybrid estimation strategy (Tier 1: Empirical, Tier 2: 
    Meta-analytical, Tier 3: Synthetic priors)
    \item Five complementary estimation methods (Econometric, LLM Monte Carlo, 
    Experimental, Meta-analytic, Milgrom Diagnostic)
    \item The complete Parameter Dictionary with current best estimates
    \item Bayesian updating protocols for continuous calibration
    \item Uncertainty quantification and context heterogeneity
\end{itemize}

\textbf{Key parameters estimated in Appendix EST:}
\begin{center}
\begin{tabular}{lcl}
\toprule
\textbf{Parameter Class} & \textbf{Count} & \textbf{Primary Source} \\
\midrule
$\Gamma = [\gamma_{dd'}]$ & 15 & Tier 1/2 (PSID, ESS) \\
$\Lambda = [\lambda_{kl}]$ & 28 & Tier 2 (Cross-country) \\
$\delta_{dd'k}$ & 120 & Tier 2/3 (Meta + LLM-MC) \\
$\Psi^*_k$ (thresholds) & 9 & Tier 2 (Literature) \\
\bottomrule
\end{tabular}
\end{center}

See Appendix EST for the complete methodological framework and current parameter 
values.


%==============================================================================
% SUMMARY BOX (Required)
%==============================================================================


%==============================================================================
% KEY RESULTS (Required for Core Content compliance)
%==============================================================================

% -----------------------------------------------------------------------------
% APPENDIX SCOPE BOX
% -----------------------------------------------------------------------------

\begin{tcolorbox}[
    colback=orange!5!white,
    colframe=orange!75!black,
    title={\textbf{Appendix Scope: Ziel / In-Scope / Out-of-Scope / Constraints}},
    fonttitle=\bfseries
]

\textbf{Ziel (Objective):}
\begin{itemize}[nosep]
    \item [TODO: Define objective for Unknown Title]
\end{itemize}

\textbf{In-Scope:}
\begin{itemize}[nosep]
    \item Key Results
\end{itemize}

\textbf{Out-of-Scope (delegiert an andere Appendices):}
\begin{itemize}[nosep]
    \item Aggregation Levels $\rightarrow$ Appendix WHO
    \item FEPSDE Utility Structure $\rightarrow$ Appendix WAT
    \item 8Ψ Context Dimensions $\rightarrow$ Appendix CTW
    \item Parameter Calibration $\rightarrow$ Appendix EST
    \item CORE-HIERARCHY $\rightarrow$ Appendix HIE
\end{itemize}

\textbf{Constraints (Anwendungsgrenzen):}
\begin{itemize}[nosep]
    \item [TODO: Define when this appendix does NOT apply]
    \item [TODO: Define required prerequisites]
    \item [TODO: Define known limitations]
\end{itemize}

\textbf{Lieferobjekte (Deliverables):}
\begin{itemize}[nosep]
    \item 18 Table(s)
    \item 146 Equation(s)
    \item Worked Example(s)
\end{itemize}

\end{tcolorbox}


\section{Key Results}
\label{sec:B-key-results}

This section summarizes the main findings from the complementarity analysis.

\subsection{Result 1: The $\Gamma$-Matrix Structure}

\begin{tcolorbox}[colback=green!5!white, colframe=green!75!black]
\textbf{Finding:} The majority of FEPSDE dimension pairs exhibit positive 
complementarity ($\gamma_{dd'} > 0$), indicating that improvements in one 
dimension enhance the marginal value of improvements in others.

\textbf{Evidence:} 
\begin{itemize}[nosep]
    \item 12 of 15 $\gamma$-coefficients are positive
    \item Strongest: Financial-Physical ($\gamma_{FP} \approx 0.09$)
    \item Weakest/Negative: Financial-Ecological under scarcity
\end{itemize}
\end{tcolorbox}

\subsection{Result 2: Context Modulation is Substantial}

\begin{tcolorbox}[colback=green!5!white, colframe=green!75!black]
\textbf{Finding:} Context ($\boldsymbol{\Psi}$) modulates complementarity by 
up to $\pm 50\%$ of baseline values.

\textbf{Evidence:}
\begin{itemize}[nosep]
    \item Average $|\delta|$ effect: 0.03--0.05 per $\Psi$-dimension
    \item Cumulative effect across 8 dimensions: substantial
    \item Dominant modulators vary by $\gamma$-pair (see Part III)
\end{itemize}
\end{tcolorbox}

\subsection{Result 3: Supermodularity Enables Tractability}

\begin{tcolorbox}[colback=green!5!white, colframe=green!75!black]
\textbf{Finding:} When the utility function is supermodular, optimization is 
tractable despite high dimensionality.

\textbf{Implications:}
\begin{itemize}[nosep]
    \item Monotone comparative statics hold (Theorem B.1)
    \item Equilibria can be computed by lattice iteration
    \item Bundled interventions are generically optimal
\end{itemize}
\end{tcolorbox}

\subsection{Result 4: Governance Should Match Complementarity}

\begin{tcolorbox}[colback=green!5!white, colframe=green!75!black]
\textbf{Finding:} Organizational structure should reflect complementarity type:
\begin{itemize}[nosep]
    \item \textbf{Strong complementarity} $\Rightarrow$ Unified governance
    \item \textbf{Supermodular complementarity} $\Rightarrow$ Distributed ecosystems
    \item \textbf{Thin crossing points} $\Rightarrow$ Contractual solutions
\end{itemize}

\textbf{Evidence:} Baldwin (2024) DSMC framework; empirical cases in Part VIII.
\end{tcolorbox}



\section{Summary}
\label{sec:B-summary}

\begin{tcolorbox}[colback=green!10!white, colframe=green!75!black, 
    title=\textbf{Appendix HOW: Complete Summary}]

\textbf{Core Question:} HOW do utility dimensions interact?

\textbf{Answer:} Through \textbf{complementarity}---the systematic pattern where 
changing one dimension affects the marginal value of changing another.

\textbf{Main Contributions:}
\begin{enumerate}[nosep]
    \item \textbf{$\Gamma$-Matrix:} Complete specification of utility$\times$utility 
    complementarity (15 parameters for FEPSDE dimensions)
    \item \textbf{$\Lambda$-Matrix:} Context$\times$context interactions 
    (28 parameters for 8Ψ dimensions)
    \item \textbf{$\delta$-Tensor:} How context modulates utility complementarity 
    (120 parameters)
    \item \textbf{Applications:} Mechanism design, matching, large games, organizations
\end{enumerate}

\textbf{Key Outputs:}
\begin{itemize}[nosep]
    \item $\Gamma = [\gamma_{dd'}]$: Symmetric $6 \times 6$ matrix
    \item $\Lambda = [\lambda_{kl}]$: Symmetric $8 \times 8$ matrix
    \item $\delta = [\delta_{dd',k}]$: $15 \times 8$ modulation tensor
    \item Total: 163 interaction parameters
\end{itemize}

\textbf{Key Results:}
\begin{itemize}[nosep]
    \item Most FEPSDE pairs are complements ($\gamma > 0$)
    \item Context modulates complementarity through $\delta$-coefficients
    \item Supermodularity enables tractable optimization (Theorem B.1)
    \item Organizations should match governance to complementarity type
\end{itemize}

\textbf{Integration:} Connects to AAA (levels), C (dimensions), V (context), 
BBB (calibration), forming the complete EBF interaction specification.

\end{tcolorbox}

%==============================================================================
% CRITICAL FOUNDATIONS / OBJECTIONS (Required for CORE)
%==============================================================================

\section{Critical Foundations and Objections}
\label{sec:B-foundations-objections}

\subsection{Objection 1: Supermodularity is Too Restrictive}

\textbf{Objection:} Real utility functions may not satisfy supermodularity. 
Requiring $\partial^2 U / \partial x_i \partial x_j \geq 0$ for all pairs 
is empirically implausible.

\textbf{Response:} 
\begin{enumerate}[nosep]
    \item We do \textit{not} require global supermodularity---only local 
    supermodularity in the relevant decision region.
    \item The ordinal equivalence theorem (Chambers \& Echenique 2009) shows 
    that under monotonicity, supermodularity is equivalent to the weaker 
    quasisupermodularity condition.
    \item Our framework explicitly allows substitutes ($\gamma < 0$) where 
    empirically justified.
    \item The $\Gamma$-matrix is calibrated from data, not assumed 
    \textit{a priori}.
\end{enumerate}

\subsection{Objection 2: Too Many Parameters}

\textbf{Objection:} 163 parameters is too many to estimate reliably. The model 
is over-parameterized and risks overfitting.

\textbf{Response:}
\begin{enumerate}[nosep]
    \item Most parameters have strong theoretical priors (sign, magnitude range).
    \item The three-tier epistemic system (Appendix EST) explicitly distinguishes 
    empirical estimates from theoretical placeholders.
    \item Symmetry and sparsity reduce effective dimensionality.
    \item LLM-MC methodology (Appendix LLM1) provides synthetic priors that 
    regularize estimation.
    \item Sensitivity analysis identifies which parameters matter for 
    specific applications.
\end{enumerate}

\subsection{Objection 3: Context Modulation is Ad Hoc}

\textbf{Objection:} The linear modulation formula 
$\gamma(\Psi) = \bar{\gamma} + \sum_k \delta_k (\Psi_k - 0.5)$ is arbitrary. 
Why linear? Why centered at 0.5?

\textbf{Response:}
\begin{enumerate}[nosep]
    \item Linear modulation is a first-order Taylor approximation around 
    the global mean---standard in economics.
    \item Centering at 0.5 means $\bar{\gamma}$ represents complementarity 
    at ``average'' context.
    \item Non-linear terms ($\Psi^2$) can be added when empirically warranted.
    \item The framework is modular: improved modulation functions can be 
    substituted without changing the overall architecture.
\end{enumerate}

\subsection{Objection 4: Ignores Dynamics}

\textbf{Objection:} Complementarity may change over time as individuals learn, 
habits form, and contexts evolve. A static $\Gamma$-matrix misses this.

\textbf{Response:}
\begin{enumerate}[nosep]
    \item Part IV addresses dynamic context$\times$context complementarity.
    \item The $\Lambda$-matrix captures how context dimensions evolve together.
    \item Adaptive mechanism design (Part V) explicitly models learning.
    \item For most policy applications, the relevant time horizon is short 
    enough that $\Gamma$ can be treated as approximately stable.
\end{enumerate}

\subsection{Behavioral Game Theory Foundation: The Beauty Contest}
\label{sec:B-beauty-contest}

The mathematical theory of complementarity (Topkis, Milgrom-Roberts) has a direct
behavioral counterpart in experimental economics. Mauersberger \& Nagel (2018)
establish that the \textbf{Keynesian Beauty Contest} is the canonical form underlying
strategic complementarity and substitutability.

\begin{tcolorbox}[colback=yellow!5!white, colframe=yellow!75!black,
    title=\textbf{Mauersberger \& Nagel (2018): Strategic Complementarity in Action}]

\textbf{Core Finding:}
Many seemingly different games share a common \textit{best-response structure}:
\[
x^* = f\bigl(\mathbb{E}[\bar{x}_{-i}]\bigr)
\]
``My optimal decision is a function of my expectation about others' behavior.''

\textbf{Connection to $\gamma$:}
\begin{itemize}[nosep]
    \item $\gamma > 0$ (complements): $f'(\cdot) > 0$ --- I do more when others do more
    \item $\gamma < 0$ (substitutes): $f'(\cdot) < 0$ --- I do less when others do more
\end{itemize}

\textbf{Key Insight:} The sign and magnitude of $\gamma$ determines whether
behavioral deviations \textbf{amplify} (complements) or \textbf{dampen} (substitutes).

\end{tcolorbox}

\paragraph{Implications for EBF Complementarity.}
The Mauersberger/Nagel framework provides three key insights for the $\Gamma$-matrix:

\begin{enumerate}[nosep]
    \item \textbf{Amplification vs. Dampening:} When $\gamma_{ij} > 0$, small deviations
    from equilibrium grow; when $\gamma_{ij} < 0$, they shrink. This explains why
    some interventions ``cascade'' while others ``fizzle out.''

    \item \textbf{Level-$k$ Heterogeneity:} Agents differ in how many steps of
    strategic reasoning they perform (typically Levels 0--3). This creates
    \textit{structured heterogeneity}---not noise, but systematic variation
    related to cognitive bounds. See Appendix AWA (Awareness).

    \item \textbf{Design $>$ Equilibrium:} Anchors, signals, and framing shift
    behavior massively---often without changing the equilibrium. This validates
    EBF's emphasis on context $\Psi$ over pure payoff structure.
\end{enumerate}

\paragraph{Games Covered by the Beauty Contest Structure.}
The following games exhibit the canonical best-response structure and are thus
directly interpretable through the $\gamma$-lens:

\begin{center}
\renewcommand{\arraystretch}{1.2}
\begin{tabular}{@{}p{4.5cm}p{3cm}p{5cm}@{}}
\toprule
\textbf{Game/Market} & \textbf{$\gamma$ Sign} & \textbf{Behavioral Implication} \\
\midrule
Coordination games & $\gamma > 0$ & Deviations amplify; multiple equilibria \\
Public Goods games & $\gamma > 0$ & Free-riding cascades \\
Asset markets (bubbles) & $\gamma > 0$ & Speculative spirals \\
Cournot competition & $\gamma < 0$ & Deviations dampen; unique equilibrium \\
Auctions (first-price) & $\gamma < 0$ & Strategic underbidding \\
\bottomrule
\end{tabular}
\end{center}

\textbf{Cross-Reference:} For the complete treatment of Mauersberger \& Nagel (2018)
including the ``generative grammar'' interpretation and worked examples, see
Appendix MSC Section 11.

%==============================================================================
% OPEN ISSUES / RESEARCH AGENDA (Required)
%==============================================================================

\section{Open Issues and Research Agenda}
\label{sec:B-open-issues}

\subsection{Methodological Priorities}

\begin{enumerate}
    \item \textbf{Higher-Order Complementarity:} Current framework focuses on 
    pairwise interactions ($\gamma_{dd'}$). Three-way and higher-order 
    complementarities ($\gamma_{dd'd''}$) may matter in some applications.
    
    \item \textbf{Non-Linear Modulation:} When does the linear $\delta$-modulation 
    break down? Identifying threshold effects and regime changes.
    
    \item \textbf{Aggregation:} How does individual-level $\Gamma$ aggregate to 
    group-level? Heterogeneity in complementarity across agents.
    
    \item \textbf{Identification:} Econometric strategies for identifying 
    $\gamma$ from observational data without experiments.
\end{enumerate}

\subsection{Empirical Priorities}

\begin{enumerate}
    \item \textbf{Direct $\Gamma$ Estimation:} More RCT-based estimates of 
    specific $\gamma_{dd'}$ coefficients, especially for understudied pairs.
    
    \item \textbf{Cross-Context Validation:} Testing whether $\delta$-coefficients 
    estimated in one country predict effects in another.
    
    \item \textbf{Temporal Stability:} How stable is $\Gamma$ over time? 
    Decade-scale panel studies needed.
    
    \item \textbf{Milgrom Tests at Scale:} Applying pairwise supermodularity 
    diagnostics to large datasets.
\end{enumerate}

\subsection{Theoretical Extensions}

\begin{enumerate}
    \item \textbf{Stochastic Complementarity:} How does uncertainty affect 
    optimal bundling of interventions?
    
    \item \textbf{Evolutionary Foundations:} Why do observed complementarities 
    have the sign structure they do? Adaptive origins.
    
    \item \textbf{Network Complementarity:} Richer models of how complementarity 
    propagates through social networks.
    
    \item \textbf{Mechanism Design with Unknown $\Gamma$:} Robust mechanisms 
    that perform well across a range of complementarity structures.
\end{enumerate}



\subsection{Notation and Symbols}
\label{sec:B-notation}

\begin{table}[h]
\centering
\small
\begin{tabular}{p{2.5cm}p{10cm}}
\toprule
\textbf{Symbol} & \textbf{Definition} \\
\midrule
\multicolumn{2}{l}{\textit{Core Structures}} \\
$\Gamma = [\gamma_{dd'}]$ & Utility-Complementarity Matrix: captures pairwise interactions between utility dimensions \\
$\gamma_{dd'}$ & Complementarity coefficient between dimensions $d$ and $d'$; positive = complements, negative = substitutes \\
$\Lambda = [\lambda_{kl}]$ & Context-Complementarity Matrix: captures interactions between context dimensions \\
$\lambda_{kl}$ & Interaction coefficient between context dimensions $k$ and $l$ \\
$\boldsymbol{\psi} = (\psi_1, \ldots, \psi_8)$ & Context vector in 8Ψ-Dimensions framework \\
\midrule
\multicolumn{2}{l}{\textit{Lattice Operations}} \\
$x \vee y$ & Join operator: component-wise maximum of $x$ and $y$ \\
$x \wedge y$ & Meet operator: component-wise minimum of $x$ and $y$ \\
$\mathcal{L}$ & Lattice: partially ordered set with join and meet \\
$\mathcal{S} \subseteq \mathcal{L}$ & Sublattice: subset closed under $\vee$ and $\wedge$ \\
\midrule
\multicolumn{2}{l}{\textit{Utility and Value}} \\
$U(x)$, $V(x)$ & Utility/Value function \\
$U_d$ & Utility in dimension $d \in \{F, E, P, S, D, E\}$ (FEPSDE) \\
$\frac{\partial^2 U}{\partial x_i \partial x_j}$ & Cross-partial derivative (complementarity measure) \\
\midrule
\multicolumn{2}{l}{\textit{Supermodularity}} \\
$f: \mathcal{L} \to \mathbb{R}$ supermodular & $f(x \vee y) + f(x \wedge y) \geq f(x) + f(y)$ for all $x, y$ \\
Increasing differences & $f(x', y') - f(x', y) \geq f(x, y') - f(x, y)$ for $x' > x$, $y' > y$ \\
$\infty$-supermodular & Supermodular under all monotone transformations \\
$r$-decomposable & Representable as maximum of $r$ modular functions \\
\midrule
\multicolumn{2}{l}{\textit{Game Theory}} \\
$\bar{a}^*, \underline{a}^*$ & Extremal Nash equilibria (highest and lowest) \\
$\mu^*$ & Equilibrium distribution in mean-field games \\
$\gamma^{\text{pop}}$ & Population-level complementarity parameter \\
$1/(1 - \gamma^{\text{pop}})$ & Social multiplier under complementarity \\
\midrule
\multicolumn{2}{l}{\textit{Organizational}} \\
DSMC & Distributed Supermodular Complementarity \\
$V_X(x; y)$ & Value for actor $X$, treating $y$ as parameter \\
$C_{\text{int}}$ & Integration costs \\
$V^{FB}$ & First-best value under central planning \\
\midrule
\multicolumn{2}{l}{\textit{Indices and Sets}} \\
$d, d' \in \mathcal{D}$ & Utility dimensions (typically FEPSDE: $|\mathcal{D}| = 6$) \\
$k, l \in \mathcal{K}$ & Context dimensions (8Ψ: $|\mathcal{K}| = 8$) \\
$i, j \in \mathcal{N}$ & Agents/players \\
$t \in \mathcal{T}$ & Time periods \\
\bottomrule
\end{tabular}
\caption{Notation guide for Appendix HOW}
\label{tab:notation}
\end{table}

%------------------------------------------------------------------------------

%==============================================================================
% GLOSSARY OF SYMBOLS (Required - links to Appendix GLS)
%==============================================================================

\section{Glossary of Symbols}
\label{sec:B-glossary}

\begin{tcolorbox}[colback=blue!5!white, colframe=blue!75!black]
\textbf{Central Reference:} For complete symbol definitions across the EBF 
framework, see \textbf{Appendix GLS} (\textit{Glossary and Notation}).

This section lists symbols introduced or specialized in this appendix.
\end{tcolorbox}

\vspace{0.5em}

\begin{longtable}{@{}clp{6cm}c@{}}
\toprule
\textbf{Symbol} & \textbf{Name} & \textbf{Definition} & \textbf{In G?} \\
\midrule
\endfirsthead
\midrule
\endhead
\midrule
\endfoot
\bottomrule
\endlastfoot

\multicolumn{4}{l}{\textbf{Core Complementarity Symbols}} \\
\midrule
$\gamma_{dd'}$ & Complementarity coefficient & $\partial^2 U / \partial x_d \partial x_{d'}$ & \checkmark \\
$\Gamma$ & Gamma matrix & $[\gamma_{dd'}]_{d,d' \in \text{FEPSDE}}$, $6 \times 6$ symmetric & \checkmark \\
$\bar{\gamma}_{dd'}$ & Baseline complementarity & Global average at $\Psi_k = 0.5$ & \checkmark \\
$\gamma_{dd'}(\Psi)$ & Context-adjusted & $\bar{\gamma} + \sum_k \delta_k (\Psi_k - 0.5)$ & \checkmark \\
\midrule

\multicolumn{4}{l}{\textbf{Context Interaction Symbols}} \\
\midrule
$\lambda_{kl}$ & Context interaction & How context dimensions $k, l$ interact & \checkmark \\
$\Lambda$ & Lambda matrix & $[\lambda_{kl}]_{k,l \in \{1..8\}}$, $8 \times 8$ symmetric & \checkmark \\
$\delta_{dd',k}$ & Modulation coefficient & Effect of $\Psi_k$ on $\gamma_{dd'}$ & \checkmark \\
$\delta$ & Delta tensor & $[\delta_{dd',k}]$, $15 \times 8$ tensor & \checkmark \\
\midrule

\multicolumn{4}{l}{\textbf{Supermodularity Symbols}} \\
\midrule
$\vee$ & Join (lattice) & $x \vee y = \max(x, y)$ componentwise & NEW \\
$\wedge$ & Meet (lattice) & $x \wedge y = \min(x, y)$ componentwise & NEW \\
$f^{SM}$ & Supermodular & $f(x \vee y) + f(x \wedge y) \geq f(x) + f(y)$ & NEW \\
$f^{QSM}$ & Quasisupermodular & Ordinal version of supermodularity & NEW \\
$r$ & Decomposition rank & For r-decomposable optimization & NEW \\
\midrule

\multicolumn{4}{l}{\textbf{Game-Theoretic Symbols}} \\
\midrule
$\mathcal{G}$ & Supermodular game & Game with strategic complementarity & NEW \\
$x^*$ & Nash equilibrium & Fixed point of best-response & \checkmark \\
$\underline{x}, \bar{x}$ & Extreme equilibria & Smallest and largest equilibria & NEW \\
$\kappa$ & Tipping point & Critical $\Psi$ for regime change & NEW \\
\midrule

\multicolumn{4}{l}{\textbf{Organizational Symbols}} \\
\midrule
DSMC & Distributed Supermodular & Baldwin's ecosystem condition & NEW \\
TCP & Thin Crossing Point & Low-cost interface between modules & NEW \\

\end{longtable}

\textbf{Note:} \checkmark = Defined in Appendix GLS; NEW = Introduced in this appendix 
(will be added to G in next update).



\subsection{Complete Axiom Index}
\label{sec:B-axiom-index}

The following table lists all axioms defined in this appendix:

\begin{longtable}{p{2cm}p{4cm}p{6.5cm}p{2cm}}
\toprule
\textbf{Axiom} & \textbf{Name} & \textbf{Description} & \textbf{Section} \\
\midrule
\endfirsthead
\toprule
\textbf{Axiom} & \textbf{Name} & \textbf{Description} & \textbf{Section} \\
\midrule
\endhead
\midrule
\multicolumn{4}{r}{\textit{Continued on next page}} \\
\endfoot
\bottomrule
\endlastfoot

\multicolumn{4}{l}{\textbf{Teil I: Foundations}} \\
B-MONO & Monotonicity & Utility increasing in each dimension & \ref{sec:B-foundations} \\
B-SUPER & Supermodularity & Cross-partials non-negative & \ref{sec:B-foundations} \\
B-ORD & Ordinal Equivalence & Super $\Leftrightarrow$ Quasisuper under monotonicity & \ref{sec:B-foundations} \\
B-$\infty$SUP & Infinite Supermodularity & Preserved under monotone transforms & \ref{sec:B-foundations} \\
B-COMP & Computational Tractability & r-decomposability for optimization & \ref{sec:B-foundations} \\
\midrule

\multicolumn{4}{l}{\textbf{Teil II: Utility × Utility}} \\
B-SYM & Symmetry & $\gamma_{dd'} = \gamma_{d'd}$ by Young's theorem & Teil II \\
B-SGN & Sign Structure & Expected signs of $\Gamma$-matrix entries & Teil II \\
B-FP & Financial-Physical & $\gamma_{FP} > 0$ (health-wealth complementarity) & Teil II \\
B-FE & Financial-Emotional & Sign depends on income level & Teil II \\
B-EP & Emotional-Physical & Bidirectional reinforcement & Teil II \\
\midrule

\multicolumn{4}{l}{\textbf{Teil III: Utility × Context}} \\
B-AWX & Awareness-Complementarity & $\gamma_{dd'}(\boldsymbol{\psi})$ context-dependent & Teil III \\
B-VIS & Visceral State-Dependence & Hot states modulate $\gamma$ & Teil III \\
B-CTX & Context Modulation & Two channels: direct and indirect & Teil III \\
B-BELIEF & Motivated Beliefs & Belief filters affect $\gamma$ perception & Teil III \\
B-REF & Reference Dependence & Complementarity relative to reference point & Teil III \\
\midrule

\multicolumn{4}{l}{\textbf{Teil IV: Context × Context}} \\
B-DYN & Dynamic Context Complementarity & $\lambda_{kl}$ over time & Teil IV \\
B-SEQ & Intervention Sequencing & Optimal ordering under $\Lambda$ & Teil IV \\
B-INF & Information Complementarity & Knowledge × Trust interaction & Teil IV \\
B-LVL & Level Complementarity & Context dimensions reinforce across levels & Teil IV \\
\midrule

\multicolumn{4}{l}{\textbf{Teil V: Mechanism Design}} \\
B-ATT & Attention-Compatible Mechanisms & Simplicity, salience, chunking & Teil V \\
B-ADP & Adaptive Mechanism Design & Learn $\Gamma$ over time & Teil V \\
B-BND & Bundled Nudges & Superadditive returns from joint nudges & Teil V \\
\midrule

\multicolumn{4}{l}{\textbf{Teil VI: Large Games}} \\
B-NET & Network Complementarity & Local + global structure & Teil VI \\
B-PROP & Dynamic Propagation & Monotone spread through network & Teil VI \\
B-MIX & Mixed Equilibria & Coexistence of multiple equilibria & Teil VI \\
\midrule

\multicolumn{4}{l}{\textbf{Teil VII: Matching}} \\
B-MATCH & Intervention-Context Matching & Optimize $\boldsymbol{\psi} \cdot \Gamma \cdot \mathbf{d}$ & Teil VII \\
B-REMATCH & Dynamic Re-Matching & Adjust to context evolution & Teil VII \\
B-ASSIGN & Assignment Efficiency & Stable matching under complementarity & Teil VII \\
\midrule

\multicolumn{4}{l}{\textbf{Teil VIII: Organizations}} \\
B-DSMC & Distributed Supermodular Complementarity & Three conditions for ecosystems & Teil VIII \\
B-GOV & Governance Mode Selection & Strong $\Rightarrow$ Unified, Super $\Rightarrow$ Distributed & Teil VIII \\
B-CHERRY & Anti-Cherry-Picking & Don't adopt practices in isolation & Teil VIII \\
B-COH & Coherent Search & Move all variables in same direction & Teil VIII \\
B-TEST & Pairwise Testing & If all pairs pass, system is supermodular & Teil VIII \\
B-SEG & Segmentation & Match governance to complementarity type & Teil VIII \\


\end{longtable}

%------------------------------------------------------------------------------
\subsection{Cross-References to Other Appendices}
\label{sec:B-crossref}

This appendix connects to the broader EBF framework through the following 
relationships:

\begin{table}[h]
\centering
\small
\begin{tabular}{p{2.5cm}p{3cm}p{7cm}}
\toprule
\textbf{Appendix} & \textbf{CORE Question} & \textbf{Connection to Appendix HOW} \\
\midrule
AAA & WHO has utility? & Aggregation levels $L$ determine which $\Gamma$-matrices apply \\
C & WHAT is utility? & FEPSDE dimensions $d$ define rows/columns of $\Gamma$ \\
V & WHEN does context matter? & 8Ψ-Dimensions define context vector $\boldsymbol{\psi}$ \\
AU & Axiom Formalization & B-axioms integrate with master axiom system \\
AN & LLM Monte Carlo & Methodology for estimating $\gamma$ coefficients \\
\textbf{NC} & \textbf{Non-Concavity \& Fit} & \textbf{$\gamma > 0 \Rightarrow$ non-concave landscape; $K$ vs $Q$ distinction (Ch.7)} \\
\bottomrule
\end{tabular}
\caption{Cross-references to other EBF appendices}
\label{tab:crossref}
\end{table}




%------------------------------------------------------------------------------
% CENTRAL REFERENCES (Required per Template)
%------------------------------------------------------------------------------

\begin{tcolorbox}[colback=blue!5!white, colframe=blue!75!black, 
    title=\textbf{Central References}]

\textbf{Glossary:} For complete symbol definitions across the EBF framework, 
see \textbf{Appendix GLS} (\textit{Glossary and Notation}).

\textbf{Bibliography:} For full reference details, see the 
\textbf{Master Bibliography} (\texttt{bcm2\_master\_references.bib}) 
in the repository.

All symbols used in this appendix are consistent with Appendix GLS definitions.

\end{tcolorbox}

\section{References}
\label{sec:B-references}

% Note: This appendix references the following key works. Full bibliographic
% entries are provided in appendix_B_references.bib

\subsubsection{Foundational Works on Supermodularity}

The mathematical foundations of complementarity and supermodularity were 
established by:

\begin{itemize}
    \item \textbf{Topkis (1979, 1998)}: Foundational theory of supermodular 
    functions and games, lattice-theoretic approach
    \item \textbf{Milgrom \& Shannon (1994)}: Monotone comparative statics, 
    single-crossing property, ordinal approach
    \item \textbf{Milgrom \& Roberts (1990)}: Strategic complementarity in games, 
    learning and rationalizability
    \item \textbf{Vives (1990, 1999)}: Nash equilibrium with strategic 
    complementarities, oligopoly applications
    \item \textbf{Tarski (1955)}: Fixed-point theorem for lattices
\end{itemize}

\subsubsection{Ordinal Foundations and Extensions}

\begin{itemize}
    \item \textbf{Chambers \& Echenique (2009)}: Ordinal equivalence between 
    supermodularity and quasisupermodularity under monotonicity
    \item \textbf{Chateauneuf et al. (2017)}: Infinite supermodularity and 
    complete characterization of complementarity orders
    \item \textbf{Chen et al. (2025)}: r-Decomposability and computational 
    tractability of supermodular optimization
    \item \textbf{Amir (2019)}: Survey of recent advances in supermodularity
\end{itemize}

\subsubsection{Behavioral Economics and Context}

\begin{itemize}
    \item \textbf{Sims (2003)}: Rational inattention and information processing 
    costs
    \item \textbf{Loewenstein (1996, 2005)}: Visceral factors, hot-cold empathy 
    gap, affect and deliberation
    \item \textbf{Bénabou \& Tirole (2016)}: Motivated beliefs, self-signaling, 
    identity economics
    \item \textbf{Bordalo et al. (2012)}: Salience theory, context-dependent 
    attention
    \item \textbf{Kőszegi \& Rabin (2006)}: Reference-dependent preferences
    \item \textbf{Mullainathan \& Shafir (2013)}: Scarcity and bandwidth tax
    \item \textbf{Becker \& Murphy (1988)}: Rational addiction and habit formation
\end{itemize}

\subsubsection{Game Theory and Strategic Interactions}

\begin{itemize}
    \item \textbf{Balbus et al. (2019)}: Supermodular games with continuum of 
    players
    \item \textbf{Jackson (2008, 2015)}: Social and economic networks, games on 
    networks
    \item \textbf{Morris \& Shin (2002)}: Global games, coordination, public 
    information
    \item \textbf{Acemoglu \& Ozdaglar (2011)}: Opinion dynamics in networks
    \item \textbf{Barthel \& Hoffmann (2019)}: Mixed-monotonic games
\end{itemize}

\subsubsection{Mechanism Design and Information}

\begin{itemize}
    \item \textbf{Kushnir \& Liu (2019)}: Equivalence of BIC and DSIC under 
    complementarity
    \item \textbf{Kamenica \& Gentzkow (2011)}: Bayesian persuasion
    \item \textbf{Bergemann \& Morris (2019)}: Information design unified 
    perspective
    \item \textbf{Pavan et al. (2014)}: Dynamic mechanism design
    \item \textbf{Carroll (2015)}: Robustness in contract theory
\end{itemize}

\subsubsection{Matching Theory}

\begin{itemize}
    \item \textbf{Becker (1973)}: Assortative matching, theory of marriage
    \item \textbf{Gale \& Shapley (1962)}: Stable matching algorithm
    \item \textbf{Kelso \& Crawford (1982)}: Gross substitutes condition
    \item \textbf{Johnson (2019)}: Matching with incomplete information
    \item \textbf{Villani (2009)}: Optimal transport theory
    \item \textbf{Chiappori et al. (2017)}: Multi-dimensional matching
\end{itemize}


\subsubsection{Organizational Economics}

\begin{itemize}
    \item \textbf{Baldwin (2024)}: Design Rules Vol. 2, Distributed Supermodular 
    Complementarity, thin crossing points, ecosystem architecture
    \item \textbf{Baldwin \& Clark (2000)}: Design Rules Vol. 1, modularity as 
    source of optionality
    \item \textbf{Brynjolfsson \& Milgrom (2013)}: Ten theorems, empirical methods, 
    Matrix of Change, organizational dynamics
    \item \textbf{Williamson (1985)}: Transaction cost economics, asset specificity, 
    governance modes
    \item \textbf{Hart \& Moore (1990)}: Property rights theory, residual control, 
    holdup problems
    \item \textbf{Milgrom \& Roberts (1990)}: Complementarities in modern manufacturing, 
    original supermodularity applications to organizations
\end{itemize}

\subsubsection{Platform Economics and Networks}

\begin{itemize}
    \item \textbf{Rochet \& Tirole (2003)}: Two-sided markets, platform competition, 
    indirect network effects
    \item \textbf{Coase (1960)}: Social cost theorem, side payments, externality 
    internalization
    \item \textbf{Athey \& Ellison (2014)}: Platform dynamics, search design
\end{itemize}

\subsubsection{Empirical Studies}

\begin{itemize}
    \item \textbf{Ichniowski et al. (1997)}: Steel finishing lines, bundled HR 
    practices, direct complementarity evidence
    \item \textbf{Brynjolfsson \& Hitt (2003)}: IT productivity paradox, long-run 
    vs. short-run returns
    \item \textbf{Bloom \& Van Reenen (2007)}: Cross-country management practices, 
    US firms in UK
    \item \textbf{Bartel et al. (2007)}: IT, skills, and incentives as three-way 
    complements
\end{itemize}



% For LaTeX compilation with BibTeX:
% \bibliographystyle{aer}
% \bibliography{appendix_B_references}

\end{document}
